% =====================================================================
%  Anatomía de un LLM — Fase VI — Módulo 26
%  Control óptimo, gemelos digitales y aprendizaje por refuerzo para procesos
%  Compilar: latexmk -r ../latexmkrc M26-control-rl.tex
% =====================================================================
\documentclass[11pt,a4paper,twoside,openany]{book}
\usepackage{../estilo-apuntes}

\newcommand{\ApuntesCurso}{Anatomía de un LLM}
\newcommand{\ApuntesModulo}{Módulo 26}

% Notación propia del módulo
\newcommand{\Sset}{\mathcal{S}}
\newcommand{\Aset}{\mathcal{A}}
\newcommand{\Xset}{\mathbb{X}}
\newcommand{\Uset}{\mathbb{U}}
\newcommand{\Dt}{\Delta t}
\newcommand{\Text}{T_{\mathrm{ext}}}
\newcommand{\COP}{\mathrm{COP}}
\newcommand{\KL}{D_{\mathrm{KL}}}
\newcommand{\Var}{\operatorname{Var}}
\newcommand{\RMSE}{\operatorname{RMSE}}
\newcommand{\diag}{\operatorname{diag}}
\newcommand{\clip}{\operatorname{clip}}
\newcommand{\sat}{\operatorname{sat}}
\newcommand{\proy}{\operatorname{proy}}
\newcommand{\tr}{\operatorname{tr}}
\newcommand{\T}{^{\top}}
\makeatletter\g@addto@macro\UrlBreaks{\do\-\do\/}\makeatother

\hypersetup{pdftitle={Módulo 26. Control óptimo, gemelos digitales y aprendizaje por refuerzo para procesos},pdfauthor={Anatomía de un LLM}}

\begin{document}
\frontmatter
% ------------------------------------------------------------------ Portada
\begin{titlepage}
\centering
\vspace*{2.5cm}
{\large\scshape Anatomía de un LLM}\par
\vspace{0.4cm}
{\normalsize Fundamentos matemáticos de los modelos de lenguaje}\par
\vspace{2.2cm}
\rule{\textwidth}{0.8pt}\par\vspace{0.5cm}
{\small\scshape Fase VI \textperiodcentered\ Diseño y construcción de optimizadores con IA}\par\vspace{0.3cm}
{\Huge\bfseries Módulo 26\par}\vspace{0.3cm}
{\LARGE Control óptimo, gemelos digitales y aprendizaje por refuerzo para procesos\par}
\vspace{0.2cm}
{\large PID, identificación de sistemas, control predictivo, SAC y TD3, y puesta en producción en planta\par}
\vspace{0.5cm}\rule{\textwidth}{0.4pt}\par
\vspace{1.6cm}
{\large Apuntes del curso}\par
\vfill
{\small Edición de octubre de 2026}\par
\end{titlepage}

% ------------------------------------------------------- Presentación
\chapter*{Presentación del módulo}
\addcontentsline{toc}{chapter}{Presentación del módulo}
\markboth{Presentación del módulo}{}

Un proceso industrial es un sistema dinámico: la temperatura de un edificio, la de un reactor o la energía almacenada en una batería evolucionan en el tiempo según leyes físicas, responden con retraso a las acciones y están sujetas a perturbaciones y a límites de seguridad. Optimizar un proceso así no consiste en resolver un problema estático una vez, sino en decidir continuamente qué hacer a partir de lo que se mide. Este módulo construye, de abajo arriba, los tres grandes enfoques para hacerlo: el control realimentado clásico (PID), el control predictivo basado en modelo (MPC), que resuelve en cada instante un problema de optimización sobre un horizonte, y el aprendizaje por refuerzo (RL), que aprende una política por interacción con un simulador. Entre los tres se sitúa una pieza común: el modelo del proceso, identificado a partir de datos de planta con las redes neuronales de la Fase~II, que hace de \emph{gemelo digital}. El módulo termina con lo que separa un prototipo de un sistema en planta: la supervisión humana, las alarmas, la validación, la ciberseguridad de la tecnología de operación y las normas.

\section*{Objetivos}
Al finalizar el módulo, el estudiante debe ser capaz de:
\begin{enumerate}[label=\arabic*.]
  \item formular un proceso como sistema en espacio de estados, discretizarlo de forma exacta con retenedor de orden cero, analizar su estabilidad mediante autovalores y diseñar y ajustar un controlador PID discreto con antisaturación;
  \item identificar un modelo de un proceso a partir de datos (ARX, NARX, redes neuronales), validarlo en simulación libre y no solo a un paso, reconocer los problemas de identificación en lazo cerrado y construir con él un gemelo digital;
  \item formular un problema de control predictivo con restricciones como programa lineal o cuadrático, demostrar sus propiedades básicas (convexidad, penalización exacta, estabilidad con ingredientes terminales) y usar un modelo aprendido dentro del MPC;
  \item pasar del PPO del módulo~21 a los métodos fuera de política para acciones continuas (DDPG, TD3, SAC), diseñar recompensas sin alterar la política óptima, empaquetar un proceso como entorno Gymnasium y aplicar técnicas de RL seguro y aleatorización del dominio;
  \item comparar PID, MPC y RL con un protocolo justo en casos de climatización, reactores por lotes y baterías, y analizar con datos publicados el caso de la refrigeración de centros de datos de Google;
  \item planificar la puesta en producción en planta: arquitectura por capas, modos de operación con supervisión humana, alarmas, validación estadística, filtros de seguridad, ciberseguridad OT y marco normativo.
\end{enumerate}

\section*{Requisitos previos}
Del módulo~0.5 («Cálculo de una variable») se usa el polinomio de Taylor (tema~6) para linealizar; del módulo~0.6, las series geométricas (tema~5) y la integral para resolver ecuaciones lineales; del módulo~0.7, los multiplicadores de Lagrange (§5.4); del módulo~0.9, los intervalos de confianza y el contraste $t$ (§4.4 y §5.3). Del módulo~3 se usan los autovalores y la diagonalización (§3.1 y §3.2); del módulo~4, el jacobiano y la regla de la cadena (tema~2) y la convexidad (tema~5); del módulo~5, la esperanza, la varianza y la estimación de Monte Carlo (§5.3); del módulo~7, la convexidad y el descenso de gradiente (temas~1 y~2) y Adam (§5.3); del módulo~8, las ecuaciones normales de mínimos cuadrados (§2.2), la división entrenamiento-validación-prueba (§1.4) y la regularización \emph{ridge} (§6.1); del módulo~9, el perceptrón multicapa (tema~1); del módulo~10, los adjuntos como multiplicadores de Lagrange (§5.4); del módulo~12, las redes recurrentes y la LSTM (§5.2 y §5.4), y del módulo~21, el proceso de decisión de Markov, la ecuación de Bellman y PPO (temas~1 a~3).

\section*{Organización}
\begin{center}\small
\begin{tabularx}{\textwidth}{@{}c>{\raggedright\arraybackslash}X>{\raggedright\arraybackslash}p{4.1cm}@{}}
\toprule
Tema & Contenido & Fuentes principales \\
\midrule
1 & Procesos dinámicos y control PID & \cite{am}, \cite{seborg}, \cite{mit10450}, \cite{skogestad} \\
2 & Identificación de sistemas y gemelos digitales & \cite{ljung}, \cite{mit6435}, \cite{narendra}, \cite{kritzinger} \\
3 & Control predictivo basado en modelo & \cite{rmd}, \cite{ee363}, \cite{nw}, \cite{lazic} \\
4 & Aprendizaje por refuerzo para control continuo & \cite{sb}, \cite{cs285}, \cite{td3}, \cite{sac}, \cite{ng} \\
5 & Casos de estudio: edificio, reactor, batería y centro de datos & \cite{drgona}, \cite{mit10450}, \cite{evansgao}, \cite{gamblegao} \\
6 & Puesta en producción en planta & \cite{am}, \cite{nist}, \cite{iec62443}, \cite{isa182}, \cite{cs329s} \\
\bottomrule
\end{tabularx}
\end{center}

Cada tema combina teoría formal concisa con al menos un caso de estudio desarrollado con datos numéricos, una tabla de tecnologías y una observación sobre los errores frecuentes en la práctica. Termina con una lista de lecturas y con ejercicios graduados en tres niveles: $\star$ comprensión, $\star\star$ formulación y cálculo a mano y $\star\star\star$ demostración o diseño razonado. Los ejercicios marcados con \emph{Fuente} están adaptados de cursos y libros universitarios; el enunciado se ha redactado de nuevo y, en algunos casos, se han cambiado los datos. Las soluciones completas están en el \cref{ap:soluciones}. El cuaderno \texttt{M26-practica.ipynb} implementa los casos de estudio: el modelo de la zona y su PI (temas~1), la identificación (tema~2), el MPC económico (tema~3), el entorno Gymnasium y el agente (tema~4), la comparación, la batería y el reactor (tema~5) y el filtro de seguridad (tema~6). Todos los resultados numéricos de los casos de climatización proceden de su ejecución.

\section*{Sobre las fuentes}
El texto es de redacción propia. El tema~1 sigue los capítulos~3 («System Modeling»), 6 («Linear Systems») y~11 («PID Control») de Åström y Murray~\cite{am}, el libro de Seborg \emph{et al.}~\cite{seborg} y las lecciones~5 («The Heated Tank») y~6 («The Exothermic Reactor») del curso 10.450 del MIT~\cite{mit10450}; las reglas de ajuste proceden de Skogestad~\cite{skogestad} y de Ziegler y Nichols~\cite{zn}. El tema~2 sigue el libro de Ljung~\cite{ljung}, texto del curso 6.435 del MIT~\cite{mit6435}, cuya hoja de problemas~5 aporta dos ejercicios, y el artículo de Narendra y Parthasarathy~\cite{narendra}; la clasificación de los gemelos digitales es la de Kritzinger \emph{et al.}~\cite{kritzinger}. El tema~3 sigue los capítulos~1 y~2 de Rawlings, Mayne y Diehl~\cite{rmd}, las clases~16 a~18 de EE363 de Stanford~\cite{ee363}, el capítulo~17 de Nocedal y Wright~\cite{nw} para la penalización exacta y la revisión de Drgoňa \emph{et al.}~\cite{drgona} para edificios. El tema~4 sigue Sutton y Barto~\cite{sb}, las clases~7, 8 y~10 de CS285 de Berkeley~\cite{cs285} y su tarea~3~\cite{cs285hw3}, y los artículos originales de DDPG~\cite{ddpg}, TD3~\cite{td3}, SAC~\cite{sac,sac2}, la conformación de recompensas~\cite{ng}, los MDP con restricciones~\cite{altman,cpo} y la aleatorización del dominio~\cite{tobin,peng}. El caso de los centros de datos usa solo información publicada por Google y DeepMind~\cite{evansgao,gamblegao} y el artículo de Lazic \emph{et al.}~\cite{lazic}. El tema~6 se apoya en el capítulo~15 de Åström y Murray~\cite{am}, en la guía NIST SP~800-82~\cite{nist}, en las series IEC~62443~\cite{iec62443}, IEC~61511~\cite{iec61511} e ISA-18.2~\cite{isa182}, en el Reglamento europeo de inteligencia artificial~\cite{aiact}, en Sculley \emph{et al.}~\cite{sculley} y en el curso CS329S de Stanford~\cite{cs329s}. Las normas citadas son documentos de pago; se usan solo sus contenidos difundidos públicamente por los organismos que las publican.

\section*{Notación}
\begin{center}\small
\begin{tabularx}{\textwidth}{@{}l>{\raggedright\arraybackslash}X@{}}
\toprule
Símbolo & Significado \\
\midrule
$x\in\R^{n}$, $u\in\R^{m}$, $w$, $y$ & estado, entrada manipulada, perturbación y salida medida \\
$\dot x=f(x,u,w)$; $x_{k+1}=f_d(x_k,u_k,w_k)$ & modelo en tiempo continuo y en tiempo discreto \\
$A,B,E,C$; $A_d,B_d$ & matrices del modelo lineal continuo y del discretizado \\
$\Dt$, $k$ & periodo de muestreo y índice de instante ($t=k\Dt$) \\
$\rho(A)$ & radio espectral: mayor módulo de los autovalores de $A$ \\
$r$, $e=r-y$ & consigna y error de control \\
$K_c$, $\tau_I$, $\tau_D$ & ganancia, tiempo integral y tiempo derivativo del PID \\
$\theta$ (tema~2) & vector de parámetros de un modelo identificado \\
$N$, $\ell(x,u)$, $V_f$ & horizonte, coste de etapa y coste terminal del MPC \\
$\Xset$, $\Uset$ & conjuntos de estados y entradas admisibles \\
$\Sset$, $\Aset$, $\gamma$, $\pi$ & estados, acciones, descuento y política (como en el módulo~21) \\
$Q^{\pi}$, $Q_\phi$, $\pi_\psi$ & función de valor de acción, crítico y actor con parámetros $\phi$ y $\psi$ \\
$\alpha$ (tema~4) & temperatura de la entropía en SAC \\
$T$, $\Text$, $P$, $G$ & temperaturas interior y exterior (°C), potencia eléctrica y ganancias (kW) \\
$R$, $C$ (casos de estudio) & resistencia (K/kW) y capacidad (kWh/K) térmicas \\
$\COP$ & coeficiente de rendimiento de la bomba de calor \\
\bottomrule
\end{tabularx}
\end{center}

\tableofcontents

\mainmatter
% =====================================================================
\chapter{Procesos dinámicos y control PID}
\label{tema:1}
% =====================================================================

Antes de optimizar un proceso hay que describirlo como sistema dinámico y entender qué hace con él un controlador realimentado. Este tema presenta los modelos en espacio de estados, su discretización exacta y su estabilidad, siguiendo los capítulos~3 y~6 de Åström y Murray~\cite{am}, y el controlador PID con sus reglas de ajuste, según el capítulo~11 del mismo libro, el texto de Seborg \emph{et al.}~\cite{seborg}, la lección~5 del curso 10.450 del MIT~\cite{mit10450} y el método SIMC de Skogestad~\cite{skogestad}. El caso de estudio es la zona de un edificio calentada por una bomba de calor, que reaparece en todo el módulo.

\section{Modelos en espacio de estados}

\begin{definicion}[Sistema en espacio de estados]\label{def:ee}
Un \emph{sistema dinámico en espacio de estados} en tiempo continuo es un par de ecuaciones
\[
\dot x(t)=f\bigl(x(t),u(t),w(t)\bigr),\qquad y(t)=h\bigl(x(t)\bigr),
\]
donde $x(t)\in\R^{n}$ es el \emph{estado} (la información mínima que, junto con las entradas futuras, determina la evolución futura), $u(t)\in\R^{m}$ la \emph{entrada manipulada}, $w(t)$ la \emph{perturbación} (entrada que no se elige) e $y(t)$ la \emph{salida medida}. El sistema es \emph{lineal e invariante en el tiempo} (LTI) si $f(x,u,w)=Ax+Bu+Ew$ y $h(x)=Cx$ con matrices constantes.
\end{definicion}

Los modelos de procesos salen de balances de materia y energía, y casi nunca son lineales. Su comportamiento cerca de un punto de operación se aproxima por uno lineal.

\begin{proposicion}[Linealización]\label{prop:lineal}
Sea $(\bar x,\bar u,\bar w)$ un \emph{punto de equilibrio}, $f(\bar x,\bar u,\bar w)=0$, con $f$ diferenciable con continuidad. Las desviaciones $\delta x=x-\bar x$, $\delta u=u-\bar u$, $\delta w=w-\bar w$ cumplen
\[
\dot{\delta x}=A\,\delta x+B\,\delta u+E\,\delta w+o\bigl(\norm{(\delta x,\delta u,\delta w)}\bigr),
\]
con $A=\partial f/\partial x$, $B=\partial f/\partial u$ y $E=\partial f/\partial w$ evaluados en el equilibrio.
\end{proposicion}

\begin{proof}
Es el desarrollo de Taylor de primer orden de $f$ en el equilibrio (módulo~0.5, tema~6, y módulo~4, tema~2): $f(x,u,w)=f(\bar x,\bar u,\bar w)+A\delta x+B\delta u+E\delta w+o(\cdot)$, y el primer término es nulo. Como $\bar x$ es constante, $\dot{\delta x}=\dot x$.
\end{proof}

\begin{ejemplo}[Modelo RC de una zona]\label{ej:rc}
Una zona de un edificio se modela con una sola capacidad térmica $C$ (kWh/K), que agrupa el aire y la masa que se calienta con él, unida al exterior por una resistencia térmica $R$ (K/kW). Una bomba de calor consume una potencia eléctrica $P$ (kW) y entrega $\COP\cdot P$ de calor, y el sol y los ocupantes aportan unas ganancias $G$ (kW). El balance de energía es
\begin{equation}\label{eq:rc}
C\,\dot T=\frac{\Text-T}{R}+\COP\cdot P+G .
\end{equation}
Con $T$ como estado, $P$ como entrada y $(\Text,G)$ como perturbaciones, es un sistema LTI escalar con $A=-1/(RC)$, $B=\COP/C$ y $E=(1/(RC),\ 1/C)$. La cantidad $\tau=RC$ es su \emph{constante de tiempo}. Con los datos del caso de estudio, $R=5$ K/kW y $C=3$ kWh/K, resulta $\tau=15$ h. En régimen permanente ($\dot T=0$) con $\Text=5$ °C, $G=0$ y $\COP=2{,}8$, mantener $T=21$ °C exige $P=(21-5)/(R\cdot\COP)=16/14=1{,}143$ kW. La \emph{ganancia estática} de $P$ a $T$ es $R\cdot\COP=14$ K/kW.
\end{ejemplo}

\section{Discretización}

Un controlador digital mide cada $\Dt$ y mantiene su salida constante entre dos muestras: es el \emph{retenedor de orden cero} (ZOH). Con esa hipótesis el modelo discreto es exacto.

\begin{proposicion}[Discretización exacta con ZOH]\label{prop:zoh}
Sea $\dot x=Ax+Bu$ con $u(t)=u_k$ constante en $[k\Dt,(k+1)\Dt)$. Entonces $x_{k+1}=A_dx_k+B_du_k$, con
\[
A_d=e^{A\Dt},\qquad B_d=\Bigl(\int_0^{\Dt}e^{A\sigma}\,d\sigma\Bigr)B,
\]
y, si $A$ es invertible, $B_d=A^{-1}\bigl(e^{A\Dt}-I\bigr)B$.
\end{proposicion}

\begin{proof}
La solución de $\dot x=Ax+Bu$ desde $x(t_0)$ es
\[
x(t)=e^{A(t-t_0)}x(t_0)+\int_{t_0}^{t}e^{A(t-s)}Bu(s)\,ds,
\]
como se comprueba derivando: $\frac{d}{dt}e^{A(t-t_0)}=Ae^{A(t-t_0)}$ y la derivada de la integral es $Bu(t)+A\int_{t_0}^{t}e^{A(t-s)}Bu(s)\,ds$. Con $t_0=k\Dt$, $t=(k+1)\Dt$ y $u(s)=u_k$, el cambio $\sigma=(k+1)\Dt-s$ da $\int_0^{\Dt}e^{A\sigma}d\sigma\,Bu_k$. Si $A$ es invertible, $\int_0^{\Dt}e^{A\sigma}d\sigma=A^{-1}(e^{A\Dt}-I)$, porque la derivada de $A^{-1}e^{A\sigma}$ es $e^{A\sigma}$.
\end{proof}

La alternativa más sencilla, el método de \emph{Euler explícito}, sustituye la derivada por un cociente incremental: $x_{k+1}=(I+A\Dt)x_k+B\Dt\,u_k$. Es solo una aproximación de primer orden, y puede volver inestable un sistema estable (\cref{cor:euler}).

\begin{ejemplo}[Discretización de la zona]\label{ej:zoh}
Para el modelo~\eqref{eq:rc} con $\Dt=0{,}25$ h, la \cref{prop:zoh} da $a=e^{-\Dt/\tau}=e^{-1/60}=0{,}983471$ y, para la entrada de calor $Q$, $b=A^{-1}(a-1)\cdot(1/C)=R(1-a)=5\cdot 0{,}016529=0{,}08264$ K/kWh. El modelo discreto de la zona es
\begin{equation}\label{eq:rc-d}
T_{k+1}=a\,T_k+(1-a)\bigl(\Text{}_{,k}+R\,(\COP_k P_k+G_k)\bigr).
\end{equation}
Euler daría $a_E=1-\Dt/\tau=0{,}983333$: la diferencia es pequeña porque $\Dt\ll\tau$.
\end{ejemplo}

\section{Estabilidad}

\begin{definicion}[Estabilidad asintótica]\label{def:estab}
El sistema discreto $x_{k+1}=A_dx_k$ es \emph{asintóticamente estable} si $x_k\to 0$ para todo estado inicial $x_0$. Se llama \emph{radio espectral} de $A_d$ a $\rho(A_d)=\max_i\abs{\lambda_i}$, con $\lambda_i$ sus autovalores (complejos en general).
\end{definicion}

\begin{teorema}[Criterio de estabilidad en tiempo discreto]\label{teo:estab}
Sea $A_d$ diagonalizable. El sistema $x_{k+1}=A_dx_k$ es asintóticamente estable si y solo si $\rho(A_d)<1$.
\end{teorema}

\begin{proof}
Sea $A_d=V\Lambda V^{-1}$ con $\Lambda=\diag(\lambda_1,\dots,\lambda_n)$ (módulo~3, §3.2). Entonces $x_k=A_d^{k}x_0=V\Lambda^{k}V^{-1}x_0$. Si $\rho(A_d)<1$, cada $\lambda_i^{k}\to0$, luego $\Lambda^{k}\to0$ y $x_k\to0$ para todo $x_0$. Recíprocamente, si $\abs{\lambda_j}\ge1$ para algún $j$, se toma $x_0=v_j$, su autovector: $x_k=\lambda_j^{k}v_j$ y $\norm{x_k}=\abs{\lambda_j}^{k}\norm{v_j}\ge\norm{v_j}>0$ no tiende a cero.
\end{proof}

El resultado es cierto también sin diagonalizabilidad (con la forma de Jordan), y su análogo continuo exige $\operatorname{Re}\lambda_i(A)<0$~\cite[cap.~6]{am}.

\begin{corolario}[ZOH conserva la estabilidad; Euler no siempre]\label{cor:euler}
Sea $A$ diagonalizable con autovalores reales negativos $\lambda_i$. (i) $A_d=e^{A\Dt}$ tiene autovalores $e^{\lambda_i\Dt}\in(0,1)$: el modelo ZOH es estable para todo $\Dt>0$. (ii) $I+A\Dt$ tiene autovalores $1+\lambda_i\Dt$, y el modelo de Euler es estable si y solo si $\Dt<2/\max_i\abs{\lambda_i}$.
\end{corolario}

\begin{proof}
Si $A=V\Lambda V^{-1}$, entonces $e^{A\Dt}=Ve^{\Lambda\Dt}V^{-1}$ e $I+A\Dt=V(I+\Lambda\Dt)V^{-1}$, de modo que los autovalores son $e^{\lambda_i\Dt}$ y $1+\lambda_i\Dt$. Para (i), $\lambda_i\Dt<0$ implica $0<e^{\lambda_i\Dt}<1$. Para (ii), $\abs{1+\lambda_i\Dt}<1$ equivale a $-2<\lambda_i\Dt<0$, es decir, $\Dt<2/\abs{\lambda_i}$ para todo $i$; se aplica el \cref{teo:estab}.
\end{proof}

\begin{ejemplo}[Un modelo rígido]\label{ej:rigido}
Un modelo de la zona con dos nodos, aire ($C_a=0{,}5$ kWh/K) y masa ($C_m=5$ kWh/K), unidos por $R_{am}=0{,}5$ K/kW, con el aire unido al exterior por $R_{ao}=5$ K/kW, tiene
\[
A=\begin{pmatrix}-4{,}4&4\\0{,}4&-0{,}4\end{pmatrix}\ \mathrm{h}^{-1},\qquad \lambda_1=-4{,}766,\quad \lambda_2=-0{,}0336,
\]
es decir, constantes de tiempo de $0{,}21$ h y $29{,}8$ h (\cref{eje:1-rigidez}). El sistema es estable, pero Euler explícito solo lo es con $\Dt<2/4{,}766=0{,}420$ h: con $\Dt=0{,}5$ h el autovalor discreto $1-4{,}766\cdot0{,}5=-1{,}383$ produce oscilaciones crecientes que no existen en la planta. Con ZOH los autovalores son $e^{-2{,}383}=0{,}092$ y $e^{-0{,}0168}=0{,}983$. Un sistema con constantes de tiempo muy distintas se llama \emph{rígido}; es la situación habitual en procesos y la razón para usar discretizaciones exactas o integradores implícitos.
\end{ejemplo}

\section{El controlador PID}

\begin{figure}[h]
\centering
\begin{tikzpicture}[>=Stealth,node distance=1.5cm,every node/.style={font=\small},
  bloque/.style={draw,rectangle,minimum height=0.9cm,minimum width=1.9cm,align=center},
  suma/.style={draw,circle,inner sep=1.5pt}]
\node (r) {$r$};
\node[suma,right=0.9cm of r] (s) {$\Sigma$};
\node[bloque,right=1.0cm of s] (c) {controlador\\PID};
\node[bloque,right=1.4cm of c] (p) {proceso};
\node[right=1.1cm of p] (y) {$y$};
\node[above=0.7cm of p] (w) {$w$};
\node[bloque,below=0.6cm of p] (m) {sensor};
\draw[->] (r) -- (s);
\draw[->] (s) -- node[above]{$e$} (c);
\draw[->] (c) -- node[above]{$u$} (p);
\draw[->] (w) -- (p);
\draw[->] (p) -- (y);
\draw[->] ($(p.east)+(0.5,0)$) |- (m);
\draw[->] (m) -| node[pos=0.9,left]{$-$} (s);
\end{tikzpicture}
\caption{Lazo de control realimentado: el controlador actúa sobre el error entre la consigna y la medida.}
\label{fig:lazo}
\end{figure}

\begin{definicion}[Controlador PID]\label{def:pid}
Dado el error $e=r-y$ entre la consigna $r$ y la medida $y$ (\cref{fig:lazo}), el controlador PID calcula
\[
u(t)=K_c\Bigl(e(t)+\frac{1}{\tau_I}\int_0^{t}e(s)\,ds+\tau_D\,\frac{de}{dt}\Bigr),
\]
con ganancia $K_c$, tiempo integral $\tau_I$ y tiempo derivativo $\tau_D$. En forma discreta con periodo $\Dt$ (forma \emph{posicional}),
\[
u_k=K_ce_k+I_k+K_c\tau_D\frac{e_k-e_{k-1}}{\Dt},\qquad I_{k+1}=I_k+\frac{K_c\Dt}{\tau_I}\,e_k .
\]
\end{definicion}

Los tres términos tienen papeles distintos: el proporcional reacciona al error actual, el integral acumula el error pasado y el derivativo anticipa su tendencia. En procesos térmicos lentos y con medidas ruidosas el término derivativo se suele anular (controlador PI) o filtrar~\cite[cap.~11]{am}.

\begin{proposicion}[Error permanente con control proporcional]\label{prop:offset}
Sea una planta de primer orden estable de ganancia estática $k$ y un controlador proporcional $u=\bar u+K_ce$ con $kK_c>-1$. Ante un escalón de consigna de amplitud $\Delta r$, el error permanente es $e_\infty=\Delta r/(1+kK_c)$.
\end{proposicion}

\begin{proof}
En variables de desviación, en régimen permanente $y_\infty=k\,u_\infty=kK_ce_\infty$, y $e_\infty=\Delta r-y_\infty=\Delta r-kK_ce_\infty$. Despejando, $e_\infty(1+kK_c)=\Delta r$. La condición $kK_c>-1$ garantiza que el lazo cerrado de primer orden es estable y que el régimen permanente existe.
\end{proof}

\begin{proposicion}[La acción integral elimina el error permanente]\label{prop:integral}
Sea un lazo con un controlador que contiene la acción integral $I_{k+1}=I_k+(K_c\Dt/\tau_I)e_k$, con $K_c\ne0$. Si, con consigna y perturbaciones constantes, el lazo converge a un equilibrio en el que la salida del controlador no está saturada, entonces el error en ese equilibrio es nulo.
\end{proposicion}

\begin{proof}
En un equilibrio todas las variables son constantes; en particular $I_{k+1}=I_k$, luego $(K_c\Dt/\tau_I)\,e_\infty=0$ y, como $K_c\Dt/\tau_I\ne0$, $e_\infty=0$. La hipótesis de no saturación es necesaria: si la salida está en su límite, el integrador puede seguir cambiando sin que cambie la entrada aplicada y no hay equilibrio del controlador (\cref{eje:1-integral}).
\end{proof}

La proposición no dice nada sobre si el lazo converge: eso depende del ajuste. Por ejemplo, con la zona del \cref{ej:rc} ($k=14$ K/kW) y $K_c=1{,}43$ kW/K, un control solo proporcional dejaría el $1/(1+20)=4{,}8$~\% del escalón de consigna sin corregir; el PI lo elimina.

\begin{definicion}[Saturación y antisaturación]\label{def:windup}
Todo actuador tiene límites: $u\in[u_{\min},u_{\max}]$. Si el error persiste mientras la salida está saturada, el término integral crece sin efecto sobre la planta (\emph{windup}) y, cuando el error cambia de signo, el controlador tarda en salir de la saturación, lo que produce sobreoscilaciones. Un mecanismo de \emph{antisaturación} lo evita; el más sencillo es la \emph{integración condicional}: solo se actualiza $I_k$ si la salida no está saturada o si el error tiende a desaturarla~\cite[cap.~11]{am}.
\end{definicion}

\section{Ajuste del PID}

\begin{definicion}[Modelo de primer orden con retardo]\label{def:fopdt}
Muchos procesos se aproximan por una respuesta de primer orden con retardo (FOPDT), $y(s)/u(s)=k\,e^{-\theta s}/(\tau s+1)$: ganancia $k$, constante de tiempo $\tau$ y retardo $\theta$, que pueden leerse de un ensayo en escalón. En un lazo digital, el muestreo y el retenedor añaden un retardo efectivo del orden de $\Dt/2$ a $\Dt$.
\end{definicion}

Las reglas de Ziegler y Nichols~\cite{zn}, de 1942, fueron las primeras reglas sistemáticas; producen lazos agresivos y poco amortiguados. Las reglas SIMC de Skogestad~\cite{skogestad} parten del FOPDT y de un único parámetro de diseño, la constante de tiempo deseada en lazo cerrado $\tau_c$:
\begin{equation}\label{eq:simc}
K_c=\frac{\tau}{k\,(\tau_c+\theta)},\qquad \tau_I=\min\bigl\{\tau,\ 4(\tau_c+\theta)\bigr\}.
\end{equation}
La elección $\tau_c=\theta$ da un ajuste «firme» con buen compromiso entre rapidez y robustez; valores mayores de $\tau_c$ dan lazos más lentos y robustos. El mínimo en $\tau_I$ evita que, en procesos muy lentos como los térmicos, el integrador tarde demasiado en rechazar perturbaciones.

\begin{ejemplo}[Caso de estudio: control PI de la temperatura de una zona]\label{ej:caso-pi}
\emph{Planta.} La zona del \cref{ej:rc}, con $\COP=2{,}5+0{,}06\,\Text$ (el rendimiento de una bomba de calor aire-agua cae con el frío), $P\in[0,3]$ kW, $\Dt=0{,}25$ h y ruido de medida de $0{,}1$ K. El clima es invernal sintético, con media de $3{,}35$ °C en la semana de prueba. La \emph{banda de confort} es $[20,23]$ °C en horario de ocupación (laborables de 8 a 18~h) y $[17,25]$ °C fuera de él. La tarifa tiene tres periodos: valle a $0{,}10$ €/kWh (de 0 a 8~h y fines de semana), llano a $0{,}15$ y punta a $0{,}22$ (de 10 a 14~h y de 18 a 22~h).

\emph{Ajuste.} Con $\COP(5\,^\circ\mathrm C)=2{,}8$, el FOPDT es $k=14$ K/kW, $\tau=15$ h y $\theta=\Dt=0{,}25$ h. Con $\tau_c=0{,}5$ h, \eqref{eq:simc} da $K_c=15/(14\cdot 0{,}75)=1{,}429$ kW/K y $\tau_I=\min\{15,\,3\}=3$ h. La consigna es $21$ °C de 6 a 18~h en laborables (dos horas de precalentamiento) y $18$ °C el resto; el PI usa integración condicional.

\emph{Resultado} (cuaderno, Parte~2). En la semana de prueba, el PI consume $175{,}4$ kWh, cuesta $21{,}06$ € y no sale de la banda de confort. Es una referencia sólida y sin modelo, pero ciega al precio: calienta igual a las 11~h, en punta, que a las 7~h, en valle. Los temas~3 y~4 parten de esta referencia.
\end{ejemplo}

\begin{table}[htbp]
\centering\footnotesize
\begin{tabularx}{\textwidth}{@{}>{\raggedright\arraybackslash}p{3.1cm}>{\raggedright\arraybackslash}X>{\raggedright\arraybackslash}p{2.0cm}>{\raggedright\arraybackslash}X@{}}
\toprule
Herramienta & Para qué sirve & Licencia o tipo & Cuándo elegirla \\
\midrule
python-control~\cite{pycontrol} & funciones de transferencia, espacio de estados, discretización, respuesta en frecuencia, LQR & libre (BSD-3) & análisis y diseño lineal en Python; docencia \\
\texttt{scipy.signal} & sistemas LTI, \texttt{cont2discrete}, simulación & libre (BSD-3) & cuando basta con lo esencial y no se quiere otra dependencia \\
Modelica / OpenModelica~\cite{modelica} & modelado físico acausal por componentes (térmico, eléctrico, fluidos) & lenguaje abierto; entorno libre & modelos de planta complejos y reutilizables \\
FMI / FMPy~\cite{fmi} & intercambio de modelos (FMU) entre herramientas y cosimulación & estándar abierto; FMPy libre & conectar un modelo de Modelica o de un fabricante con Python \\
MATLAB/Simulink & diseño de control, autoajuste PID, generación de código & comercial & entornos industriales que ya lo usan \\
DCS y PLC (IEC 61131-3) & ejecución en planta de los PID con bloques normalizados & comercial & siempre: es donde vive el control básico \\
\bottomrule
\end{tabularx}
\caption{Tecnologías para el modelado y el control clásico.}
\label{tab:tec1}
\end{table}

\begin{observacion}[Errores frecuentes en la práctica]\label{obs:errores1}
(i) \emph{Ignorar la saturación}: un PI sin antisaturación sobreoscila tras cada arranque. (ii) \emph{Derivar la medida ruidosa}: la acción derivativa sin filtro amplifica el ruido; además, derivar el error en lugar de la medida produce un pico en cada cambio de consigna. (iii) \emph{Discretizar con un paso demasiado grande} con Euler explícito en un modelo rígido (\cref{ej:rigido}). (iv) \emph{Signo de la acción}: en refrigeración, o con válvulas que cierran al aumentar la señal, la ganancia del proceso es negativa y el controlador debe ser de acción inversa. (v) \emph{Ajustar sobre un modelo linealizado lejos del punto de operación}: la ganancia de la zona cambia con el COP, es decir, con la temperatura exterior. (vi) \emph{Evaluar el ajuste con el mismo ensayo} con el que se identificó el FOPDT.
\end{observacion}

\begin{lecturas}
\item K.~J. Åström y R.~M. Murray~\cite{am}, caps.~3 «System Modeling», 6 «Linear Systems» y~11 «PID Control».
\item MIT 10.450~\cite{mit10450}, lección~5 «The Heated Tank».
\item S.~Skogestad~\cite{skogestad}, «Simple analytic rules for model reduction and PID controller tuning».
\item D.~E. Seborg \emph{et al.}~\cite{seborg}, capítulos sobre control realimentado y ajuste de controladores.
\end{lecturas}

\section*{Ejercicios del tema 1}
\addcontentsline{toc}{section}{Ejercicios del tema 1}

\begin{ejercicio}{1}{eje:1-variables}
En un reactor por lotes refrigerado por camisa (\cref{tema:5}) se miden la temperatura del reactor y la de la camisa; se manipula la apertura de la válvula del agua de refrigeración; la temperatura del agua de red varía a lo largo del día, y la concentración de reactivo no se mide. Clasifique cada magnitud como estado, entrada manipulada, perturbación o salida medida. ¿Qué estado no es medible y qué consecuencia tiene para el control?
\end{ejercicio}

\begin{ejercicio}{2}{eje:1-zoh}
Discretice la zona del \cref{ej:rc} ($R=5$ K/kW, $C=3$ kWh/K) con $\Dt=1$ h: calcule $a$ y el coeficiente $b$ de la entrada de calor $Q$ con ZOH, y los coeficientes correspondientes con Euler explícito. ¿Cuál es el error relativo de Euler en cada coeficiente?
\end{ejercicio}

\begin{ejercicio}{2}{eje:1-rigidez}
Para el modelo de dos nodos del \cref{ej:rigido}: (a) escriba los balances de energía de los dos nodos sin entradas y obtenga la matriz $A$; (b) calcule sus autovalores; (c) calcule los autovalores del modelo de Euler con $\Dt=0{,}25$ h y $\Dt=0{,}5$ h, y diga en qué caso es estable; (d) calcule el paso máximo de Euler.
\end{ejercicio}

\begin{ejercicio}{2}{eje:1-tanque}
Un tanque agitado de $V=2$ m³ recibe un caudal total de agua $F=0{,}5$ m³/h (un 60~\% por la corriente~1) y se calienta con vapor a $T_c$ a través de un intercambiador con $UA=1{,}5$ kW/K. Con $\rho c_p=4{,}18$ MJ/(m³·K), el modelo linealizado es $\tau\dot T'=-T'+K_1T_1'+K_2T_2'+K_3T_c'$, con $\tau_R=V/F$, $\beta=UA/(\rho c_pF)$, $\tau=\tau_R/(1+\beta)$, $K_1=(F_1/F)/(1+\beta)$ y $K_3=\beta/(1+\beta)$. (a) Calcule $\tau_R$, $\beta$, $\tau$, $K_1$ y $K_3$. (b) Con un controlador proporcional sobre $T_c$ de ganancia $K_c=2$ K/K, ¿qué fracción de un escalón de consigna queda como error permanente?
\fuente{adaptado de MIT 10.450, primavera de 2006, lección~5 «The Heated Tank», §5.1 y §5.10~\cite{mit10450}.}
\end{ejercicio}

\begin{ejercicio}{2}{eje:1-offset}
La zona del \cref{ej:rc} se controla en torno a $21$ °C con $k=14$ K/kW. (a) Con un controlador P de $K_c=0{,}5$ kW/K, ¿qué error permanente deja un cambio de consigna de $+2$ K? (b) ¿Qué $K_c$ haría falta para dejar menos de $0{,}1$ K? (c) Explique por qué esa ganancia no es una buena solución y qué aporta la acción integral.
\end{ejercicio}

\begin{ejercicio}{2}{eje:1-simc}
Un ensayo en escalón de un intercambiador da un FOPDT con $k=2$ K por cada \% de apertura de válvula, $\tau=10$ min y $\theta=2$ min. (a) Calcule el PI con las reglas SIMC y $\tau_c=\theta$. (b) Repita con $\tau_c=3\theta$. (c) Escriba el PI discreto en forma posicional con $\Dt=0{,}5$ min para el caso~(a) y calcule $u_0$ e $I_1$ si el error inicial es $e_0=1$ K y $I_0=0$.
\fuente{adaptado de S.~Skogestad, «Simple analytic rules for model reduction and PID controller tuning», \emph{Journal of Process Control}, 2003, reglas SIMC para procesos de primer orden con retardo~\cite{skogestad}.}
\end{ejercicio}

\begin{ejercicio}{3}{eje:1-integral}
Considere la planta escalar $x_{k+1}=ax_k+bu_k+w$, con $\abs a<1$, $b>0$, perturbación constante $w$ y salida $y_k=x_k$, controlada por un PI discreto sin saturación. (a) Escriba el sistema en lazo cerrado con estado $(x_k,I_k)$. (b) Demuestre que, si el lazo cerrado es estable, el equilibrio tiene $y=r$ para cualquier $w$. (c) Muestre con un ejemplo que, si $u$ está saturada en $u_{\max}$ y $w$ es tan grande que $ax+bu_{\max}+w=x$ tiene solución $x<r$, el integrador crece indefinidamente sin integración condicional.
\fuente{adaptado de K.~J. Åström y R.~M. Murray, \emph{Feedback Systems}, 2.ª ed., cap.~11 «PID Control», apartados sobre la acción integral y la saturación del integrador~\cite{am}.}
\end{ejercicio}

\begin{ejercicio}{3}{eje:1-estab}
(a) Demuestre que, para $\dot x=-x/\tau$ con $\tau>0$, el modelo de Euler explícito es asintóticamente estable si y solo si $\Dt<2\tau$, y que para $\Dt\in(\tau,2\tau)$ su respuesta cambia de signo en cada paso, cosa que nunca hace la solución exacta. (b) Demuestre que el método de Euler implícito, $x_{k+1}=x_k+\Dt\,Ax_{k+1}$, es estable para todo $\Dt>0$ cuando $A$ es diagonalizable con autovalores reales negativos. (c) ¿Qué método elegiría para simular el modelo rígido del \cref{ej:rigido} dentro del bucle de un MPC, y por qué?
\end{ejercicio}

% =====================================================================
\chapter{Identificación de sistemas y gemelos digitales}
\label{tema:2}
% =====================================================================

El MPC necesita un modelo para predecir y el aprendizaje por refuerzo necesita un simulador en el que equivocarse sin coste. En planta casi nunca se dispone de un modelo físico calibrado: hay que construirlo con datos. Este tema formaliza la \emph{identificación de sistemas} siguiendo a Ljung~\cite{ljung}, texto del curso 6.435 del MIT~\cite{mit6435}, extiende los modelos lineales a los no lineales con redes neuronales como en Narendra y Parthasarathy~\cite{narendra} y define el \emph{gemelo digital} según la clasificación de Kritzinger \emph{et al.}~\cite{kritzinger}. El mensaje central es metodológico: un modelo para control se valida simulando, no prediciendo un paso.

\section{El problema de identificación}

Se dispone de un registro $Z^{N}=\{(u_k,y_k)\}_{k=0}^{N}$ de entradas y salidas. Se elige una \emph{estructura de modelo}, una familia de predictores $\hat y_{k+1}(\theta)$ que dependen de parámetros $\theta$ y de los datos hasta el instante $k$, y se estima $\theta$ minimizando el error de predicción~\cite{ljung}:
\begin{equation}\label{eq:pem}
\hat\theta_N=\arg\min_\theta\ V_N(\theta),\qquad V_N(\theta)=\frac1N\sum_{k=0}^{N-1}\bigl(y_{k+1}-\hat y_{k+1}(\theta)\bigr)^{2}.
\end{equation}

\begin{definicion}[Modelo ARX]\label{def:arx}
Un modelo ARX$(n_a,n_b)$, \emph{autorregresivo con entradas exógenas}, es
\[
y_{k+1}=\sum_{i=1}^{n_a}a_iy_{k+1-i}+\sum_{j=1}^{n_b}b_ju_{k+1-j}+\varepsilon_{k+1},
\]
con $\varepsilon$ ruido blanco. Con el \emph{vector de regresores} $\varphi_k=(y_k,\dots,y_{k+1-n_a},u_k,\dots,u_{k+1-n_b})$ y $\theta=(a,b)$, se escribe $y_{k+1}=\varphi_k\T\theta+\varepsilon_{k+1}$: es una regresión lineal.
\end{definicion}

\begin{proposicion}[Estimador de mínimos cuadrados del ARX]\label{prop:ls}
Sea $\Phi$ la matriz cuyas filas son los $\varphi_k\T$ e $Y$ el vector de los $y_{k+1}$. Si $\Phi\T\Phi$ es invertible, el minimizador de \eqref{eq:pem} es $\hat\theta_N=(\Phi\T\Phi)^{-1}\Phi\T Y$. Si además los datos proceden de un ARX con parámetro $\theta_0$, entonces
\[
\hat\theta_N-\theta_0=\Bigl(\frac1N\sum_k\varphi_k\varphi_k\T\Bigr)^{-1}\frac1N\sum_k\varphi_k\varepsilon_{k+1},
\]
y $\hat\theta_N\to\theta_0$ cuando $\frac1N\sum\varphi_k\varphi_k\T$ converge a una matriz invertible y $\frac1N\sum\varphi_k\varepsilon_{k+1}\to0$.
\end{proposicion}

\begin{proof}
La primera parte son las ecuaciones normales (módulo~8, §2.2). Para la segunda, se sustituye $Y=\Phi\theta_0+\varepsilon$: $\hat\theta_N=(\Phi\T\Phi)^{-1}\Phi\T(\Phi\theta_0+\varepsilon)=\theta_0+(\Phi\T\Phi)^{-1}\Phi\T\varepsilon$, y se multiplica y divide por $N$. El último límite se sigue de la continuidad de la inversa.
\end{proof}

La condición $\frac1N\sum\varphi_k\varepsilon_{k+1}\to0$ se cumple si $\varepsilon$ es blanco, porque $\varphi_k$ solo contiene datos hasta $k$. Si el ruido está coloreado (por ejemplo, ruido de medida sobre una salida autorregresiva), $y_k$ y $\varepsilon_{k+1}$ están correlacionados y el ARX es sesgado; la solución es estimar también un modelo del ruido (ARMAX, error de salida) o usar variables instrumentales~\cite{ljung}. La condición sobre $\frac1N\sum\varphi_k\varphi_k\T$ exige que la entrada sea suficientemente rica.

\begin{definicion}[Excitación persistente]\label{def:pe}
Una señal $u$ es \emph{persistentemente excitante de orden $n$} si la matriz de covarianzas $\bigl(R_u(i-j)\bigr)_{i,j=1}^{n}$, con $R_u(\tau)=\lim_N\frac1N\sum_ku_ku_{k-\tau}$, es definida positiva.
\end{definicion}

Una constante es persistentemente excitante solo de orden~1; una sinusoide, de orden~2; un ruido blanco o una secuencia binaria pseudoaleatoria, de cualquier orden. En planta se usan escalones aleatorios de duración variable, que respetan los límites del actuador y excitan las frecuencias que interesan.

\begin{ejemplo}[Un ARX a mano]\label{ej:arx-mano}
Los datos $u=(1,0,1)$ y $y=(0;\ 0{,}5;\ 0{,}45;\ 0{,}905)$ proceden de $y_{k+1}=0{,}9y_k+0{,}5u_k$. Con $\varphi_k=(y_k,u_k)$:
\[
\Phi\T\Phi=\begin{pmatrix}0{,}4525&0{,}45\\0{,}45&2\end{pmatrix},\qquad \Phi\T Y=\begin{pmatrix}0{,}63225\\1{,}405\end{pmatrix},\qquad \det=0{,}7025,
\]
y $\hat a=(2\cdot0{,}63225-0{,}45\cdot1{,}405)/0{,}7025=0{,}9$, $\hat b=(0{,}4525\cdot1{,}405-0{,}45\cdot0{,}63225)/0{,}7025=0{,}5$. Sin ruido, tres ecuaciones bastan para recuperar los dos parámetros exactos.
\end{ejemplo}

\section{Modelos no lineales: NARX y redes neuronales}

\begin{definicion}[Modelo NARX]\label{def:narx}
Un modelo NARX es $y_{k+1}=g(\varphi_k;\theta)+\varepsilon_{k+1}$ con los mismos regresores que el ARX y una función $g$ no lineal: un polinomio, una red de base radial o un perceptrón multicapa (módulo~9). Narendra y Parthasarathy~\cite{narendra} propusieron en 1990 identificar y controlar sistemas dinámicos con redes de este tipo.
\end{definicion}

Hay tres familias principales, de menos a más estructura física:
\begin{enumerate}
  \item \emph{Caja negra}: $g$ es un perceptrón multicapa sobre los regresores. Es flexible, pero solo es fiable dentro del rango de los datos.
  \item \emph{Modelos de estado neuronales}: un estado latente $h_{k+1}=F(h_k,u_k)$ e $y_k=H(h_k)$ con $F$ y $H$ redes, es decir, una red recurrente (módulo~12, §5.2) entrenada en simulación sobre secuencias, con retropropagación en el tiempo. Captura dinámicas de orden alto sin elegir $n_a$, a costa de un entrenamiento más delicado.
  \item \emph{Caja gris}: la física aporta la estructura y los datos los parámetros, o una red aprende solo el residuo del modelo físico. Si se conoce la forma no lineal (por ejemplo, un producto entre entrada y perturbación), basta con añadir el regresor correspondiente y el modelo sigue siendo lineal en $\theta$.
\end{enumerate}

\begin{ejemplo}[Modelos de Hammerstein]\label{ej:hammerstein}
Si una no linealidad estática actúa sobre la entrada, $y_{k+1}=ay_k+b\,f(u_k)$ con $f(u)=u+cu^{2}$, el modelo es lineal en $(a,\ b,\ bc)$ con regresores $(y_k,\ u_k,\ u_k^{2})$: se estima por mínimos cuadrados como un sistema de dos entradas y se recupera $c=\widehat{bc}/\hat b$. Solo es identificable el producto $b\,f$: por eso se fija a~1 el coeficiente lineal de $f$ (\cref{eje:2-hammerstein}).
\end{ejemplo}

\section{Validación: predicción a un paso y simulación libre}

\begin{definicion}[Dos formas de validar]\label{def:validacion}
Sea $\hat f$ un modelo de un paso, $\hat x_{k+1}=\hat f(x_k,u_k)$. La \emph{predicción a un paso} usa en cada instante el estado medido: $\hat x_{k+1}=\hat f(x_k,u_k)$. La \emph{simulación libre} (validación en lazo abierto del modelo, o \emph{free run}) se realimenta con su propia salida: $\tilde x_{k+1}=\hat f(\tilde x_k,u_k)$, $\tilde x_0=x_0$.
\end{definicion}

El MPC predice $N$ pasos y el agente de RL se entrena durante episodios enteros en el modelo: ambos usan el modelo en simulación libre. La siguiente proposición explica por qué el error a un paso es engañoso.

\begin{proposicion}[Acumulación del error en simulación]\label{prop:acumula}
Sea $x_{k+1}=f(x_k,u_k)$ el sistema real y $\hat f$ un modelo tal que $\abs{\hat f(x,u)-f(x,u)}\le\varepsilon$ para todo $(x,u)$ y $\hat f$ es lipschitziana en $x$ con constante $L$. Entonces el error de simulación libre cumple
\[
\abs{\tilde x_k-x_k}\le\varepsilon\,\frac{1-L^{k}}{1-L}\quad(L\ne1),\qquad\text{y, si }L<1,\quad \abs{\tilde x_k-x_k}\le\frac{\varepsilon}{1-L}.
\]
\end{proposicion}

\begin{proof}
Sea $d_k=\abs{\tilde x_k-x_k}$, con $d_0=0$. Por la desigualdad triangular,
$d_{k+1}=\abs{\hat f(\tilde x_k,u_k)-f(x_k,u_k)}\le\abs{\hat f(\tilde x_k,u_k)-\hat f(x_k,u_k)}+\abs{\hat f(x_k,u_k)-f(x_k,u_k)}\le Ld_k+\varepsilon$.
Por inducción, $d_k\le\varepsilon(1+L+\dots+L^{k-1})=\varepsilon(1-L^{k})/(1-L)$, suma geométrica (módulo~0.6, tema~5), acotada por $\varepsilon/(1-L)$ si $L<1$.
\end{proof}

En procesos lentos $L$ es casi~1. Para la zona, $L=a=0{,}9835$ y $1/(1-L)\approx61$: un error de modelo de $0{,}005$ K por paso, invisible a un paso, puede convertirse en $0{,}30$ K en simulación, y en $0{,}24$ K tras las 96 muestras de un horizonte de 24~h.

\begin{proposicion}[Identificabilidad en lazo cerrado]\label{prop:lazo}
Sea $y_{k+1}=ay_k+bu_k+\varepsilon_{k+1}$ y supóngase que los datos se registran con un controlador proporcional $u_k=-Ky_k$ y sin excitación externa. Entonces $\Phi\T\Phi$ es singular y todos los pares $(a',b')$ con $a'-Kb'=a-Kb$ dan el mismo error de predicción: $a$ y $b$ no son identificables por separado.
\end{proposicion}

\begin{proof}
Las filas de $\Phi$ son $(y_k,\ -Ky_k)$, proporcionales a $(1,-K)$: la segunda columna es $-K$ veces la primera y $\Phi\T\Phi$ tiene rango~1. El predictor vale $a'y_k+b'u_k=(a'-Kb')y_k$, que solo depende de $a'-Kb'$.
\end{proof}

Es la razón por la que no basta con descargar el histórico de un lazo que funciona bien: la entrada es una función de la salida y no aporta información independiente. Las soluciones son añadir excitación (escalones o ruido en la consigna o en la salida del controlador), alternar entre controladores distintos o registrar periodos en manual~\cite{ljung}. El histórico de un PID también es poco informativo por otra razón: el controlador mantiene la planta cerca de un único punto de operación. Lazic \emph{et al.}~\cite{lazic} lo comprobaron en un centro de datos: tres horas de exploración aleatoria y acotada produjeron un modelo mejor que una semana de datos históricos de los PID (\cref{tema:5}).

\begin{observacion}[Partición temporal de los datos]\label{obs:particion}
Las muestras de una serie temporal no son independientes: barajarlas antes de separar entrenamiento y validación (módulo~8, §1.4) filtra información, porque cada muestra de validación tiene vecinas casi idénticas en el entrenamiento. Se separan \emph{bloques de tiempo} (días o semanas completos) y se reserva el último bloque, con condiciones distintas (otra estación del año, otro régimen de producción), como prueba final.
\end{observacion}

\section{El gemelo digital}

\begin{definicion}[Modelo, sombra y gemelo digital]\label{def:gemelo}
Kritzinger \emph{et al.}~\cite{kritzinger} distinguen tres niveles según la automatización del flujo de datos entre el objeto físico y su representación digital: un \emph{modelo digital} intercambia datos con la planta solo de forma manual; una \emph{sombra digital} recibe automáticamente los datos de la planta pero no actúa sobre ella; un \emph{gemelo digital} intercambia datos automáticamente en los dos sentidos.
\end{definicion}

\begin{figure}[h]
\centering
\begin{tikzpicture}[>=Stealth,every node/.style={font=\small},
  caja/.style={draw,rounded corners,minimum height=1cm,minimum width=2.6cm,align=center}]
\node[caja] (planta) {planta\\(sensores, DCS)};
\node[caja,right=2.4cm of planta] (est) {estimación de estado\\y calibración};
\node[caja,right=1.2cm of est] (mod) {modelo\\identificado};
\node[caja,below=1.1cm of est] (mpc) {MPC / política};
\node[caja,below=1.1cm of mod] (usos) {entrenamiento RL,\\análisis «¿y si...?»};
\draw[->] (planta) -- node[above]{medidas} (est);
\draw[->] (est) -- (mod);
\draw[->] (mod) -- (mpc);
\draw[->] (mod) -- (usos);
\draw[->] (mpc) -| node[pos=0.25,below]{consignas} (planta);
\draw[->,dashed] (est.north) to[bend left=20] node[above]{reidentificación} (mod.north);
\end{tikzpicture}
\caption{Flujo de datos de un gemelo digital usado para control.}
\label{fig:gemelo}
\end{figure}

Para control, el gemelo tiene cuatro componentes (\cref{fig:gemelo}): un \emph{modelo} (físico, de datos o híbrido); un \emph{estimador de estado}, porque no todos los estados se miden (el filtro de Kalman o la estimación de horizonte móvil, \cref{tema:3}); un procedimiento de \emph{recalibración} cuando los residuos derivan, y una \emph{interfaz} estándar para que lo usen otros programas: el estándar FMI~\cite{fmi}, que empaqueta un modelo como unidad intercambiable (FMU), o la interfaz de entornos de Gymnasium~\cite{gymnasium}. Para edificios, el simulador EnergyPlus~\cite{crawley} y bibliotecas de Modelica~\cite{modelica} son la referencia física, y marcos como BOPTEST~\cite{boptest} o Sinergym~\cite{sinergym} los envuelven como entornos de prueba para controladores.

\begin{ejemplo}[Caso de estudio: identificación del gemelo de la zona]\label{ej:caso-id}
\emph{Datos.} Se registran 14 días de la zona del \cref{ej:caso-pi} en lazo abierto con escalones aleatorios de potencia (de 1 a 4~h de duración y entre 0 y 3~kW), con ruido de medida de $0{,}1$ K; 1344 muestras. Los días 1 a 10 se usan para estimar y los días 11 a 14 para validar.

\emph{Modelos.} (a) ARX con regresores $(T_k,P_k,\Text{}_{,k},G_k,1)$; (b) NARX con física: añade el regresor $P_k\Text{}_{,k}$, porque por \eqref{eq:rc-d} el coeficiente de $P_k$ es $(1-a)R(2{,}5+0{,}06\,\Text)$; (c) NARX neuronal: un perceptrón de dos capas ocultas de 32 unidades sobre las mismas entradas que el ARX, que predice el incremento de temperatura.

\emph{Resultados} (cuaderno, Parte~3):
\begin{center}\small
\begin{tabular}{@{}lcc@{}}
\toprule
Modelo & RMSE a un paso (K) & RMSE en simulación libre, 4 días (K) \\
\midrule
ARX & $0{,}137$ & $0{,}287$ \\
NARX con física & $0{,}136$ & $0{,}180$ \\
NARX neuronal & $0{,}137$ & $0{,}340$ \\
\bottomrule
\end{tabular}
\end{center}
A un paso los tres son indistinguibles: el error está dominado por el ruido de medida de $T_k$ y $T_{k+1}$, $0{,}1\sqrt{1+a^{2}}=0{,}140$ K. En simulación, el modelo con la estructura correcta es claramente mejor, y recupera los parámetros físicos (por ejemplo, $\hat a=0{,}9827$ frente a $0{,}9835$, es decir, $\hat\tau=14{,}3$ h frente a 15~h). La red neuronal, sin esa estructura y con solo diez días de datos, simula peor que el ARX: más capacidad no compensa menos información. Este modelo con física es el que usa el MPC del \cref{tema:3}.
\end{ejemplo}

\begin{table}[htbp]
\centering\footnotesize
\begin{tabularx}{\textwidth}{@{}>{\raggedright\arraybackslash}p{3.1cm}>{\raggedright\arraybackslash}X>{\raggedright\arraybackslash}p{2.0cm}>{\raggedright\arraybackslash}X@{}}
\toprule
Herramienta & Para qué sirve & Licencia o tipo & Cuándo elegirla \\
\midrule
NumPy / SciPy & mínimos cuadrados, ARX, filtrado & libre (BSD-3) & modelos lineales en los parámetros \\
SysIdentPy & modelos NARMAX, selección de regresores & libre (BSD-3) & NARX polinómicos interpretables \\
scikit-learn & perceptrones, \emph{ridge}, validación & libre (BSD-3) & NARX neuronales pequeños \\
PyTorch & redes recurrentes, modelos de estado neuronales, entrenamiento en simulación & libre (tipo BSD) & modelos neuronales con dinámica latente o híbridos con física \\
MATLAB System Identification Toolbox & ARX, ARMAX, error de salida, subespacios & comercial & entornos ya basados en MATLAB \\
FMI / FMPy, Modelica & encapsular y ejecutar el gemelo & estándar abierto / libre & conectar el gemelo con otras herramientas \\
EnergyPlus~\cite{crawley} & simulación energética detallada de edificios & libre (tipo BSD) & gemelo físico de un edificio; generar datos \\
\bottomrule
\end{tabularx}
\caption{Tecnologías para la identificación y los gemelos digitales.}
\label{tab:tec2}
\end{table}

\begin{observacion}[Errores frecuentes en la práctica]\label{obs:errores2}
(i) \emph{Validar solo a un paso}: todos los modelos parecen buenos (\cref{ej:caso-id}). (ii) \emph{Barajar las muestras} antes de separar entrenamiento y validación (\cref{obs:particion}). (iii) \emph{Identificar con datos en lazo cerrado sin excitación} (\cref{prop:lazo}). (iv) \emph{Extrapolar}: un modelo de caja negra entrenado en invierno no sabe nada del verano, y el optimizador buscará precisamente las regiones donde el modelo se equivoca a su favor; es la versión en control del sobreajuste del modelo sustituto. (v) \emph{Incoherencias de muestreo}: el modelo identificado con $\Dt=15$ min no sirve con $\Dt=5$ min, y los retardos de los registros históricos (marcas de tiempo desalineadas) se convierten en dinámica falsa. (vi) \emph{Fuga de información del futuro}: usar como regresor una medida que en tiempo real aún no existiría, como la temperatura exterior real en lugar de su previsión.
\end{observacion}

\begin{lecturas}
\item L.~Ljung~\cite{ljung}, \emph{System Identification: Theory for the User}, en particular los capítulos sobre estructuras de modelos, métodos de error de predicción y diseño de experimentos.
\item MIT 6.435~\cite{mit6435}, temario y hojas de problemas.
\item K.~S. Narendra y K.~Parthasarathy~\cite{narendra}, «Identification and control of dynamical systems using neural networks».
\item W.~Kritzinger \emph{et al.}~\cite{kritzinger}, «Digital Twin in manufacturing: A categorical literature review and classification».
\end{lecturas}

\section*{Ejercicios del tema 2}
\addcontentsline{toc}{section}{Ejercicios del tema 2}

\begin{ejercicio}{1}{eje:2-unpaso}
Explique con sus palabras por qué un modelo puede tener un error a un paso muy pequeño y, sin embargo, ser inútil para un MPC con horizonte de 24~h. ¿Qué métrica de validación usaría y sobre qué datos?
\end{ejercicio}

\begin{ejercicio}{2}{eje:2-arx}
Con los datos $y=(1;\ 1{,}2;\ 0{,}95;\ 1{,}18)$ y $u=(1,0,1)$, estime por mínimos cuadrados el ARX $y_{k+1}=ay_k+bu_k$: escriba $\Phi\T\Phi$ y $\Phi\T Y$, resuelva el sistema y calcule los residuos. Simule después el modelo estimado desde $y_0=1$ y compare con los datos.
\end{ejercicio}

\begin{ejercicio}{2}{eje:2-escalar}
Sea el proceso $y_k=\theta_0u_k+e_k$, con $e$ ruido blanco gaussiano de varianza $\lambda=1$, independiente de $u$, y el predictor $\hat y_k=\theta u_k$. (a) Obtenga el estimador de mínimos cuadrados $\hat\theta_N$ y su límite cuando $N\to\infty$. (b) Calcule su varianza. (c) Compare la desviación típica con $N=400$ para $u_k=\pm1$ alternante y para $u_k=0{,}1\sin(\omega k)$. ¿Qué enseña sobre la amplitud de la excitación?
\fuente{adaptado de MIT 6.435 \emph{System Identification}, hoja de problemas~5, problema~2~\cite{mit6435}.}
\end{ejercicio}

\begin{ejercicio}{2}{eje:2-hammerstein}
Un actuador tiene una característica no lineal y el proceso sigue $y_{k+1}=ay_k+b\,f(u_k)$ con $f(u)=u+cu^{2}$. (a) Escriba el problema como regresión lineal con dos entradas. (b) Si la estimación da $\hat\theta=(0{,}85;\ 0{,}30;\ -0{,}06)$, obtenga $a$, $b$ y $c$. (c) Explique por qué no se podrían identificar por separado $b$ y los dos coeficientes de $f$ si no se fijara uno de ellos.
\fuente{adaptado de MIT 6.435 \emph{System Identification}, hoja de problemas~5, problema~3 (modelos de tipo Hammerstein)~\cite{mit6435}.}
\end{ejercicio}

\begin{ejercicio}{2}{eje:2-cota}
Un modelo de la zona tiene $L=0{,}98$ y un error de modelo de a lo sumo $\varepsilon=0{,}01$ K por paso. (a) Acote el error de simulación tras 4, 24 y 96 pasos y en el límite. (b) Si el MPC necesita un error menor que $0{,}2$ K en todo el horizonte de 96 pasos, ¿qué $\varepsilon$ máximo puede admitirse? (c) Proponga un margen de seguridad para las restricciones de confort a partir de esa cota.
\end{ejercicio}

\begin{ejercicio}{2}{eje:2-lazo}
Una planta $y_{k+1}=0{,}9y_k+0{,}5u_k+\varepsilon_{k+1}$ funciona con un controlador $u_k=-0{,}4y_k$. (a) ¿Qué modelo de lazo cerrado se observa? (b) Dé tres pares $(a',b')$ que expliquen los datos igual de bien. (c) Si se suma a la salida del controlador una señal $v_k$ blanca, independiente del ruido, de varianza $\sigma_v^{2}>0$, demuestre que $\Phi\T\Phi/N$ converge a una matriz invertible.
\end{ejercicio}

\begin{ejercicio}{3}{eje:2-sesgo}
Sea $x_{k+1}=ax_k+bu_k$ un sistema sin ruido de proceso, con $u$ blanco de varianza $\sigma_u^{2}$ y $\abs a<1$, y sea la medida $y_k=x_k+v_k$ con $v$ blanco de varianza $\sigma_v^{2}$, independiente de $u$. Se estima el ARX $y_{k+1}=\alpha y_k+\beta u_k$. (a) Demuestre que $\hat\alpha$ converge a $a\,\sigma_x^{2}/(\sigma_x^{2}+\sigma_v^{2})$, con $\sigma_x^{2}=b^{2}\sigma_u^{2}/(1-a^{2})$ la varianza estacionaria de $x$. (b) Interprete el resultado (atenuación por error en las variables) y calcule el sesgo de la constante de tiempo para $a=0{,}9$, $b=1$, $\sigma_u=1$ y $\sigma_v=1$.
\end{ejercicio}

\begin{ejercicio}{3}{eje:2-diseno}
Diseñe el protocolo de identificación y validación del gemelo digital de una nave industrial climatizada de la que se tiene un año de datos históricos de su PID y se puede intervenir dos semanas. Indique: experimentos y excitación (con sus límites de seguridad), partición de los datos, estructuras de modelo candidatas, métricas, criterio de aceptación para usarlo en un MPC y procedimiento de recalibración.
\end{ejercicio}

% =====================================================================
\chapter{Control predictivo basado en modelo}
\label{tema:3}
% =====================================================================

El PID corrige el error presente; no sabe que mañana a las 10 la electricidad será cara, ni que la potencia de la bomba de calor está limitada. El \emph{control predictivo basado en modelo} (MPC) usa un modelo para predecir el futuro, resuelve en cada instante un problema de optimización con restricciones sobre un horizonte y aplica solo la primera decisión. Este tema sigue los capítulos~1 y~2 de Rawlings, Mayne y Diehl~\cite{rmd}, las clases~16 («Linear quadratic regulator: Discrete-time finite horizon») a~18 de EE363 de Stanford~\cite{ee363} y la clase~10 de CS285~\cite{cs285}, que presenta el control óptimo y la planificación como punto de partida del RL basado en modelos. La aplicación a edificios sigue la revisión de Drgoňa \emph{et al.}~\cite{drgona}.

\section{Control óptimo de horizonte finito: el regulador lineal cuadrático}

\begin{definicion}[Problema de control óptimo de horizonte finito]\label{def:ocp}
Dados el modelo $x_{k+1}=f(x_k,u_k)$, un coste de etapa $\ell(x,u)$, un coste terminal $V_f(x)$, conjuntos admisibles $\Xset$, $\Uset$, $\Xset_f$ y un estado inicial $x$, el problema de horizonte $N$ es
\[
\begin{aligned}
\mathbb P_N(x):\quad V_N^{0}(x)=\min_{u_0,\dots,u_{N-1}}\ &\sum_{j=0}^{N-1}\ell(x_j,u_j)+V_f(x_N)\\
\text{s.a.}\quad &x_0=x,\quad x_{j+1}=f(x_j,u_j),\quad x_j\in\Xset,\quad u_j\in\Uset,\quad x_N\in\Xset_f .
\end{aligned}
\]
\end{definicion}

Sin restricciones, con modelo lineal y costes cuadráticos, el problema tiene solución explícita.

\begin{teorema}[LQR de horizonte finito]\label{teo:lqr}
Sean $x_{k+1}=Ax_k+Bu_k$, $\ell(x,u)=x\T Qx+u\T Ru$ y $V_f(x)=x\T P_fx$, con $Q\succeq0$, $P_f\succeq0$ y $R\succ0$, sin restricciones. Se define hacia atrás $P_N=P_f$ y, para $k=N-1,\dots,0$,
\[
K_k=(R+B\T P_{k+1}B)^{-1}B\T P_{k+1}A,\qquad P_k=Q+A\T P_{k+1}A-A\T P_{k+1}B\,K_k .
\]
Entonces el coste óptimo desde $x$ en el instante $k$ es $V_k(x)=x\T P_kx$ y la decisión óptima es $u_k=-K_kx_k$.
\end{teorema}

\begin{proof}
Por inducción hacia atrás con la ecuación de Bellman $V_k(x)=\min_u\{\ell(x,u)+V_{k+1}(Ax+Bu)\}$ (módulo~21, tema~1). Para $k=N$, $V_N=x\T P_fx$. Si $V_{k+1}(z)=z\T P_{k+1}z$ con $P_{k+1}\succeq0$, la función a minimizar es
\[
\begin{aligned}
&x\T Qx+u\T Ru+(Ax+Bu)\T P_{k+1}(Ax+Bu)\\
&\qquad=x\T(Q+A\T P_{k+1}A)x+2u\T B\T P_{k+1}Ax+u\T(R+B\T P_{k+1}B)u .
\end{aligned}
\]
Como $M=R+B\T P_{k+1}B\succ0$, es estrictamente convexa en $u$ y su gradiente $2B\T P_{k+1}Ax+2Mu$ se anula en $u=-M^{-1}B\T P_{k+1}Ax=-K_kx$. Sustituyendo, el valor mínimo es $x\T(Q+A\T P_{k+1}A-A\T P_{k+1}BM^{-1}B\T P_{k+1}A)x=x\T P_kx$. Además $P_k\succeq0$, porque es el mínimo de una suma de términos no negativos.
\end{proof}

\begin{ejemplo}[LQR escalar]\label{ej:lqr}
Sea $x_{k+1}=x_k+u_k$ con $Q=R=P_f=1$ y $N=2$. Hacia atrás: $P_2=1$; $K_1=1/(1+1)=0{,}5$, $P_1=1+1-0{,}5=1{,}5$; $K_0=1{,}5/2{,}5=0{,}6$, $P_0=1+1{,}5-1{,}5\cdot0{,}6=1{,}6$. Desde $x_0=1$: $u_0=-0{,}6$, $x_1=0{,}4$, $u_1=-0{,}2$, $x_2=0{,}2$, y el coste es $1+0{,}36+0{,}16+0{,}04+0{,}04=1{,}6=P_0x_0^{2}$. Con horizonte infinito, el punto fijo $P=1+P-P^{2}/(1+P)$ da $P^{2}-P-1=0$, $P=(1+\sqrt5)/2=1{,}618$ y $K=P/(1+P)=0{,}618$.
\end{ejemplo}

\section{El horizonte deslizante}

\begin{definicion}[Control predictivo]\label{def:mpc}
En cada instante $k$ se mide (o se estima) el estado $x_k$, se resuelve $\mathbb P_N(x_k)$, se aplica la primera decisión $u_k=u_0^{*}(x_k)$, se descarta el resto y se repite en $k+1$. La ley de control implícita $\kappa_N(x)=u_0^{*}(x)$ es una realimentación del estado.
\end{definicion}

\begin{figure}[h]
\centering
\begin{tikzpicture}[>=Stealth,every node/.style={font=\small},xscale=1.05,yscale=0.9]
\draw[->] (0,0) -- (10.6,0) node[right]{tiempo};
\draw[->] (0,0) -- (0,3.2);
\draw[dashed] (3,0) -- (3,3) node[above]{$k$};
\draw[dashed] (9,0) -- (9,3) node[above]{$k+N$};
\draw[thick] (0.2,0.8) -- (0.8,1.0) -- (1.4,1.4) -- (2.0,1.5) -- (2.6,1.8) -- (3,1.9);
\node[above] at (1.4,1.55) {pasado medido};
\draw[thick,dashed] (3,1.9) -- (4,2.2) -- (5,2.35) -- (6,2.4) -- (7,2.45) -- (8,2.5) -- (9,2.5);
\node[above] at (6.3,2.5) {predicción $\hat x_{j}$};
\draw[fill=black!25] (3,0.3) rectangle (4,0.9);
\draw (4,0.3) rectangle (5,0.7); \draw (5,0.3) rectangle (6,0.6); \draw (6,0.3) rectangle (7,0.5);
\draw (7,0.3) rectangle (8,0.45); \draw (8,0.3) rectangle (9,0.42);
\node[below] at (3.5,0) {$u_0^{*}$};
\node[below] at (6.5,0) {secuencia optimizada};
\draw[->] (3.5,1.0) to[bend left=30] node[above,pos=0.6]{se aplica} (4.3,1.55);
\end{tikzpicture}
\caption{Horizonte deslizante: se optimiza una secuencia de $N$ decisiones, se aplica la primera y se repite con la nueva medida.}
\label{fig:horizonte}
\end{figure}

La realimentación procede de volver a resolver con cada nueva medida: si una perturbación desvía la planta de la predicción, el siguiente problema parte del estado real. Que la ley resultante estabilice la planta no es automático, porque el controlador optimiza un horizonte finito que se desplaza. El resultado clásico pide unos ingredientes terminales~\cite[cap.~2]{rmd}.

\begin{teorema}[Factibilidad recursiva y estabilidad]\label{teo:mpc}
Supóngase que $f(0,0)=0$, que $\ell(x,u)\ge c\norm{x}^{2}$ con $c>0$, y que existen un conjunto terminal $\Xset_f\subseteq\Xset$ y una ley local $\kappa_f$ tales que, para todo $x\in\Xset_f$: $\kappa_f(x)\in\Uset$, $f(x,\kappa_f(x))\in\Xset_f$ y $V_f\bigl(f(x,\kappa_f(x))\bigr)-V_f(x)\le-\ell\bigl(x,\kappa_f(x)\bigr)$. Si $\mathbb P_N(x_0)$ es factible, entonces: (i) $\mathbb P_N(x_k)$ es factible para todo $k$; (ii) $V_N^{0}(x_{k+1})\le V_N^{0}(x_k)-\ell\bigl(x_k,\kappa_N(x_k)\bigr)$, y (iii) $x_k\to0$.
\end{teorema}

\begin{proof}
Sea $(u_0^{*},\dots,u_{N-1}^{*})$ óptima desde $x_k$, con estados $x_0^{*}=x_k,\dots,x_N^{*}\in\Xset_f$. Tras aplicar $u_0^{*}$, $x_{k+1}=x_1^{*}$. La secuencia desplazada $(u_1^{*},\dots,u_{N-1}^{*},\kappa_f(x_N^{*}))$ es factible desde $x_{k+1}$: respeta las restricciones hasta $x_N^{*}$ y, por la invariancia de $\Xset_f$, lleva a $f(x_N^{*},\kappa_f(x_N^{*}))\in\Xset_f$. Esto prueba (i). Su coste es
\[
V_N^{0}(x_k)-\ell(x_k,u_0^{*})-V_f(x_N^{*})+\ell\bigl(x_N^{*},\kappa_f(x_N^{*})\bigr)+V_f\bigl(f(x_N^{*},\kappa_f(x_N^{*}))\bigr)\le V_N^{0}(x_k)-\ell(x_k,u_0^{*}),
\]
por la condición de decrecimiento de $V_f$; el óptimo desde $x_{k+1}$ no es mayor que el coste de esta secuencia, lo que prueba (ii). Sumando (ii) de $0$ a $K-1$, $c\sum_{k<K}\norm{x_k}^{2}\le\sum_{k<K}\ell(x_k,u_k)\le V_N^{0}(x_0)-V_N^{0}(x_K)\le V_N^{0}(x_0)$, porque $V_N^{0}\ge0$. La serie $\sum\norm{x_k}^{2}$ converge, luego $x_k\to0$.
\end{proof}

En el caso lineal-cuadrático, la elección habitual es $V_f(x)=x\T Px$ con $P$ la solución de la ecuación de Riccati de horizonte infinito y $\kappa_f=-Kx$ el LQR, y $\Xset_f$ un conjunto en el que el LQR no viola ninguna restricción. En la práctica industrial se usan a menudo horizontes largos sin ingredientes terminales explícitos, y la estabilidad se comprueba por simulación.

\section{El MPC lineal con restricciones como programa cuadrático}

Con modelo lineal, la trayectoria predicha es lineal en las decisiones. Sean $U=(u_0,\dots,u_{N-1})$ y $X=(x_1,\dots,x_N)$. Por recursión, $x_j=A^{j}x_0+\sum_{i<j}A^{j-1-i}Bu_i$, es decir,
\begin{equation}\label{eq:condensada}
X=\Phi x_0+\Gamma U,\qquad \Phi=\begin{pmatrix}A\\A^{2}\\\vdots\\A^{N}\end{pmatrix},\qquad \Gamma=\begin{pmatrix}B&0&\cdots&0\\AB&B&\cdots&0\\\vdots&&\ddots&\\A^{N-1}B&A^{N-2}B&\cdots&B\end{pmatrix}.
\end{equation}

\begin{proposicion}[El MPC lineal-cuadrático es un QP estrictamente convexo]\label{prop:qp}
Sean $\bar Q=\diag(Q,\dots,Q,P_f)\succeq0$ y $\bar R=\diag(R,\dots,R)\succ0$. El coste $\sum_{j<N}(x_j\T Qx_j+u_j\T Ru_j)+x_N\T P_fx_N$ vale
\[
J(U)=U\T HU+2x_0\T F\T U+x_0\T(Q+\Phi\T\bar Q\Phi)x_0,\qquad H=\Gamma\T\bar Q\Gamma+\bar R,\quad F=\Gamma\T\bar Q\Phi,
\]
y $H\succ0$. Con restricciones lineales sobre $U$ y $X$ (límites de actuadores, de velocidad de cambio, de estados), el problema es un programa cuadrático estrictamente convexo con, a lo sumo, una solución.
\end{proposicion}

\begin{proof}
Se sustituye \eqref{eq:condensada} en $X\T\bar QX+U\T\bar RU+x_0\T Qx_0$ y se agrupan términos. Para todo $v\ne0$, $v\T Hv=\norm{\bar Q^{1/2}\Gamma v}^{2}+v\T\bar Rv\ge v\T\bar Rv>0$. Las restricciones lineales en $X$ son lineales en $U$ por \eqref{eq:condensada}, de modo que el conjunto factible es un poliedro convexo; una función estrictamente convexa tiene a lo sumo un minimizador en un convexo (módulo~4, tema~5).
\end{proof}

\begin{ejemplo}[MPC con restricción en la entrada]\label{ej:mpc-qp}
En el \cref{ej:lqr} se añade $\abs{u_j}\le0{,}5$. Con $x_0=1$, $\Phi=(1,1)\T$ y $\Gamma=\bigl(\begin{smallmatrix}1&0\\1&1\end{smallmatrix}\bigr)$, el coste es $J=3+2(2u_0+u_1)+3u_0^{2}+2u_0u_1+2u_1^{2}$, con $H=\bigl(\begin{smallmatrix}3&1\\1&2\end{smallmatrix}\bigr)$. El óptimo sin restricciones, $(-0{,}6;\ -0{,}2)$, viola la cota. Fijando $u_0=-0{,}5$ y minimizando en $u_1$, $\partial J/\partial u_1=2+2u_0+4u_1=0$ da $u_1=-0{,}25$, con $J=1{,}625$. Es el óptimo: en ese punto $\partial J/\partial u_0=4+6u_0+2u_1=0{,}5>0$, de modo que aumentar $u_0$ empeora el coste y disminuirlo no está permitido; $0{,}5$ es el multiplicador de la restricción activa.
\end{ejemplo}

\begin{definicion}[Restricciones blandas]\label{def:blandas}
Una restricción de estado $g(x_j)\le0$ puede volver infactible el problema si una perturbación saca la planta del conjunto admisible. Se \emph{ablanda} con una holgura $s_j\ge0$: $g(x_j)\le s_j$, y se añade $\rho\sum_js_j$ al coste.
\end{definicion}

\begin{proposicion}[Penalización exacta]\label{prop:exacta}
Sea el problema convexo $\min f(z)$ s.a. $g_i(z)\le0$, con solución $z^{*}$ y multiplicadores de Karush--Kuhn--Tucker $\lambda_i^{*}\ge0$. Si $\rho\ge\max_i\lambda_i^{*}$, entonces $z^{*}$ minimiza también $f(z)+\rho\sum_i\max\{g_i(z),0\}$, que es la versión con holguras óptimas.
\end{proposicion}

\begin{proof}
Para todo $z$, como $\rho\ge\lambda_i^{*}\ge0$ y $\max\{g_i,0\}\ge g_i$,
\[
f(z)+\rho\sum_i\max\{g_i(z),0\}\ge f(z)+\sum_i\lambda_i^{*}g_i(z)=L(z,\lambda^{*})\ge L(z^{*},\lambda^{*})=f(z^{*}).
\]
La segunda desigualdad se debe a que $L(\cdot,\lambda^{*})$ es convexa y su gradiente se anula en $z^{*}$ (estacionariedad de KKT); la igualdad final, a la holgura complementaria $\lambda_i^{*}g_i(z^{*})=0$. En $z=z^{*}$ todas las penalizaciones son nulas y se alcanza la cota.
\end{proof}

Es la versión convexa del teorema~17.3 de Nocedal y Wright~\cite[§17.2]{nw}. En la práctica: si $\rho$ es suficientemente grande, las holguras solo se activan cuando el problema con restricciones duras es infactible; si es pequeño, el optimizador «compra» violaciones de confort cuando le salen baratas. En el \cref{ej:mpc-qp}, con la restricción $u_0\ge-0{,}5$ ablandada, cualquier $\rho\ge0{,}5$ reproduce la solución dura.

\begin{definicion}[MPC económico]\label{def:empc}
Cuando el coste de etapa es un coste económico (euros de energía, producción, emisiones) en lugar de una desviación respecto a una consigna, se habla de \emph{MPC económico}. Si el modelo es lineal y el coste y las restricciones son lineales, el problema es un programa lineal.
\end{definicion}

\section{Cálculo en tiempo real y modelos no lineales}

Cada problema debe resolverse en una fracción del periodo de muestreo. Los QP de MPC lineal se resuelven con métodos de punto interior, de conjunto activo o de separación de operadores como OSQP~\cite{osqp}, que se benefician de arrancar desde la solución anterior desplazada (\emph{warm start}). Con modelo no lineal, $\mathbb P_N$ es un programa no lineal (NMPC) y no convexo en general. Las técnicas estándar~\cite[cap.~8]{rmd} son el \emph{disparo múltiple} (los estados de cada paso son también variables, con la dinámica como restricción de igualdad, lo que mejora el condicionamiento) y la \emph{iteración en tiempo real}: un solo paso de programación cuadrática secuencial por instante, que aprovecha que el problema cambia poco de un instante al siguiente. CasADi~\cite{casadi} proporciona la diferenciación automática y la formulación; acados~\cite{acados} y do-mpc~\cite{dompc} construyen sobre ellas resolvedores rápidos y herramientas completas de MPC y de \emph{estimación de horizonte móvil}, que es el problema dual de estimar el estado a partir de las últimas medidas.

\section{MPC con modelo aprendido}

El modelo del MPC puede ser el identificado en el \cref{tema:2}:
\begin{enumerate}
  \item \emph{Modelo lineal en las entradas} (ARX o NARX lineal en $u$ con coeficientes que dependen de perturbaciones previstas): el MPC sigue siendo un QP o un LP. Es la opción preferida cuando basta.
  \item \emph{Red neuronal}: $\hat f_\theta$ entra en las restricciones de igualdad del NLP. Los gradientes que necesita el optimizador son productos vector-jacobiano de la red, que proporciona la retropropagación (módulo~10, §5.4, donde se ve que los adjuntos son multiplicadores de Lagrange de la dinámica). Lazic \emph{et al.}~\cite{lazic} optimizaron así un MPC con modelo ARX directamente con Adam (módulo~7, §5.3) en TensorFlow, parametrizando las acciones con una tangente hiperbólica para respetar sus límites.
  \item \emph{Robustez frente al error de modelo}: la \cref{prop:acumula} permite \emph{endurecer} las restricciones con un margen que crece con el horizonte, $\underline T+\varepsilon(1-L^{j})/(1-L)$ en el paso $j$; un margen constante es la versión más simple.
\end{enumerate}

\begin{ejemplo}[Caso de estudio: arbitraje con una batería]\label{ej:bateria-lp}
\emph{Datos.} Una batería de $1$ kWh y $1$ kW, vacía al inicio, ve los precios $(0{,}10;\ 0{,}30;\ 0{,}05;\ 0{,}25)$ €/kWh en cuatro horas consecutivas. Sean $c_t,d_t\ge0$ la energía cargada y descargada en la hora $t$ (lado red) y $e_t$ la almacenada, con rendimientos de carga y descarga $\eta_c=\eta_d=\eta$.

\emph{Formulación.} $\max\sum_tp_t(d_t-c_t)$ s.a. $e_{t+1}=e_t+\eta c_t-d_t/\eta$, $0\le e_t\le1$, $0\le c_t,d_t\le1$, $e_0=0$: un programa lineal.

\emph{Solución.} Con $\eta=1$ conviene cargar en las horas~1 y~3 y descargar en las~2 y~4: beneficio $(0{,}30-0{,}10)+(0{,}25-0{,}05)=0{,}40$ €. Con $\eta=0{,}9$, cargar $1$ kWh de red almacena $0{,}9$ kWh y devuelve $0{,}81$ kWh: el primer ciclo deja $0{,}30\cdot0{,}81-0{,}10=0{,}143$ € y el segundo $0{,}25\cdot0{,}81-0{,}05=0{,}1525$ €, en total $0{,}2955$ €. Un ciclo solo es rentable si el cociente entre precios de venta y de compra supera $1/\eta^{2}=1{,}235$ (\cref{eje:3-bateria}).

\emph{Escala real} (cuaderno, Parte~8). Con $10$ kWh, $5$ kW, $\eta=0{,}95$, un coste de degradación de $0{,}01$ €/kWh y 24 precios horarios, el LP da un beneficio de $1{,}884$ € y $1{,}90$ ciclos equivalentes en el día.
\end{ejemplo}

\begin{ejemplo}[Caso de estudio: MPC económico de la zona]\label{ej:caso-mpc}
\emph{Formulación.} Con el modelo identificado del \cref{ej:caso-id}, $\hat T_{j+1}=\theta_1\hat T_j+b_jP_j+d_j$, donde $b_j=\theta_2+\theta_3\Text{}_{,j}$ y $d_j$ agrupa los términos de $\Text$, $G$ y la constante, y con previsiones de 24~h ($N=96$) de temperatura exterior, ganancias y precio, en cada instante se resuelve
\[
\min_{P,s}\ \Dt\sum_{j=0}^{N-1}(c_jP_j+\rho\,s_j)\quad\text{s.a.}\quad 0\le P_j\le3,\quad \underline T_{j+1}+m-s_j\le\hat T_{j+1}\le\overline T_{j+1}-m+s_j,\quad s_j\ge0,
\]
con $\rho=2$ €/(K·h) y un margen $m=0{,}3$ K. Es un LP con 289 variables.

\emph{Elección del método.} El modelo es lineal en $P$ y las previsiones son conocidas: un LP se resuelve en milisegundos y su óptimo es global; no hay motivo para un NMPC.

\emph{Resultados} (cuaderno, Parte~4). El MPC consume $171{,}1$ kWh, algo menos que el PI ($175{,}4$), y cuesta $19{,}28$ € frente a $21{,}06$ €, un $8{,}4$~\% menos, sin violar el confort. El ahorro no procede de calentar menos, sino de calentar \emph{antes}: precalienta en valle hasta el borde superior admisible y deja que la temperatura descienda durante la punta, usando la masa térmica como batería. Sin margen ($m=0$) el coste baja a $18{,}76$ €, pero el error de modelo y el ruido producen $3{,}54$ K·h de violación del confort. Cada LP tarda unos 4~ms con Clarabel.
\end{ejemplo}

\begin{table}[htbp]
\centering\footnotesize
\begin{tabularx}{\textwidth}{@{}>{\raggedright\arraybackslash}p{3.1cm}>{\raggedright\arraybackslash}X>{\raggedright\arraybackslash}p{2.0cm}>{\raggedright\arraybackslash}X@{}}
\toprule
Herramienta & Para qué sirve & Licencia o tipo & Cuándo elegirla \\
\midrule
CVXPY~\cite{cvxpy} & lenguaje de modelado de problemas convexos; parámetros para resolver repetidamente & libre (Apache-2.0) & MPC lineal (LP, QP) en Python y prototipos \\
OSQP~\cite{osqp}, Clarabel, HiGHS & resolvedores de QP, cónicos y de LP & libres (Apache-2.0, MIT) & bajo CVXPY o embebidos \\
CasADi~\cite{casadi} & diferenciación automática simbólica, NLP, integradores, interfaz a IPOPT & libre (LGPL-3) & NMPC y estimación con modelos físicos o neuronales \\
acados~\cite{acados} & NMPC y estimación en tiempo real (iteración en tiempo real, HPIPM), generación de código C & libre (BSD-2) & sistemas embebidos y muestreo rápido \\
do-mpc~\cite{dompc} & MPC no lineal, robusto multietapa y estimación de horizonte móvil sobre CasADi & libre (LGPL-3) & prototipos completos de NMPC en Python \\
Pyomo & modelado algebraico de LP, MILP y NLP & libre (BSD-3) & planificación con variables enteras \\
MPC comerciales de los fabricantes de DCS & MPC multivariable industrial integrado con el sistema de control & comercial & refinerías y plantas químicas con soporte del fabricante \\
\bottomrule
\end{tabularx}
\caption{Tecnologías para el control predictivo.}
\label{tab:tec3}
\end{table}

\begin{observacion}[Errores frecuentes en la práctica]\label{obs:errores3}
(i) \emph{Horizonte demasiado corto}: si no ve el siguiente periodo caro, el MPC no precalienta. (ii) \emph{Restricciones de estado duras}: una perturbación vuelve infactible el problema y el controlador se queda sin respuesta; hay que ablandarlas y prever una acción de respaldo. (iii) \emph{Optimismo del modelo}: el optimizador lleva la planta al borde de las restricciones, justo donde el error de modelo cuesta caro; se necesitan márgenes (\cref{ej:caso-mpc}). (iv) \emph{Previsiones perfectas en la evaluación}: comparar un MPC con previsiones reales de la meteorología frente a un PID sobrestima el beneficio si en la simulación se le dan las observadas. (v) \emph{Tiempo de cálculo}: un NLP que a veces no converge en el periodo de muestreo necesita un límite de iteraciones y una solución de reserva (la anterior desplazada). (vi) \emph{Mal escalado} de variables y restricciones, que degrada la precisión de los resolvedores.
\end{observacion}

\begin{lecturas}
\item J.~B. Rawlings, D.~Q. Mayne y M.~M. Diehl~\cite{rmd}, caps.~1 «Getting Started with Model Predictive Control», 2 «Model Predictive Control: Regulation» y~8 «Numerical Optimal Control».
\item Stanford EE363~\cite{ee363}, clases~16 a~18 sobre el regulador lineal cuadrático.
\item J.~Drgoňa \emph{et al.}~\cite{drgona}, «All you need to know about model predictive control for buildings».
\item N.~Lazic \emph{et al.}~\cite{lazic}, «Data center cooling using model-predictive control».
\end{lecturas}

\section*{Ejercicios del tema 3}
\addcontentsline{toc}{section}{Ejercicios del tema 3}

\begin{ejercicio}{1}{eje:3-deslizante}
Explique por qué el MPC aplica solo la primera decisión del plan óptimo aunque el plan entero sea óptimo para el modelo. ¿Qué ocurriría si aplicara las $N$ decisiones en lazo abierto y la previsión de temperatura exterior se equivocara en $2$ K?
\end{ejercicio}

\begin{ejercicio}{2}{eje:3-lqr}
Para $x_{k+1}=2x_k+u_k$ con $Q=1$, $R=1$, $P_f=0$ y $N=2$: (a) calcule $P_1$, $K_1$, $P_0$ y $K_0$ con la recursión de Riccati; (b) calcule la secuencia óptima y el coste desde $x_0=1$ y compruebe que coincide con $P_0x_0^{2}$; (c) obtenga $P$ y $K$ de horizonte infinito y el autovalor del lazo cerrado.
\fuente{adaptado de Stanford EE363, clase~16 «Linear quadratic regulator: Discrete-time finite horizon»~\cite{ee363}.}
\end{ejercicio}

\begin{ejercicio}{2}{eje:3-condensada}
Para $x_{k+1}=ax_k+bu_k$ con $a=0{,}9$, $b=0{,}5$ y $N=3$, escriba $\Phi$ y $\Gamma$ de \eqref{eq:condensada}. Con $Q=1$, $R=0{,}1$ y $P_f=1$, calcule $H$ y compruebe que es definida positiva.
\end{ejercicio}

\begin{ejercicio}{2}{eje:3-bateria}
(a) Demuestre que, en el \cref{ej:bateria-lp}, comprar $1$ kWh a precio $p_c$ y venderlo después a $p_v$ con rendimientos $\eta$ de carga y descarga es rentable si y solo si $p_v/p_c>1/\eta^{2}$. (b) Con $\eta=0{,}9$ y los precios del ejemplo, compruebe el beneficio de $0{,}2955$ € y diga qué pasaría si el segundo par de precios fuera $(0{,}20;\ 0{,}24)$. (c) Añada un coste de degradación de $0{,}02$ € por kWh cargado o descargado (lado red) y recalcule.
\end{ejercicio}

\begin{ejercicio}{2}{eje:3-precalentar}
La zona del \cref{ej:rc} ($\tau=15$ h, $\Dt=0{,}25$ h) puede precalentarse a las 6~h en valle ($0{,}10$ €/kWh) para necesitar menos calor a las 10~h, en punta ($0{,}22$ €/kWh). (a) ¿Qué fracción de un exceso de temperatura conseguido a las 6~h queda a las 10~h? (b) ¿Cuál es el precio efectivo del calor almacenado de esa forma y conviene precalentar? (c) Repita con $\tau=5$ h (edificio ligero) y con un COP a las 6~h de $2{,}4$ frente a $2{,}9$ a las 10~h.
\end{ejercicio}

\begin{ejercicio}{2}{eje:3-blanda}
En el \cref{ej:mpc-qp}, sustituya la restricción dura $u_0\ge-0{,}5$ por una blanda con holgura $s\ge0$ y penalización $\rho s$. Resuelva el problema para $\rho=0{,}2$ y para $\rho=1$ y relacione el resultado con la \cref{prop:exacta}.
\end{ejercicio}

\begin{ejercicio}{3}{eje:3-riccati}
Demuestre que, en el \cref{teo:lqr}, la matriz $P_k$ es simétrica y semidefinida positiva para todo $k$, y que $P_k$ no disminuye cuando aumenta el horizonte restante si $P_f=0$: $x\T P_kx\le x\T P_{k-1}x$. Interprete el resultado.
\fuente{adaptado de J.~B. Rawlings, D.~Q. Mayne y M.~M. Diehl, \emph{Model Predictive Control}, 2.ª ed., cap.~1 «Getting Started with Model Predictive Control», regulador lineal cuadrático~\cite{rmd}.}
\end{ejercicio}

\begin{ejercicio}{3}{eje:3-terminal}
Para $x_{k+1}=x_k+u_k$, $\ell=x^{2}+u^{2}$ y $\abs{u}\le0{,}5$: (a) compruebe que $V_f(x)=1{,}618\,x^{2}$ y $\kappa_f(x)=-0{,}618\,x$ satisfacen la condición de decrecimiento del \cref{teo:mpc}; (b) calcule el mayor intervalo $\Xset_f=[-\bar x,\bar x]$ en el que $\kappa_f$ respeta la restricción y es invariante; (c) explique por qué, con estos ingredientes, el MPC es estable para cualquier $N\ge1$ desde cualquier estado inicial factible.
\fuente{adaptado de J.~B. Rawlings, D.~Q. Mayne y M.~M. Diehl, \emph{Model Predictive Control}, 2.ª ed., cap.~2 «Model Predictive Control: Regulation», ingredientes terminales~\cite{rmd}.}
\end{ejercicio}

% =====================================================================
\chapter{Aprendizaje por refuerzo para control continuo}
\label{tema:4}
% =====================================================================

El módulo~21 formuló el aprendizaje por refuerzo para modelos de lenguaje: acciones discretas (tokens), recompensa al final y PPO, un método \emph{en política} que descarta las muestras tras cada actualización. En control de procesos las acciones son continuas (potencias, aperturas de válvula), las recompensas llegan en cada paso y cada muestra cuesta: viene de un simulador lento o de la planta. Este tema presenta los métodos actor-crítico \emph{fuera de política} que reutilizan la experiencia (DDPG, TD3 y SAC), cómo se empaqueta un proceso como entorno, cómo se diseña la recompensa sin cambiar el problema y cómo se tratan la seguridad y el salto del simulador a la planta. Sigue a Sutton y Barto~\cite{sb}, las clases~7 («Value Function Methods»), 8 («Deep RL with Q-Functions») y~10 de CS285~\cite{cs285}, la tarea~3 de ese curso~\cite{cs285hw3} y los artículos originales citados en cada sección.

\section{Del PPO a los métodos fuera de política}

\begin{definicion}[MDP de control continuo]\label{def:mdp-cont}
Un MDP de control continuo tiene estados $s\in\Sset\subseteq\R^{n}$, acciones $a\in\Aset=[a_{\min},a_{\max}]^{m}$, dinámica $s'\sim p(\cdot\mid s,a)$, recompensa $r(s,a,s')$ y descuento $\gamma<1$. La función de valor de acción de una política $\pi$ cumple la ecuación de Bellman (módulo~21, tema~1)
\[
Q^{\pi}(s,a)=\E_{s'}\bigl[r(s,a,s')+\gamma\,\E_{a'\sim\pi(\cdot\mid s')}Q^{\pi}(s',a')\bigr].
\]
\end{definicion}

PPO sirve para acciones continuas con una política gaussiana $\pi_\psi(a\mid s)=\mathcal N(\mu_\psi(s),\diag\sigma^{2})$ y es robusto, pero cada lote de experiencia se usa para unas pocas épocas y se descarta. Los métodos fuera de política guardan las transiciones $(s,a,r,s',d)$ en un \emph{búfer de repetición} y aprenden de ellas muchas veces, porque la ecuación de Bellman para $Q^{\pi}$ se puede imponer sobre transiciones generadas por cualquier política. El precio es la inestabilidad del aprendizaje de $Q$ con aproximadores no lineales, que se mitiga con \emph{redes objetivo}: copias lentas $\phi'\leftarrow\tau\phi+(1-\tau)\phi'$, con $\tau\approx0{,}005$, que fijan los objetivos de regresión~\cite{cs285}.

\begin{definicion}[DDPG]\label{def:ddpg}
DDPG~\cite{ddpg} aprende un crítico $Q_\phi(s,a)$ y un actor determinista $a=\mu_\psi(s)$. El crítico minimiza $\E\bigl[(Q_\phi(s,a)-y)^{2}\bigr]$ sobre el búfer con el objetivo $y=r+\gamma(1-d)\,Q_{\phi'}\bigl(s',\mu_{\psi'}(s')\bigr)$, y el actor asciende por el gradiente $\nabla_\psi\E_s\bigl[Q_\phi(s,\mu_\psi(s))\bigr]=\E_s\bigl[\nabla_aQ_\phi(s,a)\rvert_{a=\mu_\psi(s)}\nabla_\psi\mu_\psi(s)\bigr]$, por la regla de la cadena (módulo~4, tema~2). La exploración se hace sumando ruido a la acción.
\end{definicion}

\section{El sesgo de sobreestimación y TD3}

\begin{proposicion}[Sobreestimación por maximización]\label{prop:sobre}
Sean $\hat Q(a)$, $a\in\Aset$ finito, estimaciones aleatorias con $\E[\hat Q(a)]=Q(a)$. Entonces $\E\bigl[\max_a\hat Q(a)\bigr]\ge\max_aQ(a)$. Si son independientes, gaussianas de varianza $\sigma^{2}$ y todos los $Q(a)$ son iguales a $q$, con dos acciones $\E[\max_a\hat Q(a)]=q+\sigma/\sqrt\pi$.
\end{proposicion}

\begin{proof}
Para cada $a_0$, $\max_a\hat Q(a)\ge\hat Q(a_0)$; tomando esperanzas, $\E[\max_a\hat Q(a)]\ge Q(a_0)$ para todo $a_0$, luego también para el que maximiza $Q$. Para el caso gaussiano, $\max(X_1,X_2)=\frac{X_1+X_2}{2}+\frac{\abs{X_1-X_2}}{2}$ y $X_1-X_2\sim\mathcal N(0,2\sigma^{2})$, cuyo valor absoluto tiene media $\sqrt2\sigma\sqrt{2/\pi}=2\sigma/\sqrt\pi$.
\end{proof}

El objetivo de Bellman contiene un máximo, explícito en Q-learning o implícito en el actor que busca la acción de mayor $Q_\phi$; el error de estimación se convierte en sesgo positivo y se propaga hacia atrás por el \emph{bootstrapping}. Con diez acciones de igual valor y $\sigma=1$, el sesgo de cada objetivo es de $1{,}54$. Sutton y Barto lo llaman \emph{sesgo de maximización} y lo corrigen con el doble aprendizaje~\cite[§6.7]{sb}.

\begin{definicion}[TD3]\label{def:td3}
TD3~\cite{td3} modifica DDPG con tres mecanismos: (i) \emph{doble Q recortado}: dos críticos $Q_{\phi_1},Q_{\phi_2}$ y el objetivo $y=r+\gamma(1-d)\min_{i=1,2}Q_{\phi_i'}(s',\tilde a')$, que sustituye el sesgo positivo por uno ligeramente negativo; (ii) \emph{suavizado de la política objetivo}: $\tilde a'=\clip\bigl(\mu_{\psi'}(s')+\clip(\epsilon,-c,c),a_{\min},a_{\max}\bigr)$ con $\epsilon\sim\mathcal N(0,\tilde\sigma^{2})$, que impide explotar picos estrechos y erróneos del crítico; (iii) \emph{actualizaciones retardadas}: el actor y las redes objetivo se actualizan una vez por cada dos actualizaciones de los críticos.
\end{definicion}

\section{Máxima entropía: SAC}

\begin{definicion}[Objetivo de máxima entropía]\label{def:maxent}
SAC~\cite{sac} maximiza $J(\pi)=\sum_t\E\bigl[r_t+\alpha\,\mathcal H(\pi(\cdot\mid s_t))\bigr]$, con $\mathcal H$ la entropía y $\alpha>0$ la \emph{temperatura}. Las funciones de valor «blandas» cumplen
\[
Q(s,a)=\E_{s'}\bigl[r+\gamma V(s')\bigr],\qquad V(s)=\E_{a\sim\pi}\bigl[Q(s,a)-\alpha\log\pi(a\mid s)\bigr].
\]
\end{definicion}

\begin{proposicion}[Política blanda óptima]\label{prop:blanda}
Para un estado con acciones finitas y valores $Q(a)$, la política que maximiza $\E_\pi[Q]+\alpha\mathcal H(\pi)$ es $\pi^{*}(a)=\exp(Q(a)/\alpha)/Z$, con $Z=\sum_a\exp(Q(a)/\alpha)$, y el valor máximo es $V=\alpha\log Z$.
\end{proposicion}

\begin{proof}
Para toda distribución $\pi$, $\E_\pi[Q]+\alpha\mathcal H(\pi)=-\alpha\sum_a\pi(a)\log\frac{\pi(a)}{e^{Q(a)/\alpha}}=\alpha\log Z-\alpha\KL(\pi\,\|\,\pi^{*})$, y la divergencia es no negativa y nula solo si $\pi=\pi^{*}$ (desigualdad de Gibbs, módulo~6). Es el mismo argumento que da la política óptima con penalización KL en el módulo~21 (tema~3), con la distribución uniforme como referencia.
\end{proof}

\begin{ejemplo}[Política blanda]\label{ej:blanda}
Con $Q=(1,2,0)$ y $\alpha=0{,}5$: $e^{2}=7{,}389$, $e^{4}=54{,}598$, $e^{0}=1$, $Z=62{,}987$, $\pi^{*}=(0{,}117;\ 0{,}867;\ 0{,}016)$ y $V=0{,}5\ln62{,}987=2{,}071$, algo mayor que $\max Q=2$: la entropía tiene valor. Con $\alpha\to0$ la política tiende a la voraz.
\end{ejemplo}

En acciones continuas, el actor de SAC es una gaussiana comprimida, $a=\tanh(\mu_\psi(s)+\sigma_\psi(s)\odot\xi)$ con $\xi\sim\mathcal N(0,I)$, escalada al rango de la acción. Esta \emph{reparametrización} permite derivar $\E[Q_\phi(s,a)-\alpha\log\pi_\psi(a\mid s)]$ respecto a $\psi$ a través de la muestra, y el cambio de variable corrige la densidad: $\log\pi(a\mid s)=\log\mathcal N(u;\mu,\sigma^{2})-\sum_i\log(1-\tanh^{2}u_i)$. SAC usa además dos críticos con el mínimo, como TD3. La versión posterior~\cite{sac2} ajusta $\alpha$ automáticamente para mantener una entropía objetivo $\bar{\mathcal H}$ (típicamente $-\dim\Aset$), descendiendo por $\E[-\alpha\log\pi(a\mid s)-\alpha\bar{\mathcal H}]$: si la entropía de la política cae por debajo del objetivo, $\alpha$ sube (\cref{eje:4-temperatura}). La tarea~3 de CS285~\cite{cs285hw3} implementa estos componentes paso a paso.

\section{Entornos: la interfaz de Gymnasium}

\begin{definicion}[Entorno Gymnasium]\label{def:gym}
Un entorno de Gymnasium~\cite{gymnasium} es un objeto con un espacio de observaciones, un espacio de acciones (por ejemplo, \texttt{Box} para vectores continuos acotados) y dos métodos: \texttt{reset(seed)}, que devuelve \texttt{(obs, info)}, y \texttt{step(accion)}, que devuelve \texttt{(obs, recompensa, terminated, truncated, info)}. El indicador \texttt{terminated} señala un estado terminal del MDP (el valor futuro es nulo) y \texttt{truncated} un corte externo, como un límite de tiempo (el valor futuro no es nulo, simplemente no se observa).
\end{definicion}

La distinción importa al calcular objetivos: el término $(1-d)$ de DDPG, TD3 y SAC debe usar $d=\texttt{terminated}$, nunca \texttt{truncated}. Un episodio de climatización que se corta a los tres días no termina el mundo: si se tratara como terminal, el último objetivo sería $y=r$ en lugar de $r+\gamma V(s')$, y el agente aprendería que justo antes del corte da igual pasar frío (\cref{eje:4-truncado}).

Con la interfaz estándar, entrenar un agente es una línea con Stable-Baselines3~\cite{sb3}, que implementa PPO, SAC, TD3 y otros algoritmos:
\begin{verbatim}
from stable_baselines3 import SAC
modelo = SAC("MlpPolicy", EntornoHVAC(), learning_rate=3e-4,
             gamma=0.99, seed=0)
modelo.learn(total_timesteps=50_000)
\end{verbatim}
Las recomendaciones prácticas de esa biblioteca coinciden con la experiencia en procesos: normalizar observaciones y acciones (preferiblemente a $[-1,1]$), usar varias semillas y evaluar con un conjunto de episodios distinto del de entrenamiento.

\section{Diseño de la recompensa}

En control de procesos la recompensa suele ser el negativo de un coste económico más penalizaciones: $r_k=-\bigl(c_kP_k\Dt+\rho\,v_{k+1}\Dt\bigr)$, con $v$ la violación del confort, como en el MPC del \cref{ej:caso-mpc}. Usar el mismo coste en MPC y en RL es lo que hace comparables los dos enfoques. Dos transformaciones de la recompensa merecen atención: sumar una constante y sumar términos de \emph{conformación} (\emph{shaping}) para guiar el aprendizaje.

\begin{teorema}[Conformación basada en potenciales]\label{teo:shaping}
Sea $\Phi:\Sset\to\R$ acotada y $F(s,a,s')=\gamma\Phi(s')-\Phi(s)$. En un MDP con $\gamma<1$, las funciones de valor de cualquier política $\pi$ con la recompensa $r'=r+F$ cumplen $Q'^{\pi}(s,a)=Q^{\pi}(s,a)-\Phi(s)$. En consecuencia, las políticas óptimas con $r$ y con $r'$ coinciden.
\end{teorema}

\begin{proof}
Para una trayectoria desde $s_0=s$, $a_0=a$, el retorno modificado es
\[
\sum_{t\ge0}\gamma^{t}\bigl(r_t+\gamma\Phi(s_{t+1})-\Phi(s_t)\bigr)=\sum_{t\ge0}\gamma^{t}r_t+\sum_{t\ge0}\bigl(\gamma^{t+1}\Phi(s_{t+1})-\gamma^{t}\Phi(s_t)\bigr)=\sum_{t\ge0}\gamma^{t}r_t-\Phi(s_0),
\]
porque la segunda suma es telescópica y $\gamma^{T}\Phi(s_T)\to0$ al estar $\Phi$ acotada. Tomando esperanzas, $Q'^{\pi}(s,a)=Q^{\pi}(s,a)-\Phi(s)$. Como la diferencia no depende de $a$, $\arg\max_aQ'^{*}(s,a)=\arg\max_aQ^{*}(s,a)$ en cada estado.
\end{proof}

El resultado es de Ng, Harada y Russell~\cite{ng}, que demostraron además que los potenciales son, en un sentido preciso, la única forma de conformación que garantiza no alterar la política óptima. Una bonificación arbitraria, como premiar estar cerca de la consigna, puede crear ciclos que el agente explota indefinidamente. Sumar una constante a todas las recompensas es inocuo en tareas continuadas, pero no en episódicas con terminación, donde altera el valor relativo de terminar pronto o tarde~\cite[ejercicios~3.15 y~3.16]{sb}.

\section{RL seguro, restricciones y del simulador a la planta}

\begin{definicion}[MDP con restricciones]\label{def:cmdp}
Un MDP con restricciones (CMDP)~\cite{altman} añade costes $c(s,a)$ y límites $d$: $\max_\pi J_r(\pi)$ s.a. $J_c(\pi)\le d$, con $J_r$ y $J_c$ los valores esperados descontados de la recompensa y del coste. Su lagrangiana es $L(\pi,\lambda)=J_r(\pi)-\lambda\bigl(J_c(\pi)-d\bigr)$, con $\lambda\ge0$.
\end{definicion}

\begin{proposicion}[Dualidad débil]\label{prop:dual}
Para todo $\lambda\ge0$, $\max_\pi L(\pi,\lambda)\ge J_r(\pi^{\star})$, con $\pi^{\star}$ óptima del CMDP. Por tanto $\min_{\lambda\ge0}\max_\pi L(\pi,\lambda)$ es una cota superior del valor óptimo.
\end{proposicion}

\begin{proof}
$\pi^{\star}$ es factible, $J_c(\pi^{\star})\le d$, luego $L(\pi^{\star},\lambda)=J_r(\pi^{\star})-\lambda(J_c(\pi^{\star})-d)\ge J_r(\pi^{\star})$, y el máximo sobre $\pi$ es mayor o igual.
\end{proof}

Los métodos lagrangianos alternan una mejora de la política con $\lambda$ fijo, con cualquier algoritmo de RL sobre $r-\lambda c$, y un paso de ascenso dual, $\lambda\leftarrow\max\{0,\lambda+\eta(\hat J_c-d)\}$: si la restricción se viola, el coste se encarece. Garantizan el cumplimiento en media al final del entrenamiento, no durante él. Los métodos de región de confianza como CPO~\cite{cpo} buscan mantener la factibilidad en cada actualización. En planta, ninguna de estas garantías estadísticas sustituye a un \emph{filtro de seguridad} determinista entre el agente y el actuador (\cref{tema:6}).

Un agente entrenado en un gemelo digital sufre el \emph{salto de la simulación a la realidad}. El siguiente resultado acota su efecto.

\begin{teorema}[Lema de simulación]\label{teo:simlema}
Sean dos MDP con los mismos estados (finitos), acciones y recompensas, $\abs{r}\le r_{\max}$, y dinámicas $p$ y $\hat p$ con $\max_{s,a}\norm{p(\cdot\mid s,a)-\hat p(\cdot\mid s,a)}_1\le\epsilon$. Para toda política $\pi$, $\norm{V^{\pi}-\hat V^{\pi}}_\infty\le\gamma\epsilon\,r_{\max}/(1-\gamma)^{2}$.
\end{teorema}

\begin{proof}
Por la ecuación de Bellman de cada MDP, para cada $s$,
\[
V^{\pi}(s)-\hat V^{\pi}(s)=\gamma\sum_a\pi(a\mid s)\sum_{s'}\Bigl[\bigl(p-\hat p\bigr)(s'\mid s,a)\,\hat V^{\pi}(s')+p(s'\mid s,a)\bigl(V^{\pi}-\hat V^{\pi}\bigr)(s')\Bigr],
\]
sumando y restando $\gamma\sum p\,\hat V^{\pi}$ (los términos de recompensa son iguales y se cancelan). El primer sumando interior está acotado por $\epsilon\norm{\hat V^{\pi}}_\infty\le\epsilon r_{\max}/(1-\gamma)$, y el segundo por $\norm{V^{\pi}-\hat V^{\pi}}_\infty$, porque $p$ suma~1. Tomando el máximo en $s$, $\Delta\le\gamma\epsilon r_{\max}/(1-\gamma)+\gamma\Delta$ con $\Delta=\norm{V^{\pi}-\hat V^{\pi}}_\infty$, y se despeja $\Delta$.
\end{proof}

El factor $1/(1-\gamma)^{2}$ dice que los errores de modelo pesan más cuanto más largo es el horizonte efectivo: es la versión estocástica de la \cref{prop:acumula}. Dos remedios complementarios son mejorar el modelo (identificación, \cref{tema:2}) y hacer la política robusta a lo que no se sabe del modelo con la \emph{aleatorización del dominio}: entrenar maximizando $\E_{\xi\sim p(\xi)}[J(\pi;\xi)]$, donde $\xi$ son los parámetros inciertos del simulador. Tobin \emph{et al.}~\cite{tobin} aleatorizaron la apariencia visual del simulador y Peng \emph{et al.}~\cite{peng}, su dinámica, y en ambos casos los modelos entrenados solo en simulación funcionaron en robots reales: un detector de objetos usado para agarrar, en el primero, y una política de empuje de objetos, en el segundo. Cuando no hay simulador fiable pero sí históricos, el RL \emph{fuera de línea} aprende de un conjunto fijo de datos (CS285, clases~15 y~16 de otoño de 2023~\cite{cs285}).

\begin{observacion}[Búsqueda de política sin gradiente]\label{obs:busqueda}
Para políticas con pocos parámetros, la búsqueda aleatoria o el método de entropía cruzada (CEM), que muestrean parámetros, evalúan episodios completos y se quedan con los mejores, son referencias sorprendentemente competitivas: Mania, Guy y Recht~\cite{ars} obtuvieron con búsqueda aleatoria de políticas lineales resultados comparables a los de métodos profundos en tareas de locomoción. Son fáciles de paralelizar, no necesitan crítico y solo usan el retorno, lo que las hace robustas a recompensas dispersas.
\end{observacion}

\begin{ejemplo}[La batería como MDP]\label{ej:bateria-mdp}
Una batería de $1$ kWh con estados $e\in\{0,1\}$, acciones cargar, descargar o esperar, rendimiento~1 y precios $(0{,}10;\ 0{,}30;\ 0{,}20)$ en tres horas es un MDP determinista de horizonte finito. Hacia atrás, con $V_3\equiv0$: $V_2(1)=\max\{0;\ 0{,}20\}=0{,}20$ y $V_2(0)=\max\{0;\ -0{,}20+V_3(1)\}=0$; $V_1(1)=\max\{V_2(1);\ 0{,}30+V_2(0)\}=0{,}30$ y $V_1(0)=\max\{V_2(0);\ -0{,}30+V_2(1)\}=0$; $V_0(0)=\max\{V_1(0);\ -0{,}10+V_1(1)\}=0{,}20$. La política óptima carga en la primera hora y descarga en la segunda. Con 24 horas y una rejilla de 81 niveles de carga, la programación dinámica del cuaderno (Parte~8) obtiene $1{,}8837$ €, frente a $1{,}8842$ € del LP: la diferencia es el coste de discretizar el estado.
\end{ejemplo}

\begin{ejemplo}[Caso de estudio: un agente para la zona]\label{ej:caso-rl}
\emph{Entorno.} El gemelo de la zona como entorno Gymnasium: observación de 8 componentes (temperatura medida, temperatura exterior, precio, seno y coseno de la hora, límite inferior actual y dentro de 2~h, ocupación en las próximas 3~h), acción $P\in[0,3]$ kW, recompensa $r_k=-\Dt(c_kP_k+\rho v_{k+1})$ con $\rho=2$ €/(K·h), episodios de 3 días truncados por tiempo. En cada \texttt{reset} se aleatorizan el clima, el día de la semana, la temperatura inicial y la planta: $R$ y $C$ con $\pm20$~\% y el COP con $\pm10$~\%.

\emph{Agente.} Política $P=P_{\max}\,\sigma(w\T\phi(o))$, con $\phi$ los 8 rasgos escalados y $w\in\R^{8}$, optimizada con CEM: 40 iteraciones con 60 candidatos, 12 de élite y 4 episodios aleatorizados por candidato.

\emph{Elección del método.} Ocho parámetros, recompensa densa y un simulador vectorizable: CEM entrena en unos 3~s, mientras que SAC con Stable-Baselines3 necesitaría decenas de miles de pasos y varios minutos para un resultado comparable (reto opcional del cuaderno).

\emph{Resultados} (cuaderno, Parte~6). En la semana de prueba, que no se usó en el entrenamiento, el agente consume $169{,}9$ kWh, cuesta $20{,}63$ € y acumula $0{,}01$ K·h de violación: mejora al PI en un $2{,}1$~\% de coste y queda por detrás del MPC. Ha aprendido a precalentar (el peso de la ocupación próxima es el mayor, $3{,}13$) y a reaccionar al precio, pero con una política lineal y sin previsiones de precio a varias horas no puede planificar tan bien como el MPC.
\end{ejemplo}

\begin{table}[htbp]
\centering\footnotesize
\begin{tabularx}{\textwidth}{@{}>{\raggedright\arraybackslash}p{3.1cm}>{\raggedright\arraybackslash}X>{\raggedright\arraybackslash}p{2.0cm}>{\raggedright\arraybackslash}X@{}}
\toprule
Herramienta & Para qué sirve & Licencia o tipo & Cuándo elegirla \\
\midrule
Gymnasium~\cite{gymnasium} & interfaz estándar de entornos, envoltorios, comprobador de entornos & libre (MIT) & siempre que se vaya a entrenar un agente \\
Stable-Baselines3~\cite{sb3} & PPO, SAC, TD3, DQN fiables y documentados & libre (MIT) & primera opción para entrenar y comparar \\
CleanRL & implementaciones de un solo archivo & libre (MIT) & entender o modificar un algoritmo \\
RLlib (Ray), TorchRL & entrenamiento distribuido, muchos entornos en paralelo & libres (Apache-2.0, MIT) & simuladores lentos y grandes presupuestos de cómputo \\
Safety-Gymnasium & entornos con costes y restricciones & libre (Apache-2.0) & evaluar métodos de RL seguro \\
Sinergym~\cite{sinergym}, BOPTEST~\cite{boptest} & edificios simulados (EnergyPlus, Modelica) como entornos y bancos de prueba & libres & control de climatización con modelos físicos detallados \\
\bottomrule
\end{tabularx}
\caption{Tecnologías para el aprendizaje por refuerzo en control.}
\label{tab:tec4}
\end{table}

\begin{observacion}[Errores frecuentes en la práctica]\label{obs:errores4}
(i) \emph{Tratar la truncación como terminación} al calcular objetivos (\cref{def:gym}). (ii) \emph{Recompensa explotable}: el agente maximiza lo que se mide, no lo que se quiere; por ejemplo, si el confort solo se penaliza en horario de ocupación, puede dejar enfriar el edificio hasta dañar instalaciones. (iii) \emph{Escalas mal elegidas}: observaciones sin normalizar y penalizaciones de órdenes de magnitud distintos desestabilizan el aprendizaje. (iv) \emph{Evaluar con las semillas de entrenamiento} o con una sola semilla de entrenamiento. (v) \emph{Explorar en la planta real}: la exploración aleatoria de un agente sin restricciones es inaceptable en un proceso; se explora en el gemelo y en planta solo dentro de límites seguros. (vi) \emph{Sobreajuste al simulador}: una política que explota un defecto del gemelo falla en la planta; la aleatorización del dominio y la validación en un gemelo distinto lo detectan.
\end{observacion}

\begin{lecturas}
\item R.~S. Sutton y A.~G. Barto~\cite{sb}, §6.7 «Maximization Bias and Double Learning» y cap.~13 «Policy Gradient Methods».
\item CS285~\cite{cs285}, clases~7, 8 y~10; tarea~3~\cite{cs285hw3}, §2.5 y §3.
\item S.~Fujimoto \emph{et al.}~\cite{td3} (TD3) y T.~Haarnoja \emph{et al.}~\cite{sac,sac2} (SAC).
\item A.~Y. Ng, D.~Harada y S.~Russell~\cite{ng}, «Policy invariance under reward transformations».
\item G.~Dulac-Arnold \emph{et al.}~\cite{dulac}, «Challenges of real-world reinforcement learning».
\end{lecturas}

\section*{Ejercicios del tema 4}
\addcontentsline{toc}{section}{Ejercicios del tema 4}

\begin{ejercicio}{1}{eje:4-truncado}
Un episodio de la zona se trunca tras 288 pasos. Las recompensas típicas son de $-0{,}05$ € por paso y $\gamma=0{,}99$. (a) Estime el valor de un estado típico. (b) Calcule el objetivo del último paso si se trata la truncación como terminación y si se hace correctamente. (c) ¿Qué comportamiento aprendería el agente en el primer caso?
\end{ejercicio}

\begin{ejercicio}{2}{eje:4-constante}
(a) En una tarea continuada con descuento $\gamma<1$, demuestre que sumar una constante $c$ a todas las recompensas suma $c/(1-\gamma)$ al valor de todos los estados y no cambia la política óptima. (b) En una tarea episódica, muestre con un ejemplo de dos acciones (terminar ya o seguir un paso más y terminar) que sumar $c$ puede cambiar la política óptima.
\fuente{adaptado de R.~S. Sutton y A.~G. Barto, \emph{Reinforcement Learning: An Introduction}, 2.ª ed., ejercicios~3.15 y~3.16~\cite{sb}.}
\end{ejercicio}

\begin{ejercicio}{2}{eje:4-sobre}
Dos acciones tienen valores verdaderos $Q(a_1)=0$ y $Q(a_2)=-0{,}5$, y sus estimaciones son independientes y gaussianas con $\sigma=1$. (a) Calcule $\E[\max(\hat Q(a_1),\hat Q(a_2))]$ y compárela con $\max_aQ(a)$. (b) Con dos estimaciones independientes $\hat Q_1,\hat Q_2$ del mismo valor $0$, calcule $\E[\min(\hat Q_1,\hat Q_2)]$. (c) Explique cómo se relacionan (a) y (b) con las curvas de valores $Q$ que se piden en la tarea al comparar un crítico con doble Q recortado.
\fuente{adaptado de CS285, primavera de 2026, tarea~3 «Q-Learning and Actor-Critic Algorithms», §3.6, pregunta~7~\cite{cs285hw3}.}
\end{ejercicio}

\begin{ejercicio}{2}{eje:4-blanda}
Con $Q=(1,2,0)$, calcule la política blanda óptima y el valor $V$ para $\alpha=1$ y $\alpha=0{,}1$. Compruebe en cada caso que $V=\E_{\pi^{*}}[Q]+\alpha\mathcal H(\pi^{*})$.
\end{ejercicio}

\begin{ejercicio}{2}{eje:4-temperatura}
Un actor de SAC con acción unidimensional tiene una entropía media estimada $\mathcal H=-1{,}5$ nats; la entropía objetivo es $\bar{\mathcal H}=-1$. (a) Calcule la derivada de $J(\alpha)=\E[-\alpha\log\pi-\alpha\bar{\mathcal H}]$ respecto a $\alpha$ y el nuevo valor de $\alpha$ tras un paso de descenso con tasa $0{,}1$ desde $\alpha=0{,}2$. (b) Calcule la entropía de una gaussiana con $\sigma=0{,}5$ y la corrección de densidad de la compresión con $\tanh$ en $u=0{,}5$.
\fuente{adaptado de CS285, primavera de 2026, tarea~3 «Q-Learning and Actor-Critic Algorithms», §3.4 y §3.5~\cite{cs285hw3}.}
\end{ejercicio}

\begin{ejercicio}{2}{eje:4-lagrange}
Tres políticas candidatas para la zona tienen, por semana, $(J_r,J_c)=(-18;\ 2{,}0)$, $(-19;\ 0{,}5)$ y $(-21;\ 0)$, con $J_r$ en euros cambiados de signo y $J_c$ en K·h; el límite es $d=1$ K·h. (a) ¿Cuál es la óptima del CMDP? (b) Para qué valores de $\lambda$ maximiza cada una la lagrangiana. (c) Simule cuatro pasos de ascenso dual con $\eta=0{,}5$ desde $\lambda=0$ y comente el comportamiento.
\end{ejercicio}

\begin{ejercicio}{3}{eje:4-shaping}
(a) Extienda el \cref{teo:shaping} a una tarea episódica con estados terminales, suponiendo $\Phi(s)=0$ en ellos, y demuestre que la política óptima no cambia. (b) Para la zona se propone $\Phi(s)=-\kappa\max\{\underline T-T,0\}$, con $\kappa>0$. Calcule $F$ en una transición de $T=19$ °C a $T'=19{,}5$ °C con $\underline T=20$ °C, $\kappa=1$ y $\gamma=0{,}99$, e interprete el signo. (c) Dé un ejemplo de bonificación que \emph{no} derive de un potencial y que el agente pueda explotar.
\fuente{adaptado de A.~Y. Ng, D.~Harada y S.~Russell, «Policy invariance under reward transformations», ICML 1999, teorema~1~\cite{ng}.}
\end{ejercicio}

\begin{ejercicio}{3}{eje:4-simlema}
(a) Evalúe la cota del \cref{teo:simlema} con $r_{\max}=1$ y $\epsilon=0{,}01$ para $\gamma=0{,}9$ y $\gamma=0{,}99$, y compárela con el rango posible de $V$. (b) Deduzca del lema una cota para la pérdida de rendimiento de la política óptima del simulador cuando se aplica en la planta, $V^{\pi^{*}}-V^{\hat\pi^{*}}$. (c) Explique qué dice el resultado sobre la conveniencia de horizontes efectivos largos cuando el modelo es imperfecto.
\end{ejercicio}

% =====================================================================
\chapter{Casos de estudio: edificio, reactor, batería y centro de datos}
\label{tema:5}
% =====================================================================

Este tema reúne los métodos de los anteriores en cuatro problemas de distinta naturaleza: un edificio, donde el objetivo es económico y las restricciones son de confort; un reactor por lotes, donde lo que manda es la seguridad; una batería, donde el problema es casi puramente de optimización, y un centro de datos, el caso industrial publicado más conocido de aprendizaje automático aplicado a la operación de una instalación. Antes se fija cómo comparar controladores de forma honesta. Las fuentes son la revisión de Drgoňa \emph{et al.}~\cite{drgona} y el marco BOPTEST~\cite{boptest} para edificios, la lección~6 de MIT 10.450~\cite{mit10450} y el texto de Seborg \emph{et al.}~\cite{seborg} para el reactor, y las publicaciones de Google y DeepMind~\cite{evansgao,gamblegao} y de Lazic \emph{et al.}~\cite{lazic} para los centros de datos. Los resultados numéricos de los tres primeros casos proceden del cuaderno del módulo.

\section{Cómo comparar controladores}

\begin{definicion}[Indicadores de rendimiento]\label{def:kpi}
Para un periodo de evaluación de $n$ pasos se definen la energía $E=\Dt\sum_kP_k$ (kWh), el coste $\Sigma=\Dt\sum_kc_kP_k$ (€), la violación del confort $D=\Dt\sum_k\bigl(\max\{\underline T_k-T_k,0\}+\max\{T_k-\overline T_k,0\}\bigr)$ (K·h) y, si se quiere un único número, el coste total $\Sigma+\rho D$ con el mismo $\rho$ que usan los controladores.
\end{definicion}

Un protocolo de comparación justo exige: (i) el mismo periodo de prueba, no usado para ajustar ningún controlador; (ii) la misma información, en particular las mismas previsiones: un MPC con la meteorología observada frente a un PID que no la conoce compara información, no métodos; (iii) un controlador de referencia bien ajustado, no un PID descuidado; (iv) los mismos ruidos, y (v) varios periodos o semillas, con una medida de la dispersión. BOPTEST~\cite{boptest} se creó precisamente para normalizar estas comparaciones en edificios: fija los modelos, los escenarios y los indicadores.

\begin{proposicion}[Números aleatorios comunes]\label{prop:crn}
Sean $X$ e $Y$ los costes de dos controladores en un mismo escenario aleatorio, con varianzas $\sigma_X^{2}$, $\sigma_Y^{2}$ y correlación $r$. Entonces $\Var(X-Y)=\sigma_X^{2}+\sigma_Y^{2}-2r\sigma_X\sigma_Y$. Si $\sigma_X=\sigma_Y=\sigma$, evaluar los dos controladores con los mismos escenarios ($r>0$) divide la varianza de la diferencia por $1-r$ respecto a usar escenarios independientes ($r=0$).
\end{proposicion}

\begin{proof}
$\Var(X-Y)=\Var X+\Var Y-2\operatorname{Cov}(X,Y)$ y $\operatorname{Cov}(X,Y)=r\sigma_X\sigma_Y$ (módulo~5, tema~2). Con $\sigma_X=\sigma_Y$, la varianza es $2\sigma^{2}(1-r)$ frente a $2\sigma^{2}$.
\end{proof}

El coste semanal de climatización depende sobre todo del tiempo meteorológico: dos controladores evaluados en la misma semana tienen costes muy correlacionados, y su diferencia se estima con mucha menos incertidumbre que cada coste por separado. El mismo principio guía el entrenamiento del agente del \cref{ej:caso-rl}, que evalúa todos los candidatos de una iteración en los mismos episodios.

\section{Climatización de un edificio}

El caso de la zona ha recorrido los temas~1 a~4. El \cref{tab:hvac} reúne los resultados en la semana de prueba (temperatura exterior media de $3{,}35$ °C) con la planta nominal y con una planta perturbada ($R$ un 20~\% mayor, $C$ un 20~\% menor y COP un 10~\% peor) que ninguno de los controladores conoce.

\begin{table}[htbp]
\centering\small
\begin{tabular}{@{}lcccccc@{}}
\toprule
& \multicolumn{4}{c}{Planta nominal} & \multicolumn{2}{c}{Planta perturbada} \\
\cmidrule(lr){2-5}\cmidrule(l){6-7}
Controlador & $E$ (kWh) & $\Sigma$ (€) & $D$ (K·h) & Ahorro & $\Sigma$ (€) & $D$ (K·h) \\
\midrule
PI & $175{,}4$ & $21{,}06$ & $0{,}00$ & --- & $18{,}75$ & $0{,}00$ \\
MPC, margen $0{,}3$ K & $171{,}1$ & $19{,}28$ & $0{,}00$ & $8{,}4$~\% & $17{,}28$ & $0{,}00$ \\
MPC, sin margen & $166{,}9$ & $18{,}76$ & $3{,}54$ & --- & --- & --- \\
RL (CEM) & $169{,}9$ & $20{,}63$ & $0{,}01$ & $2{,}1$~\% & $18{,}50$ & $0{,}00$ \\
\bottomrule
\end{tabular}
\caption{Climatización de la zona: comparación en la semana de prueba (cuaderno, Partes~2, 4, 6 y~7): energía $E$, coste $\Sigma$ y violación del confort $D$ (\cref{def:kpi}). PI con ajuste SIMC y consigna programada; MPC económico; agente CEM con aleatorización del dominio. El ahorro es en coste respecto al PI.}
\label{tab:hvac}
\end{table}

Los resultados ilustran rasgos generales que la literatura sobre edificios documenta con casos reales~\cite{drgona,wanghong}:
\begin{enumerate}
  \item \emph{El MPC gana porque planifica}: conoce el precio y las ganancias de las próximas 24~h y usa la masa térmica para desplazar el consumo de punta a valle. Su ventaja se mantiene con la planta perturbada (un $7{,}8$~\% menos de coste que el PI), porque realimenta el estado medido en cada paso.
  \item \emph{El margen es parte del diseño}: sin él, el MPC «ahorra» un $2{,}7$~\% más a costa de violar el confort.
  \item \emph{El agente aprende sin modelo explícito, pero con un simulador}: con una política lineal de 8 parámetros mejora al PI y es robusto a la perturbación de la planta gracias a la aleatorización, pero queda lejos del MPC, que dispone de previsiones y de un horizonte de planificación explícito. Una política más rica (una red con previsiones como entrada) y un algoritmo como SAC pueden acortar la distancia, con más datos de simulación y más riesgo de explotar defectos del gemelo.
  \item \emph{El PI sigue siendo la referencia}: no necesita modelo, ni previsiones, ni entrenamiento, y en este caso solo pierde un $8$~\% de coste. En la práctica, el MPC o el agente se colocan \emph{encima} de los PI existentes, enviándoles consignas (\cref{tema:6}).
\end{enumerate}

\begin{table}[htbp]
\centering\small
\begin{tabularx}{\textwidth}{@{}l>{\raggedright\arraybackslash}X>{\raggedright\arraybackslash}X>{\raggedright\arraybackslash}X@{}}
\toprule
& PID & MPC & RL \\
\midrule
Necesita & ajuste local (FOPDT) & modelo identificado y previsiones & simulador fiable (gemelo) y entrenamiento \\
Restricciones & solo límites del actuador & explícitas, en cada paso & penalizadas en la recompensa o con un filtro externo \\
Coste en línea & despreciable & un problema de optimización por paso & una evaluación de la política \\
Adaptación a cambios & reajuste manual & cambiar el modelo o el coste & reentrenar \\
Interpretabilidad & alta & alta (el plan es visible) & baja \\
Cuándo elegirlo & lazos básicos, sin objetivo económico & objetivos económicos con buen modelo y previsiones & dinámicas difíciles de modelar con un simulador rico o decisiones discretas complejas \\
\bottomrule
\end{tabularx}
\caption{Comparación cualitativa de los tres enfoques.}
\label{tab:enfoques}
\end{table}

\section{Control de temperatura de un reactor por lotes}

En un reactor por lotes con una reacción exotérmica, la temperatura sigue una receta (calentar, mantener, enfriar) y el riesgo dominante es el \emph{embalamiento térmico}: si el calor generado por la reacción crece con la temperatura más deprisa de lo que la camisa puede retirarlo, la temperatura se dispara. El análisis de estabilidad de la lección~6 de MIT 10.450~\cite{mit10450}, que compara las curvas de calor generado y calor retirado, se resume en el criterio siguiente.

\begin{definicion}[Modelo del reactor]\label{def:reactor}
Para una reacción de primer orden $\mathrm A\to\mathrm B$ en un reactor de volumen $V$ refrigerado por una camisa a temperatura $T_j$:
\[
\begin{gathered}
\dot C_A=-k(T)C_A,\qquad \rho c_pV\,\dot T=Q_g(T,C_A)-Q_r(T),\\
Q_g=(-\Delta H)\,k(T)\,C_AV,\qquad Q_r=UA\,(T-T_j),
\end{gathered}
\]
con $k(T)=k_{\mathrm{ref}}\exp\bigl(-\frac{E}{R}(\frac1T-\frac1{T_{\mathrm{ref}}})\bigr)$ (ley de Arrhenius).
\end{definicion}

\begin{proposicion}[Criterio de Semenov]\label{prop:semenov}
Sea $C_A$ aproximadamente constante (fase inicial del lote, con conversión pequeña) y sea $\bar T$ un equilibrio de la ecuación de energía, $Q_g(\bar T)=UA(\bar T-T_j)$. El equilibrio es asintóticamente estable si $Q_g'(\bar T)<UA$ e inestable si $Q_g'(\bar T)>UA$. Además, $Q_g'(T)=Q_g(T)\,E/(RT^{2})$.
\end{proposicion}

\begin{proof}
Con $C_A$ fijo, la temperatura cumple la ecuación escalar $\dot T=F(T)=(Q_g(T)-UA(T-T_j))/(\rho c_pV)$. Por la \cref{prop:lineal}, el linealizado en $\bar T$ es $\dot{\delta T}=\lambda\,\delta T$ con $\lambda=F'(\bar T)=(Q_g'(\bar T)-UA)/(\rho c_pV)$; si $\lambda<0$ las desviaciones decaen y si $\lambda>0$ crecen (teorema de linealización de Lyapunov~\cite[cap.~5]{am}). La derivada de la exponencial da $k'(T)=k(T)\,E/(RT^{2})$, de donde se sigue la última fórmula.
\end{proof}

\begin{ejemplo}[Caso de estudio: margen térmico de un reactor]\label{ej:caso-reactor}
\emph{Datos.} $V=1$ m³, $\rho c_pV=4000$ kJ/K, $-\Delta H=100$ kJ/mol, $C_{A,0}=2000$ mol/m³, $E/R=10\,000$ K, $k_{\mathrm{ref}}=10^{-4}$ s$^{-1}$ a $T_{\mathrm{ref}}=350$ K y $UA=2$ kW/K.

\emph{Punto de operación.} A $350$ K y conversión nula, $Q_g=100\cdot10^{-4}\cdot2000=20$ kW; el equilibrio exige $T_j=350-20/2=340$ K. La pendiente es $Q_g'=20\cdot10^{4}/350^{2}=1{,}633$ kW/K, menor que $UA=2$ kW/K: el punto es estable, con constante de tiempo $4000/(2-1{,}633)=10\,889$ s, unas 3~h.

\emph{Margen.} Con $T_j=340$ K hay un segundo equilibrio en $T_u=355{,}23$ K, donde $Q_g'=2{,}414>UA$: es inestable y separa las condiciones iniciales que vuelven a $350$ K de las que se embalan. En el cuaderno (Parte~9), partiendo de $354$ K la temperatura vuelve hacia $350$ K, y partiendo de $356$ K supera los $400$ K en unas $4{,}6$ h. Además, existe una temperatura de camisa crítica: las curvas $Q_g$ y $Q_r$ son tangentes cuando $Q_g(T)=UA\,T^{2}R/E$, es decir, en $T^{*}=352{,}69$ K con $T_j^{*}=T^{*}-RT^{*2}/E=340{,}26$ K; por encima de ella no existe ningún equilibrio y la fase inicial se embala desde cualquier temperatura. El margen de operación es de solo $0{,}26$ K de camisa y $5{,}2$ K de reactor.

\emph{Elección del método de control.} El esquema clásico es una \emph{cascada}: un PID maestro sobre la temperatura del reactor fija la consigna de un PID esclavo sobre la camisa, que rechaza rápido las perturbaciones del agua de refrigeración~\cite{mit10450,seborg}. Un NMPC con el modelo de la \cref{def:reactor} puede seguir la receta minimizando el tiempo de lote con restricciones explícitas $T\le T_{\max}$ y $T_j\le T_j^{*}-\delta$. Para un agente de RL, cualquier exploración en planta es inaceptable; se entrena en el gemelo y se despliega detrás de un \emph{filtro de seguridad} que limita $T_j$ y fuerza refrigeración máxima si $T$ se acerca a $T_u$ (cuaderno, Parte~9; \cref{tema:6}).
\end{ejemplo}

\section{Gestión de una batería con precios horarios}

El \cref{ej:bateria-lp} y el \cref{ej:bateria-mdp} resuelven el mismo problema por dos vías. Con precios conocidos, el LP es exacto y rápido: $1{,}8842$ € en el día del cuaderno. La programación dinámica sobre una rejilla de 81 niveles obtiene $1{,}8837$ €, y su interés aparece cuando los precios son \emph{inciertos}: la PD estocástica maximiza el beneficio esperado con el precio como parte del estado, y el RL la aproxima cuando el estado crece (varias baterías, una planta fotovoltaica, la degradación como estado). Tres decisiones de modelado dominan el resultado: el rendimiento de ida y vuelta (un ciclo solo compensa si el cociente de precios supera $1/\eta^{2}$), el coste de degradación por kWh cíclico y la condición terminal (sin $e_{24}=e_0$, el optimizador vacía la batería al final del día porque el modelo no ve el día siguiente).

\section{Refrigeración de centros de datos: el caso publicado de Google y DeepMind}

Lo que sigue usa exclusivamente la información publicada; los detalles internos del sistema no se han hecho públicos.

\emph{2016: recomendaciones.} Según la publicación de DeepMind de julio de 2016~\cite{evansgao}, un conjunto de redes neuronales profundas se entrenó con datos históricos de miles de sensores (temperaturas, potencias, velocidades de bombas, consignas) para predecir la eficacia en el uso de la energía (PUE, cociente entre la energía total del edificio y la consumida por los equipos informáticos) futura; otros dos conjuntos de redes predecían la temperatura y la presión de la hora siguiente para comprobar que las acciones recomendadas respetaban los límites de operación. En pruebas en un centro de datos en funcionamiento, alternando periodos con el sistema activado y desactivado, se publicó una reducción de hasta el 40~\% de la energía de refrigeración, equivalente a un 15~\% de reducción del exceso de PUE sobre la unidad.

\emph{2018: control autónomo.} En agosto de 2018~\cite{gamblegao} se publicó que el sistema pasó a controlar directamente la refrigeración: cada cinco minutos, un sistema de IA en la nube toma una instantánea de miles de sensores, predice con redes neuronales el efecto de combinaciones de acciones sobre el consumo futuro y elige las que lo minimizan respetando restricciones de seguridad; el sistema de control local verifica las acciones antes de aplicarlas. Entre las salvaguardas publicadas figuran la estimación de la incertidumbre (solo se aplican acciones de alta confianza), la verificación en dos capas, la vuelta automática a un estado neutro si se violan restricciones, reglas heurísticas de respaldo y la posibilidad de que los operadores salgan del modo de IA en cualquier momento. Se publicó un ahorro medio de alrededor del 30~\% tras unos meses de operación.

\emph{Un estudio académico relacionado.} Lazic \emph{et al.}~\cite{lazic} describieron en NeurIPS 2018 un MPC con un modelo ARX disperso de historia de 5 pasos, ajustado por mínimos cuadrados recursivos tras unas 3~h de exploración aleatoria acotada, que regulaba temperaturas y presiones de un sector de servidores cada 30~s; la exploración limitaba el rango y la variación máxima de cada actuador entre pasos. En condiciones comparables, el coste del controlador fue de $116{,}7$ frente a $142{,}2$ de los controladores locales, y estimaron un ahorro de alrededor del 9~\% en el coste de refrigeración de ese sector, con la advertencia de que harían falta experimentos más largos para confirmarlo.

El caso reúne casi todo el módulo: un modelo aprendido (\cref{tema:2}), un optimizador sobre predicciones con restricciones (\cref{tema:3}), un periodo de recomendación con humanos en el lazo antes del control autónomo, salvaguardas en capas y la medición del ahorro con experimentos de activación alterna (\cref{tema:6}).

\begin{table}[htbp]
\centering\footnotesize
\begin{tabularx}{\textwidth}{@{}>{\raggedright\arraybackslash}p{3.1cm}>{\raggedright\arraybackslash}X>{\raggedright\arraybackslash}p{2.0cm}>{\raggedright\arraybackslash}X@{}}
\toprule
Herramienta & Para qué sirve & Licencia o tipo & Cuándo elegirla \\
\midrule
EnergyPlus~\cite{crawley} & simulación energética de edificios completa & libre & gemelo físico de un edificio \\
Biblioteca Buildings de Modelica~\cite{modelica} & componentes de climatización y control & libre & modelos de instalaciones a medida \\
BOPTEST~\cite{boptest} & bancos de prueba normalizados para controladores de edificios & libre & comparar controladores de forma reproducible \\
Sinergym~\cite{sinergym} & EnergyPlus como entorno Gymnasium & libre (MIT) & entrenar agentes de climatización \\
CasADi, do-mpc~\cite{casadi,dompc} & NMPC del reactor con restricciones de seguridad & libres (LGPL-3) & procesos químicos no lineales \\
CVXPY, Pyomo & programación lineal de baterías y despacho & libres & arbitraje y planificación energética \\
\bottomrule
\end{tabularx}
\caption{Tecnologías para los casos de estudio.}
\label{tab:tec5}
\end{table}

\begin{observacion}[Errores frecuentes en la práctica]\label{obs:errores5}
(i) \emph{Referencia débil}: comparar con un PID mal ajustado infla cualquier ahorro. (ii) \emph{Información desigual}: dar previsiones perfectas a uno de los controladores. (iii) \emph{Una sola semana o una sola semilla}, sin dispersión. (iv) \emph{Confundir porcentajes}: un 40~\% de la energía de refrigeración no es un 40~\% de la factura (\cref{eje:5-pue}). (v) \emph{Ignorar el coste de implantación}: sensores, modelado, mantenimiento del modelo y supervisión cuestan más que el algoritmo. (vi) \emph{En reactores, optimizar el rendimiento sin restricciones de seguridad explícitas}, o confiar a la penalización de la recompensa lo que debe garantizar un enclavamiento.
\end{observacion}

\begin{lecturas}
\item J.~Drgoňa \emph{et al.}~\cite{drgona} y Z.~Wang y T.~Hong~\cite{wanghong}, sobre MPC y RL en edificios.
\item D.~Blum \emph{et al.}~\cite{boptest}, «Building optimization testing framework (BOPTEST)».
\item MIT 10.450~\cite{mit10450}, lección~6 «The Exothermic Reactor».
\item R.~Evans y J.~Gao~\cite{evansgao}; C.~Gamble y J.~Gao~\cite{gamblegao}; N.~Lazic \emph{et al.}~\cite{lazic}.
\end{lecturas}

\section*{Ejercicios del tema 5}
\addcontentsline{toc}{section}{Ejercicios del tema 5}

\begin{ejercicio}{1}{eje:5-kpi}
Un proveedor anuncia que su controlador «ahorra un 25~\% de energía» en un edificio. Enumere las preguntas que haría antes de creerlo, relacionándolas con el protocolo de comparación de este tema.
\end{ejercicio}

\begin{ejercicio}{2}{eje:5-ahorro}
Con el \cref{tab:hvac}: (a) calcule el ahorro de energía y de coste del MPC y del agente respecto al PI; (b) calcule el coste total con $\rho=2$ €/(K·h) del MPC sin margen y diga a partir de qué valor de $\rho$ sería preferible al MPC con margen; (c) explique por qué el MPC consume casi la misma energía que el PI pero cuesta menos.
\end{ejercicio}

\begin{ejercicio}{2}{eje:5-crn}
El coste semanal de un controlador tiene desviación típica de $2$ € entre semanas. (a) Calcule la desviación típica de la diferencia de costes entre dos controladores evaluados en semanas independientes y en las mismas semanas con correlación $r=0{,}9$. (b) ¿Cuántas semanas independientes harían falta para igualar la precisión de una sola semana común?
\fuente{adaptado de D.~C. Montgomery, \emph{Design and Analysis of Experiments}, cap.~2 «Simple Comparative Experiments», ventajas del diseño de comparaciones pareadas~\cite{montgomery}.}
\end{ejercicio}

\begin{ejercicio}{2}{eje:5-semenov}
Con los datos del \cref{ej:caso-reactor}: (a) calcule $Q_g$ y $Q_g'$ a $360$ K con $C_A=C_{A,0}$; ¿qué valor de $UA$ haría falta para que un equilibrio a esa temperatura fuera estable? (b) Por ensuciamiento, $UA$ baja a $1{,}5$ kW/K. ¿Qué $T_j$ mantiene el equilibrio a $350$ K y es estable? (c) ¿Qué medida de operación propondría?
\fuente{adaptado de MIT 10.450, primavera de 2006, lección~6 «The Exothermic Reactor», §6.8 (calor generado frente a calor retirado)~\cite{mit10450}.}
\end{ejercicio}

\begin{ejercicio}{2}{eje:5-pue}
Según la publicación de 2016~\cite{evansgao}, reducir un 40~\% la energía de refrigeración equivalió a reducir un 15~\% el exceso de PUE. (a) ¿Qué fracción del exceso de PUE representaba la refrigeración? (b) Con un PUE ilustrativo de $1{,}12$, ¿cuál sería el PUE resultante y la reducción porcentual del consumo total del centro? (c) Comente la diferencia entre los tres porcentajes.
\end{ejercicio}

\begin{ejercicio}{2}{eje:5-bateria}
La batería del \cref{ej:bateria-lp} ($1$ kWh, $1$ kW, $\eta=1$) se opera con precios previstos $(0{,}10;\ 0{,}30;\ 0{,}05;\ 0{,}25)$, pero los precios reales resultan $(0{,}10;\ 0{,}30;\ 0{,}28;\ 0{,}25)$. (a) ¿Qué beneficio obtiene el plan calculado con la previsión si se aplica en lazo abierto? (b) ¿Y un MPC que en cada hora conoce el precio de esa hora y reoptimiza con la previsión para las restantes? (c) ¿Cuál habría sido el óptimo con información perfecta?
\end{ejercicio}

\begin{ejercicio}{3}{eje:5-semenov-dem}
(a) Demuestre que la condición de tangencia entre $Q_g$ y $Q_r$ conduce a $T_j^{*}=T^{*}-RT^{*2}/E$, con $T^{*}$ solución de $Q_g(T)=UA\,RT^{2}/E$. (b) Demuestre que, si existen exactamente dos equilibrios y $Q_g$ es convexa en el intervalo que los contiene, el inferior es estable y el superior inestable. (c) Explique por qué la aproximación de $C_A$ constante es conservadora para la fase inicial del lote.
\fuente{adaptado de MIT 10.450, primavera de 2006, lección~6 «The Exothermic Reactor», §6.8 y §6.10 (estados estacionarios múltiples y su estabilidad)~\cite{mit10450}.}
\end{ejercicio}

\begin{ejercicio}{3}{eje:5-protocolo}
Diseñe el protocolo para comparar en un edificio de oficinas real el PID existente con un MPC durante un invierno: cómo asignar los días a cada controlador, qué indicadores medir y con qué normalización (grados-día, ocupación), qué análisis estadístico aplicar y qué criterios de parada por seguridad o confort establecer.
\end{ejercicio}

% =====================================================================
\chapter{Puesta en producción en planta}
\label{tema:6}
% =====================================================================

Un controlador que funciona en el gemelo digital está a mitad de camino. En planta hay operadores que deben entender y poder anular lo que hace, alarmas que no deben saturarlos, un sistema de control básico y otro de seguridad que no se pueden comprometer, una red industrial que es un objetivo de ciberataques y un marco normativo. Este tema recorre esos requisitos. La arquitectura por capas sigue el capítulo~15 («Architecture and System Design») de Åström y Murray~\cite{am}; la gestión de alarmas, la norma ISA-18.2~\cite{isa182}; la seguridad funcional, IEC~61511~\cite{iec61511}; la ciberseguridad, la guía NIST SP~800-82~\cite{nist} y la serie IEC~62443~\cite{iec62443}; la validación y el seguimiento de modelos, el curso CS329S de Stanford~\cite{cs329s} y Sculley \emph{et al.}~\cite{sculley}, y la regulación, el Reglamento europeo de inteligencia artificial~\cite{aiact}. Las normas son documentos de pago: aquí se usa solo lo que sus organismos difunden públicamente.

\section{Arquitectura por capas}

\begin{figure}[h]
\centering
\begin{tikzpicture}[>=Stealth,every node/.style={font=\small},
  capa/.style={draw,minimum width=8.2cm,minimum height=0.8cm,align=center}]
\node[capa] (opt) at (0,0) {planificación y optimización económica (horas, días)};
\node[capa] (sup) at (0,-1.2) {control supervisor: MPC o política de RL (minutos)};
\node[capa,fill=black!6] (fil) at (0,-2.4) {filtro de seguridad y validación de consignas};
\node[capa] (bas) at (0,-3.6) {control básico: PID en DCS o PLC (segundos)};
\node[capa] (pro) at (0,-4.8) {proceso: sensores y actuadores};
\draw (-6.6,0.4) rectangle (-4.5,-5.2); \node[align=center] at (-5.55,-2.4) {operador\\(HMI,\\alarmas)};
\draw (4.5,-3.2) rectangle (7.2,-5.2); \node[align=center] at (5.85,-4.2) {sistema\\instrumentado\\de seguridad\\(independiente)};
\draw[->] (opt) -- (sup);
\draw[->] (sup) -- node[right,font=\footnotesize]{consignas} (fil);
\draw[->] (fil) -- (bas);
\draw[<->] (bas) -- (pro);
\draw[<->] (pro.east) -- (4.5,-4.8);
\draw[<->] (-4.5,-1.2) -- (sup.west);
\draw[<->] (-4.5,-3.6) -- (bas.west);
\end{tikzpicture}
\caption{Arquitectura por capas: la capa de IA escribe consignas acotadas sobre el control básico; el sistema de seguridad es independiente de ambas.}
\label{fig:capas}
\end{figure}

La \cref{fig:capas} resume el principio de diseño: cada capa trabaja en una escala de tiempo y confía en la inferior. El MPC o la política de RL \emph{no} mueven válvulas: envían consignas a los PID existentes, que siguen funcionando si la capa superior falla. Entre ambas, un filtro valida cada consigna. El \emph{sistema instrumentado de seguridad} (SIS), que detiene el proceso ante condiciones peligrosas, es independiente de todo lo anterior, con sus propios sensores y lógica, y se diseña según IEC~61511~\cite{iec61511}, la aplicación al sector de procesos de la seguridad funcional de IEC~61508, que asigna a cada función de seguridad un nivel de integridad (SIL, de 1 a 4). Ningún componente de aprendizaje automático forma parte del SIS. El caso de Google~\cite{gamblegao} sigue el mismo esquema: la IA en la nube propone, el sistema local verifica y aplica, y las reglas clásicas quedan como respaldo.

\section{Modos de operación y lazo humano}

\begin{definicion}[Modos de operación]\label{def:modos}
Se distinguen tres modos de creciente autonomía: \emph{asesor} (el sistema recomienda y el operador decide y ejecuta), \emph{supervisado} (el sistema ejecuta dentro de límites y el operador vigila y puede intervenir en cualquier momento) y \emph{autónomo con supervisión} (el sistema opera con salvaguardas automáticas y el operador interviene por excepción). En todos ellos el operador puede volver al control básico.
\end{definicion}

Un despliegue prudente recorre los modos en orden y solo avanza cuando hay evidencia: el caso de los centros de datos pasó de recomendaciones en 2016~\cite{evansgao} a control autónomo con supervisión en 2018~\cite{gamblegao}. El paso entre modos no debe provocar saltos en los actuadores.

\begin{proposicion}[Transferencia sin saltos]\label{prop:bumpless}
Sea un PI en forma posicional, $u_k=K_ce_k+I_k$, que toma el control en el instante $k_0$ de un sistema que aplicaba $u_{k_0-1}$. Si se inicializa $I_{k_0}=u_{k_0-1}-K_ce_{k_0}$, entonces $u_{k_0}=u_{k_0-1}$.
\end{proposicion}

\begin{proof}
Por sustitución: $u_{k_0}=K_ce_{k_0}+u_{k_0-1}-K_ce_{k_0}=u_{k_0-1}$. A partir de ahí el PI evoluciona con su ley; la salida es continua en el cambio.
\end{proof}

Con la zona del \cref{ej:caso-pi}, si la capa supervisora aplicaba $1{,}2$ kW y el error al recuperar el PI es de $0{,}4$ K, el integrador se inicializa en $1{,}2-1{,}429\cdot0{,}4=0{,}629$ kW. El mismo principio se aplica al pasar del PI al MPC: el MPC debe penalizar los cambios de la acción respecto a la última aplicada.

\section{Alarmas}

La norma ANSI/ISA-18.2, base de la internacional IEC~62682~\cite{isa182}, organiza la gestión de alarmas en un ciclo de vida (filosofía, identificación, racionalización, diseño, implantación, operación, mantenimiento, seguimiento, gestión del cambio y auditoría) y da referencias de carga: alrededor de 150 alarmas al día (una cada 10 minutos de media) se considera muy probablemente aceptable y alrededor de 300 al día (dos cada 10 minutos), el máximo manejable. Una alarma solo debe existir si exige una acción del operador. Una capa de IA añade fuentes nuevas que hay que racionalizar con el mismo criterio: residuos del modelo fuera de banda, proximidad a restricciones, salidas a modo de respaldo, pérdida de comunicación con la capa supervisora o predicciones con incertidumbre alta.

\section{Validación y despliegue}

\begin{definicion}[Etapas de validación]\label{def:validar}
(i) \emph{Validación fuera de línea}: el controlador se evalúa en el gemelo y sobre históricos, con escenarios no usados en su diseño. (ii) \emph{Modo sombra}: el controlador calcula sus decisiones con datos en tiempo real, pero no las aplica; se comparan con las del sistema existente y se comprueba que serían seguras. (iii) \emph{Experimento en planta}: se alterna el controlador nuevo y el existente por periodos (días o semanas) y se mide la diferencia con un diseño de comparación pareada~\cite[cap.~2]{montgomery}. (iv) \emph{Despliegue con seguimiento}: se vigilan los residuos del modelo, la distribución de las entradas y los indicadores, con criterios de vuelta atrás definidos de antemano.
\end{definicion}

\begin{ejemplo}[Contraste del ahorro]\label{ej:contraste}
Durante 20 días laborables se alternan el MPC y el PI en días emparejados por temperatura exterior. Las diferencias diarias de coste (PI menos MPC) tienen media $\bar d=1{,}20$ € y desviación típica $s=1{,}50$ €. El estadístico es $t=\bar d/(s/\sqrt{20})=1{,}20/0{,}3354=3{,}58$, con 19 grados de libertad: $p=0{,}002$. El intervalo de confianza del 95~\% es $1{,}20\pm2{,}093\cdot0{,}3354=[0{,}50;\ 1{,}90]$ € al día (módulo~0.9, §4.4 y §5.3). El ahorro es significativo, pero su magnitud es incierta en un factor cercano a~4: conviene seguir midiendo.
\end{ejemplo}

Los modelos degradan: cambian los equipos, el uso del edificio o el producto, y las relaciones aprendidas dejan de valer. Sculley \emph{et al.}~\cite{sculley} describen cómo un sistema de aprendizaje automático acumula «deuda técnica» por dependencias de datos ocultas, bucles de realimentación (el controlador cambia los datos con los que se reentrenará) y configuración, y las clases~10 y~11 de CS329S~\cite{cs329s} tratan la detección de cambios de distribución y el seguimiento. En control, la señal de seguimiento más útil es el residuo del modelo en simulación libre sobre ventanas recientes: un sesgo sostenido indica que hay que recalibrar el gemelo.

\section{El filtro de seguridad}

\begin{definicion}[Filtro de seguridad]\label{def:filtro}
Sea un modelo $x_{k+1}=f(x_k,u_k,w_k)$ con perturbación acotada $w_k\in W$ y un conjunto seguro $\Sset_{\mathrm{seg}}$. Para cada $x$ se define el conjunto de acciones seguras $\Uset_{\mathrm{seg}}(x)=\{u\in\Uset:\ f(x,u,w)\in\Sset_{\mathrm{seg}}\ \forall w\in W\}$. El \emph{filtro de seguridad} sustituye la acción propuesta $u^{\mathrm{p}}$ por su proyección $u=\proy_{\Uset_{\mathrm{seg}}(x)}(u^{\mathrm{p}})$.
\end{definicion}

\begin{proposicion}[Invariancia bajo el filtro]\label{prop:invariancia}
Si $\Uset_{\mathrm{seg}}(x)\ne\emptyset$ para todo $x\in\Sset_{\mathrm{seg}}$ y $x_0\in\Sset_{\mathrm{seg}}$, entonces, cualquiera que sea la política que propone las acciones, $x_k\in\Sset_{\mathrm{seg}}$ para todo $k$ y toda secuencia de perturbaciones en $W$.
\end{proposicion}

\begin{proof}
Por inducción. Si $x_k\in\Sset_{\mathrm{seg}}$, el conjunto $\Uset_{\mathrm{seg}}(x_k)$ no es vacío y la acción filtrada pertenece a él, luego $x_{k+1}=f(x_k,u_k,w_k)\in\Sset_{\mathrm{seg}}$ para todo $w_k\in W$, por definición de $\Uset_{\mathrm{seg}}$.
\end{proof}

La garantía es tan buena como el modelo y la cota de las perturbaciones: el filtro usa un modelo sencillo, conservador y verificable, no la red del agente. Cuando $\Uset_{\mathrm{seg}}$ es un intervalo, la proyección es un simple recorte. Se completa con límites de velocidad de cambio (Lazic \emph{et al.}~\cite{lazic} limitaban la variación de cada actuador entre pasos), un \emph{perro guardián} (\emph{watchdog}) que devuelve el control a los PID si la capa supervisora deja de enviar consignas válidas, y la transferencia sin saltos de la \cref{prop:bumpless}.

\begin{ejemplo}[Caso de estudio: filtro para la zona]\label{ej:filtro}
Con el modelo de la zona, $T_{k+1}=aT_k+bP_k+d_k$, $a=0{,}9835$, $b=0{,}23$ K/kW (COP de $2{,}8$) y la perturbación $d_k=(1-a)(\Text+RG)$ acotada por $\Text\in[-5,15]$ °C y $G\in[0,2]$ kW, es decir, $d\in[-0{,}0827;\ 0{,}4133]$ K, se toma $\Sset_{\mathrm{seg}}=[17,26]$ °C. El intervalo seguro es
\[
\Uset_{\mathrm{seg}}(T)=\Bigl[\max\Bigl\{0,\ \tfrac{17-aT-d_{\min}}{b}\Bigr\},\ \min\Bigl\{3,\ \tfrac{26-aT-d_{\max}}{b}\Bigr\}\Bigr].
\]
En $T=17$ °C vale $[1{,}579;\ 3]$ kW y en $T=26$ °C, $[0;\ 0{,}068]$ kW. Como los extremos son afines en $T$, no es vacío en todo $[17,26]$, y por la \cref{prop:invariancia} ninguna política, por mala que sea, saca la zona de ese intervalo mientras el modelo y las cotas sean válidos. En el reactor del \cref{ej:caso-reactor}, el filtro análogo limita la temperatura de camisa por debajo de la crítica y fuerza refrigeración máxima cerca del equilibrio inestable (cuaderno, Parte~9).
\end{ejemplo}

\section{Ciberseguridad de la tecnología de operación}

Los sistemas de control industrial (tecnología de operación, OT) priorizan la disponibilidad y la seguridad física por encima de la confidencialidad, funcionan durante décadas con equipos difíciles de actualizar y usan protocolos diseñados sin autenticación. La guía NIST SP~800-82, revisión~3, de 2023~\cite{nist}, adapta la gestión de riesgos de ciberseguridad a este entorno. La serie IEC~62443~\cite{iec62443}, desarrollada por IEC e ISA, se organiza en cuatro bloques (general; políticas y procedimientos; sistema; componentes) y aporta dos conceptos centrales: la partición del sistema en \emph{zonas} de activos con requisitos de seguridad comunes, comunicadas solo por \emph{conductos} definidos, y los \emph{niveles de seguridad} SL~1 a SL~4, ligados a la capacidad del adversario que se quiere resistir, desglosados en siete requisitos fundamentales (identificación y autenticación, control de uso, integridad del sistema, confidencialidad de los datos, flujo de datos restringido, respuesta oportuna a eventos y disponibilidad de recursos).

Para una capa de IA, de ahí se derivan reglas prácticas: (i) la capa supervisora vive en su propia zona y solo escribe consignas por un conducto autenticado y registrado, nunca acceso directo a los controladores; (ii) el sistema local valida cada consigna (el filtro de seguridad es también una defensa: limita el daño de una consigna maliciosa); (iii) si el modelo se ejecuta en la nube, la planta debe poder operar sin ella; (iv) los datos de entrenamiento son un activo: un atacante que los manipule puede inducir un modelo que se comporte mal en condiciones concretas, y (v) actualizar un modelo es un cambio en el sistema de control y pasa por la gestión del cambio.

\section{Normas y regulación}

\begin{table}[htbp]
\centering\small
\begin{tabularx}{\textwidth}{@{}l>{\raggedright\arraybackslash}X>{\raggedright\arraybackslash}X@{}}
\toprule
Referencia & Ámbito & Qué implica para una capa de IA \\
\midrule
IEC 61508 / IEC 61511~\cite{iec61511} & seguridad funcional; sistemas instrumentados de seguridad en procesos & la IA queda fuera del SIS; el análisis de riesgos debe considerar sus modos de fallo \\
ANSI/ISA-18.2, IEC 62682~\cite{isa182} & gestión de alarmas & racionalizar las alarmas nuevas; vigilar la carga del operador \\
IEC 62443~\cite{iec62443}, NIST SP 800-82~\cite{nist} & ciberseguridad de sistemas de automatización & zonas y conductos; autenticación y registro de consignas \\
Reglamento (UE) 2024/1689~\cite{aiact} & inteligencia artificial & los sistemas destinados a ser componentes de seguridad en la gestión y operación de infraestructuras críticas, como el suministro de agua, gas, calefacción o electricidad, se clasifican como de alto riesgo (anexo~III, punto~2) \\
\bottomrule
\end{tabularx}
\caption{Marco normativo de referencia.}
\label{tab:normas}
\end{table}

Los sistemas de alto riesgo del Reglamento de IA deben cumplir los requisitos de la sección~2 de su capítulo~III (artículos~8 a~15), entre ellos la gestión de riesgos, la gobernanza de los datos, el registro de eventos, la supervisión humana (artículo~14: el sistema debe diseñarse para que las personas puedan entender sus límites, detectar anomalías, no confiar en exceso en sus salidas, anular sus decisiones y detenerlo de forma segura) y la precisión, robustez y ciberseguridad. Un optimizador de climatización de oficinas no suele ser un componente de seguridad; el control de una red de calefacción urbana o de una subestación puede serlo. La clasificación de cada caso es una cuestión jurídica que debe evaluarse con especialistas; lo que el módulo aporta es que la arquitectura de la \cref{fig:capas} facilita cumplir esos requisitos.

\begin{table}[htbp]
\centering\footnotesize
\begin{tabularx}{\textwidth}{@{}>{\raggedright\arraybackslash}p{3.1cm}>{\raggedright\arraybackslash}X>{\raggedright\arraybackslash}p{2.0cm}>{\raggedright\arraybackslash}X@{}}
\toprule
Herramienta & Para qué sirve & Licencia o tipo & Cuándo elegirla \\
\midrule
OPC UA & comunicación industrial con seguridad (autenticación, cifrado) & estándar abierto & leer medidas y escribir consignas desde la capa supervisora \\
Historiadores de proceso & almacenamiento de series temporales de planta & comercial o libre & datos de identificación y de seguimiento \\
MLflow u otros registros de modelos & versionado de modelos, datos y métricas & libre (Apache-2.0) & trazabilidad de qué modelo controló la planta y cuándo \\
Contenedores y orquestación & despliegue reproducible de la capa supervisora & libre & separar entornos de prueba y producción \\
Evidently, WhyLogs & seguimiento de distribuciones y de la deriva de los datos & libres & alarmas de degradación del modelo \\
\bottomrule
\end{tabularx}
\caption{Tecnologías para la puesta en producción.}
\label{tab:tec6}
\end{table}

\begin{observacion}[Errores frecuentes en la práctica]\label{obs:errores6}
(i) \emph{Escribir directamente en los actuadores} desde la capa de IA, sin PID por debajo ni filtro. (ii) \emph{No definir el modo de respaldo} ni probar la transferencia, que falla precisamente el día en que hace falta. (iii) \emph{Inundar de alarmas} al operador con avisos del modelo que no exigen acción. (iv) \emph{Medir el ahorro antes-después} sin grupo de control: el tiempo meteorológico o la producción cambian a la vez. (v) \emph{Olvidar el seguimiento}: el modelo se degrada en silencio. (vi) \emph{Conectar la planta a la nube} sin zonas, conductos ni autenticación. (vii) \emph{Tratar la actualización de un modelo como un despliegue de software ordinario}, sin gestión del cambio ni validación en sombra.
\end{observacion}

\begin{lecturas}
\item K.~J. Åström y R.~M. Murray~\cite{am}, cap.~15 «Architecture and System Design».
\item NIST SP~800-82r3~\cite{nist}, \emph{Guide to Operational Technology (OT) Security}.
\item C.~Gamble y J.~Gao~\cite{gamblegao}, sobre salvaguardas en el control autónomo de un centro de datos.
\item D.~Sculley \emph{et al.}~\cite{sculley}, «Hidden technical debt in machine learning systems».
\item CS329S~\cite{cs329s}, clases~8 «Deployment» y~10 sobre fallos, cambios de distribución y seguimiento.
\end{lecturas}

\section*{Ejercicios del tema 6}
\addcontentsline{toc}{section}{Ejercicios del tema 6}

\begin{ejercicio}{1}{eje:6-modos}
Clasifique en asesor, supervisado o autónomo con supervisión: (a) un panel que propone cada mañana las consignas del día y el jefe de turno las introduce; (b) un MPC que escribe consignas cada 15 minutos con límites y que el operador puede desactivar; (c) el sistema de 2018 del centro de datos de Google descrito en el \cref{tema:5}. Justifique qué evidencia pediría para pasar de (a) a (b).
\end{ejercicio}

\begin{ejercicio}{2}{eje:6-capas}
Para la zona climatizada, asigne cada función a una capa de la \cref{fig:capas} y una escala de tiempo: seguimiento de la temperatura de impulsión; elección de las consignas horarias; parada de la bomba de calor por alta presión del refrigerante; previsión del precio de mañana; detección de que una sonda está averiada. Explique por qué la tercera no puede depender de la capa supervisora.
\fuente{adaptado de K.~J. Åström y R.~M. Murray, \emph{Feedback Systems}, 2.ª ed., cap.~15 «Architecture and System Design», arquitecturas jerárquicas de control~\cite{am}.}
\end{ejercicio}

\begin{ejercicio}{2}{eje:6-alarmas}
Una sala de control recibe de media 9 alarmas por hora en operación normal. (a) ¿Cuántas al día y cada 10 minutos? Compárelo con las referencias de ISA-18.2. (b) El nuevo MPC añadiría 3 alarmas por hora. ¿Qué ocurre con la carga? (c) Proponga cómo racionalizarlas.
\end{ejercicio}

\begin{ejercicio}{2}{eje:6-contraste}
En el \cref{ej:contraste}: (a) verifique $t$, el valor~$p$ y el intervalo; (b) ¿qué desviación típica de las diferencias haría el ahorro no significativo al 5~\% con la misma media? (c) ¿Cuántos días harían falta para que la semiamplitud del intervalo fuera de $0{,}3$ € si $s=1{,}5$ €? (d) ¿Por qué el diseño emparejado es preferible a comparar 20 días de MPC de este invierno con 20 días de PI del invierno pasado?
\fuente{adaptado de D.~C. Montgomery, \emph{Design and Analysis of Experiments}, cap.~2 «Simple Comparative Experiments», diseño de comparaciones pareadas~\cite{montgomery}.}
\end{ejercicio}

\begin{ejercicio}{2}{eje:6-deriva}
El residuo en simulación libre del gemelo de la zona tiene, con el modelo bien calibrado, media nula y desviación típica $0{,}14$ K por muestra; se vigila la media diaria de 96 muestras. (a) Suponiendo independencia, calcule el error típico de la media diaria y un umbral de alarma de 3 errores típicos. (b) ¿Por qué la hipótesis de independencia es optimista en este caso y cómo afecta al umbral? (c) Proponga una regla que dispare la recalibración.
\fuente{adaptado de Stanford CS329S, invierno de 2022, clase~10 «Diagnosis of ML system failures \& data distribution shifts \& monitoring»~\cite{cs329s}.}
\end{ejercicio}

\begin{ejercicio}{2}{eje:6-filtro}
Con los datos del \cref{ej:filtro}: (a) calcule $\Uset_{\mathrm{seg}}(T)$ para $T=20$ °C y $T=25{,}9$ °C; (b) el agente propone $P=0{,}5$ kW a $25{,}9$ °C y $P=0$ kW a $17$ °C: ¿qué aplica el filtro?; (c) si la cota superior de $\Text$ fuera $35$ °C, ¿seguiría siendo invariante $[17,26]$?
\end{ejercicio}

\begin{ejercicio}{3}{eje:6-invariancia}
(a) Demuestre que, para un modelo escalar $x_{k+1}=ax_k+bu_k+w_k$ con $a>0$, $b>0$, $u\in[0,\bar u]$ y $w\in[w_{\min},w_{\max}]$, el intervalo $[x_{\min},x_{\max}]$ es invariante bajo el filtro si y solo si $\Uset_{\mathrm{seg}}$ no es vacío en sus dos extremos, y dé las dos condiciones explícitas. (b) Añada un límite de velocidad $\abs{u_k-u_{k-1}}\le\delta$ y explique por qué la invariancia ya no puede comprobarse solo en función de $x$.
\end{ejercicio}

\begin{ejercicio}{3}{eje:6-plan}
Redacte el plan de puesta en producción del MPC de la zona en un edificio real: arquitectura y punto de inserción, modos y criterios de avance, alarmas nuevas, filtro y modo de respaldo, experimento de validación, seguimiento y recalibración, ciberseguridad y análisis regulatorio. Indique, para cada elemento, qué sección del módulo lo fundamenta.
\end{ejercicio}


\appendix
% =====================================================================
\chapter{Soluciones de los ejercicios}
\label{ap:soluciones}
% =====================================================================

Cada solución se presenta en cinco partes: datos, planteamiento, desarrollo, resultado y comprobación. En los pasos intermedios se redondea a cuatro cifras significativas salvo que se indique otra cosa. Los ejercicios de diseño admiten varias respuestas correctas; se da una razonada.

% ---------------------------------------------------------------------
\section{Tema 1}

\begin{solucion}{eje:1-variables}
\datos reactor con camisa; se miden $T$ y $T_j$; se manipula la válvula; varía la temperatura del agua de red; no se mide $C_A$.
\planteamiento aplicar la \cref{def:ee}: el estado es lo que hay que conocer para predecir; las entradas se dividen en elegidas y no elegidas.
\pasos
\begin{enumerate}[label=\arabic*.]
  \item Estados: $T$, $T_j$ y $C_A$ (cada uno tiene su propio balance, de energía o de materia).
  \item Entrada manipulada: la apertura de la válvula de refrigeración.
  \item Perturbación: la temperatura del agua de red.
  \item Salidas medidas: $T$ y $T_j$.
  \item $C_A$ es un estado no medido: un MPC necesita estimarlo (observador, filtro de Kalman extendido o estimación de horizonte móvil), por ejemplo a partir del balance de energía, porque el calor generado delata la velocidad de reacción.
\end{enumerate}
\resultado \textbf{estados $T,T_j,C_A$; entrada, la válvula; perturbación, el agua de red; salidas, $T$ y $T_j$}. El estado no medido obliga a usar un estimador.
\comprobacion el número de estados coincide con el número de ecuaciones diferenciales del modelo (\cref{def:reactor} más el balance de la camisa).
\end{solucion}

\begin{solucion}{eje:1-zoh}
\datos $R=5$ K/kW; $C=3$ kWh/K; $\tau=15$ h; $\Dt=1$ h.
\planteamiento \cref{prop:zoh} con $A=-1/\tau$ y $B=1/C$ para la entrada $Q$; Euler: $a_E=1-\Dt/\tau$, $b_E=\Dt/C$.
\pasos
\begin{enumerate}[label=\arabic*.]
  \item $a=e^{-1/15}=0{,}93551$.
  \item $b=A^{-1}(a-1)B=(-15)(0{,}93551-1)/3=R(1-a)=5\cdot0{,}06449=0{,}32247$ K/kWh.
  \item $a_E=1-1/15=0{,}93333$; $b_E=1/3=0{,}33333$ K/kWh.
  \item Errores relativos: $(0{,}93333-0{,}93551)/0{,}93551=-0{,}23$~\%; $(0{,}33333-0{,}32247)/0{,}32247=+3{,}37$~\%.
\end{enumerate}
\resultado $\mathbf{a=0{,}9355}$, $\mathbf{b=0{,}3225}$ K/kWh; Euler, $0{,}9333$ y $0{,}3333$, con errores del $-0{,}23$~\% y $+3{,}4$~\%.
\comprobacion la ganancia estática discreta $b/(1-a)=0{,}32247/0{,}06449=5{,}000=R$ K/kW coincide con la continua; con Euler, $b_E/(1-a_E)=0{,}3333/0{,}0667=5$ también: Euler conserva la ganancia estática pero no la dinámica.
\end{solucion}

\begin{solucion}{eje:1-rigidez}
\datos $C_a=0{,}5$, $C_m=5$ kWh/K; $R_{am}=0{,}5$, $R_{ao}=5$ K/kW.
\planteamiento balances de energía de cada nodo; autovalores de una matriz $2\times2$ por traza y determinante; \cref{cor:euler}.
\pasos
\begin{enumerate}[label=\arabic*.]
  \item $C_a\dot T_a=(T_m-T_a)/R_{am}+(\Text-T_a)/R_{ao}$ y $C_m\dot T_m=(T_a-T_m)/R_{am}$. Sin entradas:
  $A=\begin{pmatrix}-(2+0{,}2)/0{,}5&2/0{,}5\\2/5&-2/5\end{pmatrix}=\begin{pmatrix}-4{,}4&4\\0{,}4&-0{,}4\end{pmatrix}$.
  \item $\tr A=-4{,}8$, $\det A=1{,}76-1{,}6=0{,}16$; $\lambda=(-4{,}8\pm\sqrt{23{,}04-0{,}64})/2=(-4{,}8\pm4{,}7329)/2$, es decir, $\lambda_1=-4{,}7664$ y $\lambda_2=-0{,}03357$ h$^{-1}$.
  \item Euler con $\Dt=0{,}25$: $1+\lambda_i\Dt=-0{,}1916$ y $0{,}9916$; ambos de módulo menor que~1: estable.
  \item Euler con $\Dt=0{,}5$: $1-2{,}3832=-1{,}3832$ y $0{,}9832$: inestable.
  \item $\Dt_{\max}=2/4{,}7664=0{,}4196$ h.
\end{enumerate}
\resultado \textbf{$\lambda=-4{,}766$ y $-0{,}0336$ h$^{-1}$; Euler es estable con $\Dt=0{,}25$ h e inestable con $0{,}5$ h; $\Dt_{\max}=0{,}42$ h.}
\comprobacion $\lambda_1+\lambda_2=-4{,}8=\tr A$ y $\lambda_1\lambda_2=0{,}160=\det A$.
\end{solucion}

\begin{solucion}{eje:1-tanque}
\datos $V=2$ m³; $F=0{,}5$ m³/h; $F_1/F=0{,}6$; $UA=1{,}5$ kW/K; $\rho c_p=4180$ kJ/(m³·K); $K_c=2$.
\planteamiento fórmulas del enunciado (lección~5 de 10.450) con unidades coherentes; error permanente de la \cref{prop:offset} con ganancia de lazo $K_3K_c$.
\pasos
\begin{enumerate}[label=\arabic*.]
  \item $\tau_R=V/F=2/0{,}5=4$ h.
  \item $\rho c_pF=4180\cdot0{,}5/3600=0{,}5806$ kW/K; $\beta=1{,}5/0{,}5806=2{,}584$.
  \item $\tau=4/3{,}584=1{,}116$ h; $K_1=0{,}6/3{,}584=0{,}1674$; $K_3=2{,}584/3{,}584=0{,}7210$.
  \item Ganancia de lazo: $K_3K_c=1{,}442$; fracción de error permanente: $1/(1+1{,}442)=0{,}4095$.
\end{enumerate}
\resultado $\mathbf{\tau=1{,}12}$ h, $\mathbf{K_1=0{,}167}$, $\mathbf{K_3=0{,}721}$; con control P queda un $\mathbf{41}$~\% del escalón como error.
\comprobacion $\tau<\tau_R$, como debe: el intercambio de calor acelera la respuesta; y $K_3<1$ porque el vapor solo actúa a través del intercambiador.
\end{solucion}

\begin{solucion}{eje:1-offset}
\datos $k=14$ K/kW; $\Delta r=2$ K.
\planteamiento \cref{prop:offset} y \cref{prop:integral}.
\pasos
\begin{enumerate}[label=\arabic*.]
  \item (a) $e_\infty=2/(1+14\cdot0{,}5)=2/8=0{,}25$ K.
  \item (b) $2/(1+14K_c)<0{,}1\iff1+14K_c>20\iff K_c>1{,}357$ kW/K.
  \item (c) Una ganancia alta satura el actuador ante cualquier cambio, amplifica el ruido de medida y reduce el margen de estabilidad cuando hay retardos; aun así, el error nunca es nulo. La acción integral anula el error en equilibrio para cualquier ganancia que estabilice el lazo, lo que permite elegir $K_c$ por robustez y no por precisión.
\end{enumerate}
\resultado \textbf{(a) $0{,}25$ K; (b) $K_c>1{,}36$ kW/K}; (c) la solución es añadir acción integral, no subir la ganancia.
\comprobacion con $K_c=1{,}357$, $2/(1+19)=0{,}1$ K, en el límite.
\end{solucion}

\begin{solucion}{eje:1-simc}
\datos $k=2$ K/\%; $\tau=10$ min; $\theta=2$ min; $\Dt=0{,}5$ min.
\planteamiento reglas \eqref{eq:simc} y PI posicional de la \cref{def:pid}.
\pasos
\begin{enumerate}[label=\arabic*.]
  \item (a) $\tau_c=2$: $K_c=10/(2\cdot4)=1{,}25$ \%/K; $\tau_I=\min\{10,16\}=10$ min.
  \item (b) $\tau_c=6$: $K_c=10/(2\cdot8)=0{,}625$ \%/K; $\tau_I=\min\{10,32\}=10$ min.
  \item (c) $u_0=K_ce_0+I_0=1{,}25\cdot1+0=1{,}25$ \%; $I_1=I_0+K_c\Dt e_0/\tau_I=1{,}25\cdot0{,}5/10=0{,}0625$ \%.
\end{enumerate}
\resultado \textbf{(a) $K_c=1{,}25$ \%/K, $\tau_I=10$ min; (b) $K_c=0{,}625$ \%/K, $\tau_I=10$ min; (c) $u_0=1{,}25$ \%, $I_1=0{,}0625$ \%.}
\comprobacion duplicar $\tau_c+\theta$ divide $K_c$ por dos, como indica \eqref{eq:simc}; $\tau_I$ queda limitado por $\tau$ en ambos casos.
\end{solucion}

\begin{solucion}{eje:1-integral}
\datos $x_{k+1}=ax_k+bu_k+w$; $u_k=K_c(r-x_k)+I_k$; $I_{k+1}=I_k+\kappa(r-x_k)$ con $\kappa=K_c\Dt/\tau_I$.
\planteamiento escribir el lazo como sistema lineal en $(x,I)$ y estudiar su equilibrio.
\pasos
\begin{enumerate}[label=\arabic*.]
  \item (a) $\begin{pmatrix}x_{k+1}\\I_{k+1}\end{pmatrix}=\begin{pmatrix}a-bK_c&b\\-\kappa&1\end{pmatrix}\begin{pmatrix}x_k\\I_k\end{pmatrix}+\begin{pmatrix}bK_cr+w\\\kappa r\end{pmatrix}$.
  \item (b) Si la matriz tiene radio espectral menor que~1 (\cref{teo:estab}), existe un único equilibrio y el lazo converge a él. En el equilibrio, la segunda ecuación da $I=I+\kappa(r-x)$, luego $x=r$, sea cual sea $w$. La primera fija entonces $I^{*}=\bigl((1-a)r-w\bigr)/b$: el integrador «aprende» la perturbación.
  \item (c) Ejemplo: $a=0{,}9$, $b=1$, $u_{\max}=1$, $w=-0{,}2$, $r=10$. Con $u=1$, el equilibrio de la planta es $x=(1-0{,}2)/(1-0{,}9)=8<10$. El error tiende a $2$ y, sin integración condicional, $I_{k+1}=I_k+2\kappa$ crece sin límite, aunque $u$ siga en~1. Si después la perturbación desaparece, $I$ es enorme y la salida permanece saturada mucho tiempo: sobreoscilación por \emph{windup}.
\end{enumerate}
\resultado \textbf{con lazo estable y sin saturación, el equilibrio cumple $y=r$ para todo $w$; con saturación, el integrador diverge si no hay antisaturación.}
\comprobacion es la \cref{prop:integral} en un caso concreto; en el cuaderno, la comprobación del ejercicio~3 verifica que el integrador no crece con la salida saturada.
\end{solucion}

\begin{solucion}{eje:1-estab}
\datos $\dot x=-x/\tau$; Euler explícito e implícito; modelo rígido del \cref{ej:rigido}.
\planteamiento estudiar el factor de amplificación por paso; para el implícito, diagonalizar.
\pasos
\begin{enumerate}[label=\arabic*.]
  \item (a) Euler: $x_{k+1}=(1-\Dt/\tau)x_k$; $\abs{1-\Dt/\tau}<1\iff0<\Dt/\tau<2\iff\Dt<2\tau$. Si $\tau<\Dt<2\tau$, el factor está en $(-1,0)$ y el signo alterna en cada paso, mientras que la solución exacta $x_0e^{-t/\tau}$ nunca cambia de signo.
  \item (b) Implícito: $(I-\Dt A)x_{k+1}=x_k$. Con $A=V\Lambda V^{-1}$, $(I-\Dt A)^{-1}=V(I-\Dt\Lambda)^{-1}V^{-1}$, de autovalores $1/(1-\lambda_i\Dt)$; si $\lambda_i<0$, $1-\lambda_i\Dt>1$ y el autovalor está en $(0,1)$ para todo $\Dt>0$.
  \item (c) Para un modelo lineal, la discretización exacta ZOH (se calcula $e^{A\Dt}$ una vez); si el modelo es no lineal y rígido, un integrador implícito o la colocación ortogonal, que CasADi y do-mpc ofrecen. Euler explícito obligaría a pasos menores que $0{,}42$ h, dictados por la dinámica rápida del aire, aunque interese la lenta.
\end{enumerate}
\resultado \textbf{Euler explícito: estable si y solo si $\Dt<2\tau$; Euler implícito: estable para todo $\Dt$; para el MPC, ZOH exacto o integradores implícitos.}
\comprobacion con $\tau=0{,}21$ h (dinámica rápida del \cref{ej:rigido}), $2\tau=0{,}42$ h, el mismo límite obtenido con los autovalores.
\end{solucion}

% ---------------------------------------------------------------------
\section{Tema 2}

\begin{solucion}{eje:2-unpaso}
\planteamiento \cref{def:validacion} y \cref{prop:acumula}.
\pasos
\begin{enumerate}[label=\arabic*.]
  \item A un paso, el modelo parte de la medida real en cada instante: sus errores no se acumulan y, en procesos lentos, predecir $T_{k+1}\approx T_k$ ya da un error pequeño.
  \item El MPC encadena 96 predicciones: el error de cada una se suma, amplificado por la dinámica lenta (factor $1/(1-L)$, cercano a 60 en la zona).
  \item Métrica: RMSE (y error máximo) en simulación libre sobre horizontes de la longitud del MPC, en días de validación separados en el tiempo de los de estimación, y con condiciones de operación representativas.
\end{enumerate}
\resultado \textbf{el error a un paso no mide lo que el MPC usa; hay que validar en simulación libre sobre bloques temporales no vistos.}
\comprobacion en el \cref{ej:caso-id}, los tres modelos tienen $0{,}137$ K a un paso y difieren en un factor de casi~2 en simulación.
\end{solucion}

\begin{solucion}{eje:2-arx}
\datos $y=(1;\ 1{,}2;\ 0{,}95;\ 1{,}18)$; $u=(1,0,1)$.
\planteamiento \cref{prop:ls} con $\varphi_k=(y_k,u_k)$ e $Y=(y_1,y_2,y_3)$.
\pasos
\begin{enumerate}[label=\arabic*.]
  \item Filas de $\Phi$: $(1,1)$, $(1{,}2;\ 0)$, $(0{,}95;\ 1)$.
  \item Productos:
  \[
  \Phi\T\Phi=\begin{pmatrix}1+1{,}44+0{,}9025&1+0+0{,}95\\1{,}95&2\end{pmatrix}=\begin{pmatrix}3{,}3425&1{,}95\\1{,}95&2\end{pmatrix},
  \]
  \[
  \Phi\T Y=\begin{pmatrix}1{,}2+1{,}14+1{,}121\\1{,}2+0+1{,}18\end{pmatrix}=\begin{pmatrix}3{,}461\\2{,}38\end{pmatrix}.
  \]
  \item $\det=6{,}685-3{,}8025=2{,}8825$; $\hat a=(2\cdot3{,}461-1{,}95\cdot2{,}38)/2{,}8825=2{,}281/2{,}8825=0{,}7913$; $\hat b=(3{,}3425\cdot2{,}38-1{,}95\cdot3{,}461)/2{,}8825=1{,}2062/2{,}8825=0{,}4185$.
  \item Residuos: $1{,}2-(0{,}7913+0{,}4185)=-0{,}0098$; $0{,}95-0{,}7913\cdot1{,}2=0{,}0004$; $1{,}18-(0{,}7913\cdot0{,}95+0{,}4185)=0{,}0098$.
  \item Simulación desde $y_0=1$: $1{,}2098$; $0{,}7913\cdot1{,}2098=0{,}9573$; $0{,}7913\cdot0{,}9573+0{,}4185=1{,}1760$.
\end{enumerate}
\resultado $\mathbf{\hat a=0{,}791}$, $\mathbf{\hat b=0{,}418}$; simulación $(1;\ 1{,}210;\ 0{,}957;\ 1{,}176)$ frente a los datos $(1;\ 1{,}2;\ 0{,}95;\ 1{,}18)$.
\comprobacion los residuos son ortogonales a las columnas de $\Phi$: $1\cdot(-0{,}0098)+1{,}2\cdot0{,}0004+0{,}95\cdot0{,}0098\approx0$ y $-0{,}0098+0+0{,}0098=0$.
\end{solucion}

\begin{solucion}{eje:2-escalar}
\datos $y_k=\theta_0u_k+e_k$; $\Var e_k=1$; $N=400$.
\planteamiento \cref{prop:ls} con un solo parámetro.
\pasos
\begin{enumerate}[label=\arabic*.]
  \item (a) $\hat\theta_N=\sum u_ky_k/\sum u_k^{2}=\theta_0+\sum u_ke_k/\sum u_k^{2}$. Si $\frac1N\sum u_k^{2}\to\bar\sigma_u^{2}>0$ y $\frac1N\sum u_ke_k\to0$ (independencia), $\hat\theta_N\to\theta_0$.
  \item (b) Condicionado a $u$, $\Var\hat\theta_N=\sum u_k^{2}\Var e/(\sum u_k^{2})^{2}=\lambda/\sum u_k^{2}$.
  \item (c) $u_k=\pm1$: $\sum u_k^{2}=400$ y la desviación típica es $1/\sqrt{400}=0{,}05$. Sinusoide de amplitud $0{,}1$: valor cuadrático medio $0{,}01/2$, $\sum u_k^{2}\approx2$ y desviación típica $1/\sqrt2=0{,}71$.
\end{enumerate}
\resultado \textbf{$\hat\theta_N\to\theta_0$, $\Var\hat\theta_N=\lambda/\sum u_k^{2}$; desviación típica $0{,}05$ frente a $0{,}71$.}
\comprobacion la precisión depende de la energía de la excitación frente al ruido: reducir la amplitud a la décima parte (con la mitad de energía por la forma sinusoidal) multiplica la desviación por $\sqrt{200}=14{,}1$, y $0{,}05\cdot14{,}1=0{,}71$.
\end{solucion}

\begin{solucion}{eje:2-hammerstein}
\datos $y_{k+1}=ay_k+b(u_k+cu_k^{2})$; $\hat\theta=(0{,}85;\ 0{,}30;\ -0{,}06)$.
\planteamiento \cref{ej:hammerstein}.
\pasos
\begin{enumerate}[label=\arabic*.]
  \item (a) $y_{k+1}=ay_k+bu_k+(bc)u_k^{2}$: regresores $(y_k,u_k,u_k^{2})$ y parámetros $(a,b,bc)$, lineal en ellos.
  \item (b) $a=0{,}85$; $b=0{,}30$; $c=-0{,}06/0{,}30=-0{,}2$.
  \item (c) Con $f(u)=\alpha_1u+\alpha_2u^{2}$ libre, $(b,\alpha_1,\alpha_2)$ y $(b/\gamma,\gamma\alpha_1,\gamma\alpha_2)$ dan el mismo producto $bf$ para todo $\gamma\ne0$: solo son identificables $b\alpha_1$ y $b\alpha_2$. Fijar $\alpha_1=1$ elimina la ambigüedad.
\end{enumerate}
\resultado $\mathbf{a=0{,}85}$, $\mathbf{b=0{,}30}$, $\mathbf{c=-0{,}2}$; sin normalización, la identificación es única salvo un factor de escala.
\comprobacion $b\cdot c=0{,}30\cdot(-0{,}2)=-0{,}06$, el tercer parámetro estimado.
\end{solucion}

\begin{solucion}{eje:2-cota}
\datos $L=0{,}98$; $\varepsilon=0{,}01$ K.
\planteamiento \cref{prop:acumula}: $d_k\le\varepsilon(1-L^{k})/(1-L)$.
\pasos
\begin{enumerate}[label=\arabic*.]
  \item $k=4$: $0{,}01(1-0{,}9224)/0{,}02=0{,}0388$ K. $k=24$: $0{,}01(1-0{,}6158)/0{,}02=0{,}1921$ K. $k=96$: $0{,}01(1-0{,}1438)/0{,}02=0{,}4281$ K. Límite: $0{,}5$ K.
  \item (b) $\varepsilon\le0{,}2\cdot0{,}02/(1-0{,}98^{96})=0{,}004/0{,}8562=0{,}00467$ K.
  \item (c) Margen dependiente del paso, $m_j=\varepsilon(1-L^{j})/(1-L)$, que vale cero al principio y crece con el horizonte; o, más sencillo y conservador, un margen constante igual a la cota en el último paso.
\end{enumerate}
\resultado \textbf{$0{,}039$; $0{,}192$; $0{,}428$ y $0{,}5$ K; $\varepsilon\le0{,}0047$ K por paso.}
\comprobacion la cota a un paso es $\varepsilon=0{,}01$ y crece monótonamente hacia $\varepsilon/(1-L)=0{,}5$.
\end{solucion}

\begin{solucion}{eje:2-lazo}
\datos $a=0{,}9$; $b=0{,}5$; $K=0{,}4$.
\planteamiento \cref{prop:lazo}.
\pasos
\begin{enumerate}[label=\arabic*.]
  \item (a) $y_{k+1}=(0{,}9-0{,}5\cdot0{,}4)y_k+\varepsilon_{k+1}=0{,}7y_k+\varepsilon_{k+1}$.
  \item (b) Cualquier par con $a'-0{,}4b'=0{,}7$: $(0{,}7;\ 0)$, $(0{,}9;\ 0{,}5)$, $(1{,}1;\ 1)$.
  \item (c) Con $u_k=-0{,}4y_k+v_k$: $\frac1N\sum y_k^{2}\to\sigma_y^{2}$, $\frac1N\sum y_ku_k\to-0{,}4\sigma_y^{2}$ (porque $y_k$ solo depende de $v$ hasta $k-1$ y $v_k$ es independiente de $y_k$) y $\frac1N\sum u_k^{2}\to0{,}16\sigma_y^{2}+\sigma_v^{2}$. El determinante del límite es $\sigma_y^{2}(0{,}16\sigma_y^{2}+\sigma_v^{2})-0{,}16\sigma_y^{4}=\sigma_y^{2}\sigma_v^{2}>0$.
\end{enumerate}
\resultado \textbf{se observa $y_{k+1}=0{,}7y_k$; los datos no distinguen $(a,b)$; con excitación externa la matriz límite es invertible.}
\comprobacion el par verdadero $(0{,}9;\ 0{,}5)$ está entre los del apartado (b), y el determinante se anula si $\sigma_v=0$, que es el caso de la \cref{prop:lazo}.
\end{solucion}

\begin{solucion}{eje:2-sesgo}
\datos $x_{k+1}=ax_k+bu_k$; $y_k=x_k+v_k$; $u$, $v$ blancos e independientes.
\planteamiento límite del estimador de mínimos cuadrados: $\theta^{*}=R_{\varphi\varphi}^{-1}r_{\varphi y}$.
\pasos
\begin{enumerate}[label=\arabic*.]
  \item Varianza estacionaria: $\sigma_x^{2}=a^{2}\sigma_x^{2}+b^{2}\sigma_u^{2}$, luego $\sigma_x^{2}=b^{2}\sigma_u^{2}/(1-a^{2})$.
  \item $\E[y_k^{2}]=\sigma_x^{2}+\sigma_v^{2}$; $\E[y_ku_k]=0$, porque $x_k$ depende de $u$ hasta $k-1$; $\E[u_k^{2}]=\sigma_u^{2}$. La matriz $R_{\varphi\varphi}$ es diagonal.
  \item $\E[y_ky_{k+1}]=\E[(x_k+v_k)(ax_k+bu_k+v_{k+1})]=a\sigma_x^{2}$; $\E[u_ky_{k+1}]=b\sigma_u^{2}$.
  \item Por tanto $\hat\alpha\to a\sigma_x^{2}/(\sigma_x^{2}+\sigma_v^{2})$ y $\hat\beta\to b$.
  \item Con $a=0{,}9$, $b=\sigma_u=\sigma_v=1$: $\sigma_x^{2}=1/0{,}19=5{,}263$; $\hat\alpha\to0{,}9\cdot5{,}263/6{,}263=0{,}7563$.
  \item Constantes de tiempo, en muestras: la verdadera es $-1/\ln0{,}9=9{,}49$ y la estimada, $-1/\ln0{,}7563=3{,}58$.
\end{enumerate}
\resultado \textbf{$\hat\alpha\to a\,\sigma_x^{2}/(\sigma_x^{2}+\sigma_v^{2})$: atenuación hacia cero; el modelo es casi tres veces más rápido que la planta.}
\comprobacion si $\sigma_v=0$, $\hat\alpha\to a$; el ruido en un regresor atenúa su coeficiente, el fenómeno de error en las variables, y es la razón de preferir la validación en simulación libre y los métodos de error de salida.
\end{solucion}

\begin{solucion}{eje:2-diseno}
\planteamiento aplicar el tema completo: excitación (\cref{def:pe}), lazo cerrado (\cref{prop:lazo}), partición (\cref{obs:particion}), validación (\cref{def:validacion}), gemelo (\cref{def:gemelo}).
\pasos
\begin{enumerate}[label=\arabic*.]
  \item Históricos del PID: útiles para la parte estacional y las perturbaciones, pero poco informativos sobre la respuesta a la entrada (\cref{prop:lazo}). Se usan como prueba final por estación, no para estimar.
  \item Experimentos: dos semanas con escalones aleatorios de consigna o de potencia, de 1 a 4~h, dentro de una banda segura (por ejemplo, $\pm2$ K respecto a la consigna habitual y solo fuera del horario de ocupación), con parada automática si se sale de la banda de confort ampliada.
  \item Partición: semana~1 para estimar, semana~2 para validar; los históricos de otra estación, como prueba.
  \item Modelos: ARX y NARX con regresores físicos (producto potencia por temperatura exterior, radiación), modelo RC de dos nodos ajustado por mínimos cuadrados no lineales y, si hay datos, un modelo de estado neuronal.
  \item Métricas: RMSE y error máximo en simulación libre a 24~h, sesgo medio y error en la energía diaria.
  \item Criterio de aceptación: por ejemplo, RMSE en simulación a 24~h menor que $0{,}3$ K y error de energía diaria menor que el 5~\%; el margen del MPC se fija con el error observado (\cref{eje:2-cota}).
  \item Recalibración: vigilar la media diaria del residuo (\cref{eje:6-deriva}) y reestimar con ventana móvil cuando supere el umbral, validando el modelo nuevo en sombra antes de sustituir al anterior.
\end{enumerate}
\resultado \textbf{experimento acotado y excitante en lazo abierto, partición temporal, validación en simulación libre con umbrales, y recalibración disparada por el seguimiento de residuos.}
\comprobacion cada paso responde a uno de los errores frecuentes de la \cref{obs:errores2}.
\end{solucion}

% ---------------------------------------------------------------------
\section{Tema 3}

\begin{solucion}{eje:3-deslizante}
\planteamiento \cref{def:mpc}.
\pasos
\begin{enumerate}[label=\arabic*.]
  \item El plan es óptimo para la predicción hecha con la información de $k$; en $k+1$ hay una medida nueva que incorpora perturbaciones y errores de modelo. Volver a resolver es lo que convierte el plan en realimentación.
  \item En lazo abierto, un error de $2$ K en la temperatura exterior desplaza la temperatura interior durante horas (\cref{prop:acumula}) y el plan, calculado para otra trayectoria, provocaría violaciones de confort o calentamiento innecesario sin corregirlos.
\end{enumerate}
\resultado \textbf{se aplica solo la primera decisión para incorporar cada nueva medida; en lazo abierto los errores de previsión no se corrigen.}
\comprobacion el MPC del cuaderno mantiene el confort pese al ruido de medida y al error del modelo identificado, precisamente porque reoptimiza en cada paso.
\end{solucion}

\begin{solucion}{eje:3-lqr}
\datos $A=2$; $B=1$; $Q=R=1$; $P_f=0$; $N=2$; $x_0=1$.
\planteamiento recursión del \cref{teo:lqr}.
\pasos
\begin{enumerate}[label=\arabic*.]
  \item $P_2=0$; $K_1=(1+0)^{-1}\cdot0=0$; $P_1=1+4\cdot0-0=1$.
  \item $K_0=(1+1)^{-1}\cdot1\cdot2=1$; $P_0=1+4\cdot1-2\cdot1\cdot1=3$.
  \item Secuencia: $u_0=-x_0=-1$, $x_1=2-1=1$, $u_1=0$, $x_2=2$. Coste: $1+1+1+0=3=P_0x_0^{2}$.
  \item Horizonte infinito: $P=1+4P-4P^{2}/(1+P)$, es decir, $P^{2}-4P-1=0$, $P=2+\sqrt5=4{,}236$, $K=2P/(1+P)=1{,}618$, y el lazo cerrado tiene autovalor $2-1{,}618=0{,}382$.
\end{enumerate}
\resultado \textbf{$P_1=1$, $K_1=0$, $P_0=3$, $K_0=1$; coste $3$; $P=4{,}236$, $K=1{,}618$, autovalor $0{,}382$.}
\comprobacion el sistema abierto es inestable ($\abs{2}>1$) y el LQR de horizonte infinito lo estabiliza ($0{,}382<1$); con horizonte~2 y $P_f=0$, el último paso no actúa ($K_1=0$) porque el estado final no cuesta nada.
\end{solucion}

\begin{solucion}{eje:3-condensada}
\datos $a=0{,}9$; $b=0{,}5$; $N=3$; $Q=P_f=1$; $R=0{,}1$.
\planteamiento \eqref{eq:condensada} y \cref{prop:qp}.
\pasos
\begin{enumerate}[label=\arabic*.]
  \item $\Phi=(0{,}9;\ 0{,}81;\ 0{,}729)\T$; $\Gamma=\begin{pmatrix}0{,}5&0&0\\0{,}45&0{,}5&0\\0{,}405&0{,}45&0{,}5\end{pmatrix}$.
  \item $\Gamma\T\Gamma$: diagonal $0{,}25+0{,}2025+0{,}1640=0{,}6165$; $0{,}25+0{,}2025=0{,}4525$; $0{,}25$; fuera de la diagonal $0{,}225+0{,}18225=0{,}40725$, $0{,}2025$ y $0{,}225$.
  \item $H=\Gamma\T\Gamma+0{,}1I=\begin{pmatrix}0{,}7165&0{,}4073&0{,}2025\\0{,}4073&0{,}5525&0{,}225\\0{,}2025&0{,}225&0{,}35\end{pmatrix}$.
  \item Autovalores: $0{,}185$; $0{,}273$; $1{,}161$, todos positivos.
\end{enumerate}
\resultado \textbf{$H\succ0$, con autovalores entre $0{,}185$ y $1{,}161$.}
\comprobacion el menor autovalor es mayor que $R=0{,}1$, como garantiza $v\T Hv\ge v\T\bar Rv$.
\end{solucion}

\begin{solucion}{eje:3-bateria}
\datos precios $(0{,}10;\ 0{,}30;\ 0{,}05;\ 0{,}25)$; $\eta=0{,}9$; capacidad y potencia de $1$.
\planteamiento energía que vuelve a la red: $\eta^{2}$ por kWh comprado.
\pasos
\begin{enumerate}[label=\arabic*.]
  \item (a) Comprar $1$ kWh cuesta $p_c$ y devuelve $\eta^{2}$ kWh, vendidos a $p_v$: beneficio $p_v\eta^{2}-p_c>0\iff p_v/p_c>1/\eta^{2}$.
  \item (b) Ciclo~1: $0{,}30\cdot0{,}81-0{,}10=0{,}143$; ciclo~2: $0{,}25\cdot0{,}81-0{,}05=0{,}1525$; total $0{,}2955$ €. Con $(0{,}20;\ 0{,}24)$: $0{,}24/0{,}20=1{,}2<1{,}235$ y $0{,}24\cdot0{,}81-0{,}20=-0{,}0056$: el segundo ciclo no se hace y el beneficio es $0{,}143$ €.
  \item (c) Cada ciclo mueve $1+0{,}81=1{,}81$ kWh por la red y cuesta $0{,}0362$ € de degradación: $0{,}1068+0{,}1163=0{,}2231$ €.
\end{enumerate}
\resultado \textbf{umbral $1/\eta^{2}=1{,}235$; $0{,}2955$ €; $0{,}143$ € con el segundo par cambiado; $0{,}2231$ € con degradación.}
\comprobacion con $\eta=1$ y sin degradación se recupera $0{,}40$ € del \cref{ej:bateria-lp}.
\end{solucion}

\begin{solucion}{eje:3-precalentar}
\datos $\tau=15$ h; 4~h entre las 6 y las 10; precios $0{,}10$ y $0{,}22$ €/kWh.
\planteamiento el exceso de temperatura decae como $e^{-t/\tau}$ (\cref{ej:zoh}); precio efectivo $=$ precio $/$ fracción retenida.
\pasos
\begin{enumerate}[label=\arabic*.]
  \item (a) $e^{-4/15}=0{,}766$.
  \item (b) $0{,}10/0{,}766=0{,}131<0{,}22$ €/kWh: conviene.
  \item (c) Con $\tau=5$ h: $e^{-0{,}8}=0{,}449$ y $0{,}10/0{,}449=0{,}223>0{,}22$: ya no conviene (por poco).
  \item Con COP: el calor cuesta $0{,}10/2{,}4=0{,}0417$ €/kWh$_{\mathrm{t}}$ a las 6 y $0{,}22/2{,}9=0{,}0759$ a las 10. Con $\tau=15$ h el efectivo es $0{,}0417/0{,}766=0{,}0544<0{,}0759$: conviene; con $\tau=5$ h, $0{,}0417/0{,}449=0{,}0927>0{,}0759$: no conviene.
\end{enumerate}
\resultado \textbf{se retiene el 77~\%; precalentar conviene en el edificio pesado y no en el ligero.}
\comprobacion es lo que hace el MPC del \cref{ej:caso-mpc}: almacenar calor en la masa térmica solo compensa si la constante de tiempo es larga frente al desfase entre valle y punta.
\end{solucion}

\begin{solucion}{eje:3-blanda}
\datos $J=3+4u_0+2u_1+3u_0^{2}+2u_0u_1+2u_1^{2}$; $u_0\ge-0{,}5-s$, $s\ge0$; multiplicador de la restricción dura $0{,}5$.
\planteamiento con la holgura activa, la condición de estacionariedad en $u_0$ es $\partial J/\partial u_0=\rho$.
\pasos
\begin{enumerate}[label=\arabic*.]
  \item $\rho=0{,}2$: $4+6u_0+2u_1=0{,}2$ y $2+2u_0+4u_1=0$. De la segunda, $u_1=-(1+u_0)/2$; sustituyendo, $5u_0=-2{,}8$, $u_0=-0{,}56$, $u_1=-0{,}22$, $s=0{,}06$.
  \item Coste: $J=1{,}604$ y $J+\rho s=1{,}616$.
  \item $\rho=1\ge0{,}5$: la holgura no compensa; $s=0$, $u=(-0{,}5;\ -0{,}25)$ y $J=1{,}625$, la solución dura.
\end{enumerate}
\resultado \textbf{con $\rho=0{,}2$ la restricción se viola en $0{,}06$; con $\rho=1$ se recupera exactamente la solución dura.}
\comprobacion \cref{prop:exacta}: el umbral es el multiplicador $0{,}5$; resolviendo con CVXPY se obtienen los mismos valores.
\end{solucion}

\begin{solucion}{eje:3-riccati}
\planteamiento inducción sobre la recursión y comparación de costes de horizontes distintos.
\pasos
\begin{enumerate}[label=\arabic*.]
  \item Simetría: $P_N=P_f$ es simétrica. Si $P_{k+1}$ lo es, $A\T P_{k+1}A$ es simétrica y $A\T P_{k+1}BK_k=A\T P_{k+1}BM^{-1}B\T P_{k+1}A$ también, porque $M=R+B\T P_{k+1}B$ es simétrica; luego $P_k$ lo es.
  \item Semidefinición: $x\T P_kx=V_k(x)$ es el mínimo de una suma de términos $\ge0$ (\cref{teo:lqr}).
  \item Monotonía con $P_f=0$: $V_{k-1}(x)$ es el coste óptimo con un paso más que $V_k(x)$. Dada la secuencia óptima de $V_{k-1}$, sus primeros $N-k$ pasos forman una secuencia admisible para el horizonte menor, con coste no mayor (se suprimen términos no negativos); luego $V_k(x)\le V_{k-1}(x)$.
\end{enumerate}
\resultado \textbf{$P_k=P_k\T\succeq0$ y, con $P_f=0$, $x\T P_kx$ crece al aumentar el horizonte restante.}
\comprobacion en el \cref{eje:3-lqr}, $P_2=0\le P_1=1\le P_0=3$, que convergen hacia $P=4{,}236$ al alargar el horizonte.
\end{solucion}

\begin{solucion}{eje:3-terminal}
\datos $x^{+}=x+u$; $\ell=x^{2}+u^{2}$; $\abs u\le0{,}5$; $P=1{,}618$; $K=0{,}618$.
\planteamiento comprobar las hipótesis del \cref{teo:mpc}.
\pasos
\begin{enumerate}[label=\arabic*.]
  \item (a) Con $u=-Kx$, $x^{+}=(1-K)x=0{,}382x$. $V_f(x^{+})-V_f(x)=1{,}618(0{,}1459-1)x^{2}=-1{,}382x^{2}$ y $\ell=(1+0{,}382)x^{2}=1{,}382x^{2}$: se cumple con igualdad, por ser $P$ solución de Riccati.
  \item (b) $\abs{Kx}\le0{,}5\iff\abs x\le0{,}5/0{,}618=0{,}809$. Si $\abs x\le0{,}809$, $\abs{x^{+}}=0{,}382\abs x\le0{,}809$: invariante. $\Xset_f=[-0{,}809;\ 0{,}809]$.
  \item (c) $f(0,0)=0$, $\ell\ge x^{2}$, y $V_f$, $\kappa_f$ y $\Xset_f$ cumplen las hipótesis; por el \cref{teo:mpc}, para todo $N\ge1$ y todo $x_0$ desde el que $\mathbb P_N$ sea factible (los que pueden llegar a $\Xset_f$ en $N$ pasos, $\abs{x_0}\le0{,}809+0{,}5N$), el lazo cerrado es factible siempre y $x_k\to0$.
\end{enumerate}
\resultado \textbf{la condición de decrecimiento se cumple con igualdad; $\Xset_f=[-0{,}809;\ 0{,}809]$; estabilidad para todo $N$ en la región factible.}
\comprobacion $1{,}618\cdot0{,}382^{2}+1+0{,}618^{2}=0{,}236+1+0{,}382=1{,}618=P$: la ecuación de Riccati escalar.
\end{solucion}

% ---------------------------------------------------------------------
\section{Tema 4}

\begin{solucion}{eje:4-truncado}
\datos $r\approx-0{,}05$ €; $\gamma=0{,}99$; truncación a los 288 pasos.
\planteamiento valor estacionario $r/(1-\gamma)$; objetivo $y=r+\gamma(1-d)V(s')$.
\pasos
\begin{enumerate}[label=\arabic*.]
  \item (a) $V\approx-0{,}05/0{,}01=-5$ €.
  \item (b) Tratada como terminación: $y=-0{,}05$. Correctamente ($d=0$): $y=-0{,}05+0{,}99\cdot(-5)=-5{,}0$.
  \item (c) Los estados cercanos al corte parecen cien veces mejores de lo que son. Si la observación permite reconocerlos (por ejemplo, por la hora), el agente aprende que al final no importa lo que haga y deja de calentar; si no puede reconocerlos, el objetivo es ruidoso e inconsistente.
\end{enumerate}
\resultado \textbf{$V\approx-5$ €; objetivo $-0{,}05$ (incorrecto) frente a $-5{,}0$ (correcto).}
\comprobacion $-0{,}05+0{,}99\cdot(-5)=-5$, consistente con un valor estacionario: el estado siguiente vale lo mismo que el actual.
\end{solucion}

\begin{solucion}{eje:4-constante}
\planteamiento sumar la constante dentro del retorno.
\pasos
\begin{enumerate}[label=\arabic*.]
  \item (a) $G_t'=\sum_k\gamma^{k}(R_{t+k+1}+c)=G_t+c/(1-\gamma)$. Tomando esperanzas, $V'^{\pi}=V^{\pi}+c/(1-\gamma)$ y $Q'^{\pi}=Q^{\pi}+c/(1-\gamma)$ para toda política; el desplazamiento es igual para todas las acciones y no cambia el $\arg\max$.
  \item (b) Desde $s$, «terminar» da recompensa $0$ y acaba; «seguir» da $-0{,}5$ y en el paso siguiente acaba con recompensa~$0$. Con $\gamma=1$: terminar vale $0$ y seguir $-0{,}5$; es mejor terminar. Con $c=1$: terminar vale $1$ y seguir $0{,}5+1=1{,}5$; ahora es mejor seguir.
\end{enumerate}
\resultado \textbf{en tareas continuadas, $V'=V+c/(1-\gamma)$ y la política no cambia; en episódicas, la constante premia o castiga la duración y puede cambiarla.}
\comprobacion en (b) el retorno con $c$ suma $c$ multiplicado por el número de pasos, que depende de la política: es exactamente lo que no ocurre en la tarea continuada.
\end{solucion}

\begin{solucion}{eje:4-sobre}
\datos $Q=(0;\ -0{,}5)$; $\sigma=1$.
\planteamiento $\E[\max(X,Y)]$ para gaussianas independientes: con $\vartheta=\sqrt{\sigma_X^{2}+\sigma_Y^{2}}$ y $\alpha=(\mu_X-\mu_Y)/\vartheta$, vale $\mu_X\Phi(\alpha)+\mu_Y\Phi(-\alpha)+\vartheta\varphi(\alpha)$.
\pasos
\begin{enumerate}[label=\arabic*.]
  \item (a) $\vartheta=\sqrt2=1{,}414$; $\alpha=0{,}354$; $\Phi(-\alpha)=0{,}362$; $\varphi(\alpha)=0{,}375$. $\E[\max]=0-0{,}5\cdot0{,}362+1{,}414\cdot0{,}375=0{,}349>0$.
  \item (b) $\E[\min(\hat Q_1,\hat Q_2)]=-\sigma/\sqrt\pi=-0{,}564$, por simetría con la \cref{prop:sobre}.
  \item (c) Con un solo crítico, los valores $Q$ estimados tienden a superar a los retornos observados y el error crece por el \emph{bootstrapping}; con el mínimo de dos críticos, el sesgo pasa a ser negativo y los valores $Q$ quedan por debajo o cerca de los retornos, lo que se ve al comparar las curvas de $Q$ y de retorno de evaluación en la tarea.
\end{enumerate}
\resultado \textbf{sesgo $+0{,}349$ con el máximo; $-0{,}564$ con el mínimo de dos estimaciones.}
\comprobacion si las dos acciones valieran lo mismo, la fórmula daría $\vartheta\varphi(0)=1{,}414\cdot0{,}399=0{,}564=1/\sqrt\pi$, el caso de la \cref{prop:sobre}.
\end{solucion}

\begin{solucion}{eje:4-blanda}
\datos $Q=(1,2,0)$; $\alpha\in\{1;\ 0{,}1\}$.
\planteamiento \cref{prop:blanda}.
\pasos
\begin{enumerate}[label=\arabic*.]
  \item $\alpha=1$: $Z=e+e^{2}+1=11{,}107$; $\pi^{*}=(0{,}2447;\ 0{,}6652;\ 0{,}0900)$; $V=\ln11{,}107=2{,}4076$.
  \item $\E[Q]=0{,}2447+1{,}3305=1{,}5752$; $\mathcal H=0{,}8324$ nats; suma $2{,}4076$.
  \item $\alpha=0{,}1$: $\pi^{*}=(4{,}5\cdot10^{-5};\ 0{,}99995;\ 2\cdot10^{-9})$; $V=0{,}1\ln(e^{10}+e^{20}+1)=2{,}0000045$; $\mathcal H=0{,}0005$ nats.
\end{enumerate}
\resultado \textbf{$V=2{,}408$ con $\alpha=1$ y $V=2{,}000005$ con $\alpha=0{,}1$.}
\comprobacion al bajar la temperatura, la política se concentra en la acción de mayor $Q$ y $V\to\max Q=2$.
\end{solucion}

\begin{solucion}{eje:4-temperatura}
\datos $\mathcal H=-1{,}5$; $\bar{\mathcal H}=-1$; $\alpha=0{,}2$; tasa $0{,}1$; $\sigma=0{,}5$; $u=0{,}5$.
\planteamiento $\partial J/\partial\alpha=\E[-\log\pi]-\bar{\mathcal H}=\mathcal H-\bar{\mathcal H}$; entropía gaussiana $\frac12\ln(2\pi e\sigma^{2})$; corrección $-\log(1-\tanh^{2}u)$.
\pasos
\begin{enumerate}[label=\arabic*.]
  \item (a) $\partial J/\partial\alpha=-1{,}5+1=-0{,}5$; $\alpha\leftarrow0{,}2-0{,}1\cdot(-0{,}5)=0{,}25$.
  \item (b) $\mathcal H=\frac12\ln(2\pi e\cdot0{,}25)=\frac12\ln4{,}270=0{,}7258$ nats.
  \item $\tanh0{,}5=0{,}4621$; $\ln(1-0{,}4621^{2})=\ln0{,}7864=-0{,}2402$; luego $\log\pi(a)=\log\mathcal N(u)+0{,}2402$.
\end{enumerate}
\resultado \textbf{$\alpha$ sube a $0{,}25$; $\mathcal H=0{,}726$ nats; la compresión suma $0{,}240$ al logaritmo de la densidad.}
\comprobacion la entropía estaba por debajo del objetivo y la temperatura sube para premiar más la exploración; la densidad aumenta tras $\tanh$ porque la función comprime el intervalo.
\end{solucion}

\begin{solucion}{eje:4-lagrange}
\datos $(J_r,J_c)$: $\pi_1=(-18;\ 2)$, $\pi_2=(-19;\ 0{,}5)$, $\pi_3=(-21;\ 0)$; $d=1$; $\eta=0{,}5$.
\planteamiento $L_i(\lambda)=J_r-\lambda(J_c-d)$; ascenso dual $\lambda\leftarrow\max\{0,\lambda+\eta(J_c-d)\}$.
\pasos
\begin{enumerate}[label=\arabic*.]
  \item (a) Factibles: $\pi_2$ y $\pi_3$; la mejor es $\pi_2$ ($-19$). Si se admiten mezclas, $\frac13\pi_1+\frac23\pi_2$ tiene $J_c=1$ y $J_r=-18{,}67$: el óptimo de un CMDP puede ser aleatorizado~\cite{altman}.
  \item (b) $L_1=-18-\lambda$, $L_2=-19+0{,}5\lambda$, $L_3=-21+\lambda$. $L_1\ge L_2\iff\lambda\le2/3$; $L_2\ge L_3\iff\lambda\le4$. Gana $\pi_1$ si $\lambda<2/3$, $\pi_2$ si $2/3<\lambda<4$ y $\pi_3$ si $\lambda>4$.
  \item (c) $\lambda_0=0\to\pi_1$, $\lambda_1=0{,}5\to\pi_1$, $\lambda_2=1\to\pi_2$, $\lambda_3=0{,}75\to\pi_2$, $\lambda_4=0{,}5\to\pi_1$: oscila entre $\pi_1$ y $\pi_2$.
\end{enumerate}
\resultado \textbf{óptima determinista $\pi_2$; el ascenso dual oscila entre $\pi_1$ y $\pi_2$ y, en promedio, se comporta como una mezcla de ambas.}
\comprobacion la mezcla del apartado (a) tiene la restricción activa, $2/3+1/3=1$, y su valor $-18{,}67$ queda entre los de $\pi_1$ y $\pi_2$, como corresponde a un multiplicador en $[2/3,4]$.
\end{solucion}

\begin{solucion}{eje:4-shaping}
\datos $\Phi(s)=-\kappa\max\{\underline T-T,0\}$; $\kappa=1$; $\gamma=0{,}99$; $T=19\to19{,}5$ °C; $\underline T=20$ °C.
\planteamiento repetir la suma telescópica del \cref{teo:shaping} con horizonte aleatorio $T_f$.
\pasos
\begin{enumerate}[label=\arabic*.]
  \item (a) Si el episodio termina en $s_{T_f}$ terminal, $\sum_{t<T_f}\gamma^{t}F_t=\gamma^{T_f}\Phi(s_{T_f})-\Phi(s_0)=-\Phi(s_0)$, porque $\Phi$ es nula en los terminales. Como antes, $Q'=Q-\Phi(s)$ y el $\arg\max$ no cambia.
  \item (b) $\Phi(s)=-1$, $\Phi(s')=-0{,}5$; $F=0{,}99\cdot(-0{,}5)-(-1)=0{,}505>0$: se premia acercarse a la banda.
  \item (c) Una bonificación $+0{,}1$ cada vez que $T$ sube, sin penalización cuando baja: el agente puede calentar y dejar enfriar en ciclo para cobrarla indefinidamente. No deriva de un potencial: su suma sobre un ciclo cerrado no es nula.
\end{enumerate}
\resultado \textbf{la política óptima se conserva con terminales si $\Phi=0$ en ellos; $F=0{,}505$; las bonificaciones no potenciales crean ciclos explotables.}
\comprobacion sobre cualquier ciclo $s_0\to\dots\to s_0$, con $\gamma=1$ la suma de $F$ de un potencial es nula, y la de la bonificación de (c) es positiva.
\end{solucion}

\begin{solucion}{eje:4-simlema}
\datos $r_{\max}=1$; $\epsilon=0{,}01$.
\planteamiento \cref{teo:simlema}; descomposición de la pérdida de rendimiento en tres términos.
\pasos
\begin{enumerate}[label=\arabic*.]
  \item (a) $\gamma=0{,}9$: $0{,}9\cdot0{,}01/0{,}01=0{,}9$, frente a un rango de $V$ de anchura $2/(1-\gamma)=20$ (4,5~\%). $\gamma=0{,}99$: $0{,}99\cdot0{,}01/10^{-4}=99$, frente a 200 (50~\%): casi inútil.
  \item (b) $V^{\pi^{*}}-V^{\hat\pi^{*}}=(V^{\pi^{*}}-\hat V^{\pi^{*}})+(\hat V^{\pi^{*}}-\hat V^{\hat\pi^{*}})+(\hat V^{\hat\pi^{*}}-V^{\hat\pi^{*}})$. El término central es $\le0$ porque $\hat\pi^{*}$ es óptima en el simulador; cada uno de los otros está acotado por el lema. Pérdida $\le2\gamma\epsilon r_{\max}/(1-\gamma)^{2}$.
  \item (c) El error crece como $(1-\gamma)^{-2}$: con un modelo imperfecto conviene un horizonte efectivo más corto o, mejor, replanificar con realimentación frecuente como hace el MPC, que reinicia la predicción desde el estado medido.
\end{enumerate}
\resultado \textbf{cota de $0{,}9$ con $\gamma=0{,}9$ y de $99$ con $\gamma=0{,}99$; pérdida de la política del simulador $\le2\gamma\epsilon r_{\max}/(1-\gamma)^{2}$.}
\comprobacion con $\epsilon=0$ la cota es nula, como debe: modelos idénticos dan valores idénticos.
\end{solucion}

% ---------------------------------------------------------------------
\section{Tema 5}

\begin{solucion}{eje:5-kpi}
\planteamiento protocolo de la sección~5.1 y errores de la \cref{obs:errores5}.
\pasos
\begin{enumerate}[label=\arabic*.]
  \item ¿Respecto a qué controlador y con qué ajuste?
  \item ¿Mismo periodo, o antes-después con tiempo meteorológico y uso distintos? ¿Se normalizó por grados-día y ocupación?
  \item ¿Se mantuvo el confort? ¿Con qué indicador?
  \item ¿Energía o coste? ¿De qué consumo: climatización o edificio entero?
  \item ¿Cuántos días, con qué dispersión e intervalo de confianza?
  \item ¿Tenían ambos controladores la misma información (previsiones, sensores)?
\end{enumerate}
\resultado \textbf{sin referencia, protocolo, confort e incertidumbre, un porcentaje de ahorro no es interpretable.}
\comprobacion cada pregunta corresponde a uno de los cinco requisitos del protocolo.
\end{solucion}

\begin{solucion}{eje:5-ahorro}
\datos \cref{tab:hvac}; $\rho=2$ €/(K·h).
\planteamiento ahorro $=1-\text{valor}/\text{valor del PI}$.
\pasos
\begin{enumerate}[label=\arabic*.]
  \item (a) Energía: MPC $1-171{,}1/175{,}4=2{,}5$~\%; RL $1-169{,}9/175{,}4=3{,}1$~\%. Coste: MPC $1-19{,}28/21{,}06=8{,}5$~\%; RL $1-20{,}63/21{,}06=2{,}0$~\% (la tabla da $8{,}4$ y $2{,}1$~\% porque el cuaderno los calcula con los valores sin redondear).
  \item (b) Total del MPC sin margen: $18{,}76+2\cdot3{,}54=25{,}84$ €, peor que $19{,}28$ €. Sería preferible si $18{,}76+3{,}54\rho<19{,}28$, es decir, $\rho<0{,}147$ €/(K·h).
  \item (c) El MPC desplaza el consumo de punta a valle usando la masa térmica: la energía apenas cambia, el precio medio pagado sí.
\end{enumerate}
\resultado \textbf{MPC: $-2{,}5$~\% de energía y $-8{,}4$~\% de coste; RL: $-3{,}1$~\% y $-2{,}0$~\%; el MPC sin margen solo compensaría con $\rho<0{,}15$ €/(K·h).}
\comprobacion precio medio pagado: PI $21{,}06/175{,}4=0{,}120$ €/kWh; MPC $19{,}28/171{,}1=0{,}113$ €/kWh.
\end{solucion}

\begin{solucion}{eje:5-crn}
\datos $\sigma=2$ €; $r=0{,}9$.
\planteamiento \cref{prop:crn}.
\pasos
\begin{enumerate}[label=\arabic*.]
  \item (a) Independientes: $\sqrt{2\cdot4}=2{,}83$ €; comunes: $\sqrt{2\cdot4\cdot0{,}1}=0{,}894$ €.
  \item (b) Con $n$ pares independientes, $2{,}83/\sqrt n=0{,}894\Rightarrow n=8/0{,}8=10$.
\end{enumerate}
\resultado \textbf{$2{,}83$ € frente a $0{,}89$ €; una semana común equivale a diez independientes.}
\comprobacion el cociente de varianzas es $1/(1-r)=10$.
\end{solucion}

\begin{solucion}{eje:5-semenov}
\datos los del \cref{ej:caso-reactor}; $Q_g(350)=20$ kW.
\planteamiento \cref{prop:semenov}: $Q_g(T)=20\exp\bigl(10^{4}(1/350-1/T)\bigr)$ y $Q_g'=Q_gE/(RT^{2})$.
\pasos
\begin{enumerate}[label=\arabic*.]
  \item (a) $10^{4}(1/350-1/360)=0{,}7937$; $Q_g(360)=20e^{0{,}7937}=44{,}23$ kW; $Q_g'=44{,}23\cdot10^{4}/360^{2}=3{,}413$ kW/K. Haría falta $UA>3{,}41$ kW/K, con $T_j=360-44{,}23/UA$.
  \item (b) $UA=1{,}5$: $T_j=350-20/1{,}5=336{,}67$ K, pero $Q_g'(350)=1{,}633>1{,}5$: el equilibrio es inestable.
  \item (c) Bajar la temperatura de operación hasta donde $Q_g'<1{,}5$ kW/K, que ocurre por debajo de $348{,}9$ K (con $T_j=336{,}7$ K), con margen adicional; dosificar el reactivo (semilotes) para reducir $C_A$; y programar la limpieza del intercambiador.
\end{enumerate}
\resultado \textbf{$Q_g'(360)=3{,}41$ kW/K; con $UA=1{,}5$ kW/K el punto de $350$ K se vuelve inestable.}
\comprobacion la pendiente crece con $T$ como $Q_g/T^{2}$, exponencialmente, mientras que la retirada es lineal: subir $T$ siempre acerca al límite de estabilidad.
\end{solucion}

\begin{solucion}{eje:5-pue}
\datos reducción del 40~\% en refrigeración y del 15~\% en el exceso de PUE; PUE ilustrativo $1{,}12$.
\planteamiento exceso $=$ PUE $-1$; la refrigeración es una fracción $f$ del exceso.
\pasos
\begin{enumerate}[label=\arabic*.]
  \item (a) $0{,}40f=0{,}15\Rightarrow f=0{,}375$.
  \item (b) Exceso $0{,}12\to0{,}12\cdot0{,}85=0{,}102$; PUE $1{,}102$. Consumo total: $(1{,}12-1{,}102)/1{,}12=1{,}6$~\%.
  \item (c) Los tres porcentajes (40, 15 y 1,6) se refieren a bases distintas: refrigeración, exceso sobre la carga informática y consumo total.
\end{enumerate}
\resultado \textbf{la refrigeración era el 37,5~\% del exceso; PUE $1{,}102$; un 1,6~\% del consumo total.}
\comprobacion $0{,}375\cdot0{,}12\cdot0{,}40=0{,}018=1{,}12-1{,}102$.
\end{solucion}

\begin{solucion}{eje:5-bateria}
\datos previsión $(0{,}10;\ 0{,}30;\ 0{,}05;\ 0{,}25)$; precios reales $(0{,}10;\ 0{,}30;\ 0{,}28;\ 0{,}25)$; $1$ kWh, $\eta=1$.
\planteamiento plan del \cref{ej:bateria-lp}: cargar, descargar, cargar, descargar.
\pasos
\begin{enumerate}[label=\arabic*.]
  \item (a) Lazo abierto: $-0{,}10+0{,}30-0{,}28+0{,}25=0{,}17$ €.
  \item (b) MPC: hora~1, carga; hora~2, descarga a $0{,}30$; hora~3, observa $0{,}28$ y, con la previsión $0{,}25$ para la hora~4, comprar perdería $0{,}03$: espera. Beneficio $0{,}20$ €.
  \item (c) Información perfecta: $0{,}20$ € (un solo ciclo; mantener la carga hasta la hora~3 daría $0{,}18$).
\end{enumerate}
\resultado \textbf{$0{,}17$ € en lazo abierto, $0{,}20$ € con el MPC, igual al óptimo con información perfecta.}
\comprobacion la diferencia ($0{,}03$ €) es exactamente la pérdida del segundo ciclo, que el MPC evita al observar el precio real antes de decidir.
\end{solucion}

\begin{solucion}{eje:5-semenov-dem}
\planteamiento condiciones de tangencia; convexidad de $h(T)=Q_g(T)-UA(T-T_j)$.
\pasos
\begin{enumerate}[label=\arabic*.]
  \item (a) Tangencia: $Q_g(T^{*})=UA(T^{*}-T_j)$ y $Q_g'(T^{*})=UA$. Como $Q_g'=Q_gE/(RT^{2})$, la segunda da $Q_g(T^{*})=UA\,RT^{*2}/E$; sustituyendo en la primera, $T^{*}-T_j=RT^{*2}/E$.
  \item (b) $Q_g''=Q_g\bigl[(E/RT^{2})^{2}-2E/(RT^{3})\bigr]>0$ si $T<E/(2R)$ (5000~K aquí): $h$ es convexa. Si $h(T_1)=h(T_2)=0$ con $T_1<T_2$ y no hay más ceros, $h<0$ en $(T_1,T_2)$; luego $h'(T_1)\le0\le h'(T_2)$, y las desigualdades son estrictas porque, si $h'(T_1)=0$, $T_1$ sería un mínimo de la función convexa y $h\ge0$ en todas partes. Por la \cref{prop:semenov}, $T_1$ es estable y $T_2$ inestable.
  \item (c) $Q_g$ es proporcional a $C_A$, que solo disminuye durante el lote: congelarla en su valor inicial sobrestima el calor generado y su pendiente en cada temperatura, de modo que si el criterio se cumple con $C_{A,0}$ también se cumple después.
\end{enumerate}
\resultado \textbf{$T_j^{*}=T^{*}-RT^{*2}/E$; con dos equilibrios, el inferior es estable y el superior inestable; la hipótesis de $C_A$ constante es conservadora.}
\comprobacion con los datos del \cref{ej:caso-reactor}: $T^{*}=352{,}69$ K y $352{,}69-352{,}69^{2}/10^{4}=340{,}26$ K; y $Q_g'(350)=1{,}633<2<2{,}414=Q_g'(355{,}23)$.
\end{solucion}

\begin{solucion}{eje:5-protocolo}
\planteamiento \cref{prop:crn}, protocolo de la sección~5.1 y \cref{def:validar}.
\pasos
\begin{enumerate}[label=\arabic*.]
  \item Asignación: días alternos o bloques de dos días, sorteando el orden dentro de cada par (para no confundir con el ciclo semanal), durante al menos 8 semanas.
  \item Indicadores: energía y coste de climatización diarios, grados-hora fuera de banda en ocupación, potencia máxima y número de intervenciones del operador.
  \item Normalización: por grados-día de calefacción y por ocupación, o mejor, un modelo de regresión del consumo con temperatura exterior, radiación y ocupación como covariables y el controlador como factor.
  \item Análisis: contraste $t$ pareado sobre los pares de días (\cref{ej:contraste}) o el modelo de regresión; intervalo de confianza del ahorro.
  \item Parada: volver al PID si la violación diaria supera un umbral (por ejemplo, $1$ K·h en ocupación), si hay quejas o si el MPC falla en resolver; todo registrado.
\end{enumerate}
\resultado \textbf{diseño emparejado y aleatorizado, indicadores de energía, coste y confort normalizados, contraste pareado e intervalos, y criterios de parada definidos de antemano.}
\comprobacion es el método publicado de activar y desactivar el sistema que se usó en las pruebas del centro de datos~\cite{evansgao}, con un análisis estadístico explícito.
\end{solucion}

% ---------------------------------------------------------------------
\section{Tema 6}

\begin{solucion}{eje:6-modos}
\planteamiento \cref{def:modos}.
\pasos
\begin{enumerate}[label=\arabic*.]
  \item (a) Asesor: el sistema no ejecuta. (b) Supervisado: ejecuta dentro de límites con el operador vigilando. (c) Autónomo con supervisión: actúa cada cinco minutos con salvaguardas automáticas y los operadores pueden salir del modo de IA~\cite{gamblegao}.
  \item Evidencia para pasar de (a) a (b): un periodo en sombra con las recomendaciones registradas; tasa de aceptación por el operador; ninguna recomendación que hubiera violado límites; ahorro estimado con un experimento pareado; filtro y respaldo probados.
\end{enumerate}
\resultado \textbf{(a) asesor; (b) supervisado; (c) autónomo con supervisión.}
\comprobacion la secuencia coincide con la evolución publicada del caso de Google, de recomendaciones (2016) a control autónomo (2018).
\end{solucion}

\begin{solucion}{eje:6-capas}
\planteamiento \cref{fig:capas}: cada función en la capa de su escala de tiempo.
\pasos
\begin{enumerate}[label=\arabic*.]
  \item Temperatura de impulsión: control básico (PID), segundos.
  \item Consignas horarias: control supervisor (MPC o RL), minutos.
  \item Parada por alta presión: protección o sistema de seguridad, milisegundos a segundos, con lógica y sensores propios.
  \item Previsión del precio: planificación, horas o días.
  \item Detección de sonda averiada: supervisión y diagnóstico (comparación con redundancias o con el gemelo), minutos, con alarma al operador.
  \item La protección no puede depender de la capa supervisora porque esta es software complejo, puede fallar o comunicarse mal y, por diseño, no se certifica como función de seguridad; la protección debe actuar aunque todo lo demás falle.
\end{enumerate}
\resultado \textbf{básico, supervisor, seguridad, planificación y diagnóstico, respectivamente; la seguridad es independiente.}
\comprobacion la asignación respeta que cada capa trabaje más lenta que la inferior.
\end{solucion}

\begin{solucion}{eje:6-alarmas}
\datos 9 alarmas/h; referencias de ISA-18.2: unas 150/día (1 cada 10 min) aceptable; unas 300/día (2 cada 10 min) máximo manejable.
\planteamiento convertir tasas.
\pasos
\begin{enumerate}[label=\arabic*.]
  \item (a) $9\cdot24=216$ al día; $9/6=1{,}5$ cada 10 minutos: entre lo aceptable y el máximo manejable.
  \item (b) $12\cdot24=288$ al día; $2$ cada 10 minutos: en el máximo manejable, sin margen para episodios anómalos.
  \item (c) Racionalizar: solo alarmas que exijan acción (la salida al modo de respaldo y la pérdida de comunicación sí; un residuo algo alto no); agrupar las relacionadas; prioridades; suprimir las del MPC cuando está desactivado; enviar lo informativo a registros y paneles.
\end{enumerate}
\resultado \textbf{216 y 288 al día; la capa nueva lleva la carga al máximo manejable si no se racionaliza.}
\comprobacion $288/300=0{,}96$ del máximo de referencia.
\end{solucion}

\begin{solucion}{eje:6-contraste}
\datos $n=20$; $\bar d=1{,}20$ €; $s=1{,}50$ €; $t_{0{,}975;19}=2{,}093$.
\planteamiento contraste $t$ pareado (módulo~0.9, §5.3).
\pasos
\begin{enumerate}[label=\arabic*.]
  \item (a) $s/\sqrt{20}=0{,}3354$; $t=3{,}578$; $p=0{,}0020$; intervalo $[0{,}498;\ 1{,}902]$ €.
  \item (b) No significativo si $t<2{,}093$: $s>1{,}2\sqrt{20}/2{,}093=2{,}564$ €.
  \item (c) Semiamplitud $t_{0{,}975;n-1}\cdot1{,}5/\sqrt n\le0{,}3$: con la aproximación normal, $n\approx(1{,}96\cdot1{,}5/0{,}3)^{2}=96$; con la $t$, $n=99$ días.
  \item (d) El emparejamiento elimina la variabilidad entre días (tiempo meteorológico, ocupación) y evita confundir el efecto del controlador con las diferencias entre inviernos.
\end{enumerate}
\resultado \textbf{$t=3{,}58$, $p=0{,}002$, $[0{,}50;\ 1{,}90]$ €; no significativo si $s>2{,}56$ €; unos 99 días para $\pm0{,}3$ €.}
\comprobacion el intervalo no contiene el cero, coherente con $p<0{,}05$.
\end{solucion}

\begin{solucion}{eje:6-deriva}
\datos $\sigma=0{,}14$ K; $n=96$.
\planteamiento error típico de la media $\sigma/\sqrt n$.
\pasos
\begin{enumerate}[label=\arabic*.]
  \item (a) $0{,}14/\sqrt{96}=0{,}0143$ K; umbral $3\cdot0{,}0143=0{,}043$ K.
  \item (b) Los residuos de un sistema lento están muy autocorrelacionados: el número efectivo de muestras independientes es mucho menor que 96, el error típico real es mayor y el umbral de (a) daría falsas alarmas frecuentes. Hay que estimar la varianza de la media diaria con datos de días sin deriva (medias por lotes) o con la autocorrelación.
  \item (c) Por ejemplo: alarma de recalibración si la media diaria supera el umbral corregido tres días seguidos, o un CUSUM sobre las medias diarias; nunca recalibrar automáticamente sin validar el modelo nuevo en sombra.
\end{enumerate}
\resultado \textbf{error típico ingenuo $0{,}014$ K y umbral $0{,}043$ K, que hay que ampliar por la autocorrelación.}
\comprobacion con autocorrelación positiva $\Var\bar e=\frac{\sigma^{2}}{n}\bigl(1+2\sum_{j\ge1}(1-j/n)r_j\bigr)>\sigma^{2}/n$.
\end{solucion}

\begin{solucion}{eje:6-filtro}
\datos $a=0{,}9835$; $b=0{,}23$; $d\in[-0{,}0827;\ 0{,}4133]$; $[17,26]$ °C; $P\in[0,3]$.
\planteamiento fórmula de $\Uset_{\mathrm{seg}}(T)$ del \cref{ej:filtro}.
\pasos
\begin{enumerate}[label=\arabic*.]
  \item (a) $T=20$: $aT=19{,}67$; inferior $(17-19{,}67+0{,}0827)/0{,}23<0\to0$; superior $(26-19{,}67-0{,}4133)/0{,}23=25{,}7\to3$: $[0,3]$. $T=25{,}9$: $aT=25{,}473$; superior $(26-25{,}473-0{,}4133)/0{,}23=0{,}496$: $[0;\ 0{,}496]$.
  \item (b) A $25{,}9$ °C se aplica $0{,}496$ kW (recorte de $0{,}5$); a $17$ °C, $\Uset_{\mathrm{seg}}=[1{,}579;\ 3]$ y se aplica $1{,}579$ kW.
  \item (c) Con $\Text\le35$ °C, $d_{\max}=0{,}01653\cdot45=0{,}744$; en $T=26$, el extremo superior es $(26-25{,}571-0{,}744)/0{,}23=-1{,}37<0$: $\Uset_{\mathrm{seg}}$ vacío. Ninguna acción garantiza seguir en el intervalo cuando $T>25{,}68$ °C (sin refrigeración); $[17,26]$ deja de ser invariante.
\end{enumerate}
\resultado \textbf{$[0,3]$ y $[0;\ 0{,}496]$; se aplican $0{,}496$ y $1{,}579$ kW; con $35$ °C de máxima el intervalo no es invariante.}
\comprobacion en $T=25{,}9$ con $P=0{,}496$ y $d=d_{\max}$: $25{,}473+0{,}114+0{,}413=26{,}000$ °C, justo en el límite.
\end{solucion}

\begin{solucion}{eje:6-invariancia}
\planteamiento los extremos de $\Uset_{\mathrm{seg}}(x)$ son funciones afines decrecientes de $x$.
\pasos
\begin{enumerate}[label=\arabic*.]
  \item (a) $\Uset_{\mathrm{seg}}(x)=[\max\{0,\ell(x)\},\min\{\bar u,m(x)\}]$ con $\ell(x)=(x_{\min}-ax-w_{\min})/b$ y $m(x)=(x_{\max}-ax-w_{\max})/b$. Es no vacío si y solo si (i) $\ell(x)\le\bar u$, (ii) $m(x)\ge0$ y (iii) $\ell(x)\le m(x)$.
  \item (iii) equivale a $w_{\max}-w_{\min}\le x_{\max}-x_{\min}$, independiente de $x$. (i) es más exigente cuanto menor es $x$ (porque $\ell$ decrece): basta en $x_{\min}$, $(1-a)x_{\min}-w_{\min}\le b\bar u$. (ii) es más exigente cuanto mayor es $x$: basta en $x_{\max}$, $(1-a)x_{\max}\ge w_{\max}$. Luego no vacío en todo el intervalo $\iff$ no vacío en sus extremos, y por la \cref{prop:invariancia} el intervalo es invariante.
  \item (b) Con $\abs{u_k-u_{k-1}}\le\delta$, las acciones admisibles dependen también de $u_{k-1}$: el estado relevante es $(x_k,u_{k-1})$ y la invariancia debe comprobarse en ese espacio aumentado; puede ocurrir que $x$ sea seguro pero, viniendo de una acción muy distinta, no haya tiempo de corregir.
\end{enumerate}
\resultado \textbf{invariante si y solo si $(1-a)x_{\min}-w_{\min}\le b\bar u$, $(1-a)x_{\max}\ge w_{\max}$ y $w_{\max}-w_{\min}\le x_{\max}-x_{\min}$.}
\comprobacion con los datos del \cref{ej:filtro}: $0{,}0165\cdot17+0{,}0827=0{,}363\le0{,}69$; $0{,}0165\cdot26=0{,}429\ge0{,}413$; $0{,}496\le9$.
\end{solucion}

\begin{solucion}{eje:6-plan}
\planteamiento recorrer las secciones del tema y enlazar con los anteriores.
\pasos
\begin{enumerate}[label=\arabic*.]
  \item Arquitectura: el MPC escribe consignas horarias a los PID de la bomba de calor por OPC UA, a través del filtro; el SIS y las protecciones de la máquina no cambian (sección~6.1).
  \item Modos: dos semanas en sombra, cuatro en asesor, después supervisado; criterios de avance cuantitativos (sección~6.2, \cref{eje:6-modos}).
  \item Alarmas nuevas: salida a respaldo, pérdida de comunicación y fallo del resolvedor, racionalizadas (sección~6.3).
  \item Filtro: intervalo de la \cref{prop:invariancia}, límite de velocidad, perro guardián con transferencia sin saltos al PI (\cref{prop:bumpless}).
  \item Validación: experimento pareado de 8 semanas y contraste (\cref{ej:contraste}, \cref{eje:5-protocolo}).
  \item Seguimiento: residuos diarios del gemelo con umbral corregido y recalibración validada en sombra (\cref{eje:6-deriva}, \cref{tema:2}).
  \item Ciberseguridad: zona propia para la capa supervisora, conducto autenticado y registrado, operación sin nube (sección~6.6).
  \item Regulación: análisis de si el sistema es componente de seguridad de una infraestructura crítica (anexo~III del Reglamento de IA); en una oficina, normalmente no; documentar la supervisión humana igualmente (sección~6.7).
\end{enumerate}
\resultado \textbf{un plan por etapas con salvaguardas, validación estadística, seguimiento y análisis de ciberseguridad y regulación.}
\comprobacion cada elemento corresponde a un error frecuente de la \cref{obs:errores6}.
\end{solucion}


\backmatter
% =====================================================================
\begin{thebibliography}{99}
\addcontentsline{toc}{chapter}{Bibliografía}
\markboth{Bibliografía}{}
\small

\bibitem{cpo}
J.~Achiam, D.~Held, A.~Tamar y P.~Abbeel.
\newblock Constrained policy optimization.
\newblock En \emph{Proceedings of the 34th International Conference on Machine Learning (ICML)}, PMLR~70, 2017.
\newblock \url{https://proceedings.mlr.press/v70/achiam17a.html}

\bibitem{altman}
E.~Altman.
\newblock \emph{Constrained Markov Decision Processes}.
\newblock Chapman \& Hall/CRC, 1999.

\bibitem{casadi}
J.~A.~E. Andersson, J.~Gillis, G.~Horn, J.~B. Rawlings y M.~Diehl.
\newblock CasADi: a software framework for nonlinear optimization and optimal control.
\newblock \emph{Mathematical Programming Computation}, 11(1):1--36, 2019.
\newblock \url{https://web.casadi.org/}

\bibitem{am}
K.~J. Åström y R.~M. Murray.
\newblock \emph{Feedback Systems: An Introduction for Scientists and Engineers}, 2.ª ed.
\newblock Princeton University Press, 2021. Caps.~3 «System Modeling», 5 «Dynamic Behavior», 6 «Linear Systems», 11 «PID Control» y~15 «Architecture and System Design». Versión electrónica libre.
\newblock \url{https://fbsbook.org/}

\bibitem{boptest}
D.~Blum, J.~Arroyo, S.~Huang, J.~Drgoňa, F.~Jorissen, H.~T. Walnum, Y.~Chen, K.~Benne, D.~Vrabie, M.~Wetter y L.~Helsen.
\newblock Building optimization testing framework (BOPTEST) for simulation-based benchmarking of control strategies in buildings.
\newblock \emph{Journal of Building Performance Simulation}, 14:586--610, 2021.

\bibitem{sinergym}
A.~Campoy-Nieves, A.~Manjavacas, J.~Jiménez-Raboso, M.~Molina-Solana y J.~Gómez-Romero.
\newblock Sinergym -- A virtual testbed for building energy optimization with reinforcement learning.
\newblock Prepublicación, 2024.
\newblock \url{https://arxiv.org/abs/2412.08293}

\bibitem{crawley}
D.~B. Crawley, L.~K. Lawrie, F.~C. Winkelmann, W.~F. Buhl, Y.~J. Huang, C.~O. Pedersen, R.~K. Strand, R.~J. Liesen, D.~E. Fisher y M.~J. Witte.
\newblock EnergyPlus: creating a new-generation building energy simulation program.
\newblock \emph{Energy and Buildings}, 33(4):319--331, 2001.
\newblock \url{https://energyplus.net/}

\bibitem{cvxpy}
S.~Diamond y S.~Boyd.
\newblock CVXPY: A Python-embedded modeling language for convex optimization.
\newblock \emph{Journal of Machine Learning Research}, 17(83):1--5, 2016.
\newblock \url{https://www.cvxpy.org/}

\bibitem{drgona}
J.~Drgoňa, J.~Arroyo, I.~Cupeiro Figueroa, D.~Blum, K.~Arendt, D.~Kim, E.~Perarnau Ollé, J.~Oravec, M.~Wetter, D.~L. Vrabie y L.~Helsen.
\newblock All you need to know about model predictive control for buildings.
\newblock \emph{Annual Reviews in Control}, 2020.
\newblock \url{https://doi.org/10.1016/j.arcontrol.2020.09.001}

\bibitem{dulac}
G.~Dulac-Arnold, N.~Levine, D.~J. Mankowitz, J.~Li, C.~Paduraru, S.~Gowal y T.~Hester.
\newblock Challenges of real-world reinforcement learning: definitions, benchmarks and analysis.
\newblock \emph{Machine Learning}, 110, 2021.
\newblock \url{https://arxiv.org/abs/2003.11881}

\bibitem{evansgao}
R.~Evans y J.~Gao.
\newblock DeepMind AI reduces Google data centre cooling bill by 40\%.
\newblock Blog de Google DeepMind, 20 de julio de 2016.
\newblock \url{https://deepmind.google/blog/deepmind-ai-reduces-google-data-centre-cooling-bill-by-40/}

\bibitem{dompc}
F.~Fiedler, B.~Karg, L.~Lüken, D.~Brandner, M.~Heinlein, F.~Brabender y S.~Lucia.
\newblock do-mpc: Towards FAIR nonlinear and robust model predictive control.
\newblock \emph{Control Engineering Practice}, 140:105676, 2023.
\newblock \url{https://www.do-mpc.com/}

\bibitem{td3}
S.~Fujimoto, H.~van Hoof y D.~Meger.
\newblock Addressing function approximation error in actor-critic methods.
\newblock En \emph{Proceedings of the 35th International Conference on Machine Learning (ICML)}, PMLR~80:1587--1596, 2018.
\newblock \url{https://arxiv.org/abs/1802.09477}

\bibitem{pycontrol}
S.~Fuller, B.~Greiner, J.~Moore, R.~M. Murray, R.~van Paassen y R.~Yorke.
\newblock The Python Control Systems Library (python-control).
\newblock En \emph{IEEE Conference on Decision and Control (CDC)}, 2021.
\newblock \url{https://python-control.readthedocs.io/}

\bibitem{gamblegao}
C.~Gamble y J.~Gao.
\newblock Safety-first AI for autonomous data centre cooling and industrial control.
\newblock Blog de Google DeepMind, 17 de agosto de 2018.
\newblock \url{https://deepmind.google/blog/safety-first-ai-for-autonomous-data-centre-cooling-and-industrial-control/}

\bibitem{sac}
T.~Haarnoja, A.~Zhou, P.~Abbeel y S.~Levine.
\newblock Soft actor-critic: Off-policy maximum entropy deep reinforcement learning with a stochastic actor.
\newblock En \emph{International Conference on Machine Learning (ICML)}, 2018.
\newblock \url{https://arxiv.org/abs/1801.01290}

\bibitem{sac2}
T.~Haarnoja, A.~Zhou, K.~Hartikainen, G.~Tucker, S.~Ha, J.~Tan, V.~Kumar, H.~Zhu, A.~Gupta, P.~Abbeel y S.~Levine.
\newblock Soft actor-critic algorithms and applications.
\newblock Prepublicación, 2018.
\newblock \url{https://arxiv.org/abs/1812.05905}

\bibitem{cs329s}
C.~Huyen.
\newblock CS329S: Machine Learning Systems Design.
\newblock Stanford University, invierno de 2022. Clases~8 «Deployment», 10 «Diagnosis of ML system failures \& data distribution shifts \& monitoring» y~11 «Monitoring \& Continual Learning».
\newblock \url{https://stanford-cs329s.github.io/}

\bibitem{iec61511}
International Electrotechnical Commission.
\newblock IEC 61511-1:2016. \emph{Functional safety -- Safety instrumented systems for the process industry sector -- Part 1: Framework, definitions, system, hardware and application programming requirements}.
\newblock IEC, 2016.
\newblock \url{https://webstore.iec.ch/publication/24241}

\bibitem{iec62443}
International Electrotechnical Commission e ISA.
\newblock Serie IEC 62443, \emph{Security for industrial automation and control systems}.
\newblock Resumen público del comité de sistemas de energía inteligente de IEC.
\newblock \url{https://syc-se.iec.ch/deliveries/cybersecurity-guidelines/security-standards-and-best-practices/iec-62443/}

\bibitem{isa182}
International Society of Automation.
\newblock ANSI/ISA-18.2, \emph{Management of Alarm Systems for the Process Industries}, 2009, base de la norma internacional IEC~62682. Resumen público de su ciclo de vida y de sus referencias de carga en: «Implementing alarm management per the ANSI/ISA-18.2 standard», \emph{Control Engineering}.
\newblock \url{https://www.controleng.com/articles/implementing-alarm-management-per-the-ansi-isa-18-2-standard/}

\bibitem{kritzinger}
W.~Kritzinger, M.~Karner, G.~Traar, J.~Henjes y W.~Sihn.
\newblock Digital Twin in manufacturing: A categorical literature review and classification.
\newblock \emph{IFAC-PapersOnLine}, 51(11):1016--1022, 2018.

\bibitem{lazic}
N.~Lazic, T.~Lu, C.~Boutilier, M.~K. Ryu, E.~J. Wong, B.~Roy y G.~Imwalle.
\newblock Data center cooling using model-predictive control.
\newblock En \emph{Advances in Neural Information Processing Systems 31 (NeurIPS)}, 3818--3827, 2018.
\newblock \url{https://research.google/pubs/data-center-cooling-using-model-predictive-control/}

\bibitem{cs285}
S.~Levine.
\newblock CS285: Deep Reinforcement Learning.
\newblock University of California, Berkeley, otoño de 2023. Clases~7 «Value Function Methods», 8 «Deep RL with Q-Functions», 10 «Optimal Control and Planning» y~15 y~16 «Offline Reinforcement Learning».
\newblock \url{https://rail.eecs.berkeley.edu/deeprlcourse-fa23/}

\bibitem{cs285hw3}
S.~Levine y equipo docente de CS285.
\newblock CS185/285: Deep Reinforcement Learning, Decision Making, and Control. Tarea~3, «Q-Learning and Actor-Critic Algorithms».
\newblock University of California, Berkeley, primavera de 2026. §2.5 (doble Q-learning) y §3 (\emph{soft actor-critic}).
\newblock \url{https://rail.eecs.berkeley.edu/deeprlcourse/static/homeworks/hw3.pdf}

\bibitem{ddpg}
T.~P. Lillicrap, J.~J. Hunt, A.~Pritzel, N.~Heess, T.~Erez, Y.~Tassa, D.~Silver y D.~Wierstra.
\newblock Continuous control with deep reinforcement learning.
\newblock En \emph{International Conference on Learning Representations (ICLR)}, 2016.
\newblock \url{https://arxiv.org/abs/1509.02971}

\bibitem{ljung}
L.~Ljung.
\newblock \emph{System Identification: Theory for the User}, 2.ª ed.
\newblock Prentice Hall, 1999. Texto del curso 6.435 del MIT.

\bibitem{ars}
H.~Mania, A.~Guy y B.~Recht.
\newblock Simple random search of static linear policies is competitive for reinforcement learning.
\newblock En \emph{Advances in Neural Information Processing Systems 31 (NeurIPS)}, 2018.
\newblock \url{https://arxiv.org/abs/1803.07055}

\bibitem{mit6435}
M.~A. Dahleh.
\newblock 6.435: System Identification.
\newblock Massachusetts Institute of Technology, MIT OpenCourseWare, primavera de 2005. Temario y hoja de problemas~5.
\newblock \url{https://ocw.mit.edu/courses/6-435-system-identification-spring-2005/}

\bibitem{mit10450}
B.~S. Johnston.
\newblock 10.450: Process Dynamics, Operations, and Control.
\newblock Massachusetts Institute of Technology, MIT OpenCourseWare, primavera de 2006. Lecciones~5 «The Heated Tank» y~6 «The Exothermic Reactor».
\newblock \url{https://ocw.mit.edu/courses/10-450-process-dynamics-operations-and-control-spring-2006/}

\bibitem{fmi}
Modelica Association.
\newblock \emph{Functional Mock-up Interface (FMI)}, versión~3.0.
\newblock \url{https://fmi-standard.org/}

\bibitem{modelica}
Modelica Association.
\newblock \emph{Modelica Language} y bibliotecas asociadas.
\newblock \url{https://modelica.org/}

\bibitem{montgomery}
D.~C. Montgomery.
\newblock \emph{Design and Analysis of Experiments}, 9.ª ed.
\newblock Wiley, 2017. Cap.~2 «Simple Comparative Experiments», §2.5 sobre comparaciones pareadas.

\bibitem{narendra}
K.~S. Narendra y K.~Parthasarathy.
\newblock Identification and control of dynamical systems using neural networks.
\newblock \emph{IEEE Transactions on Neural Networks}, 1(1):4--27, 1990.

\bibitem{ng}
A.~Y. Ng, D.~Harada y S.~Russell.
\newblock Policy invariance under reward transformations: Theory and application to reward shaping.
\newblock En \emph{International Conference on Machine Learning (ICML)}, 278--287, 1999.

\bibitem{nw}
J.~Nocedal y S.~J. Wright.
\newblock \emph{Numerical Optimization}, 2.ª ed.
\newblock Springer, 2006. Cap.~17 «Penalty and Augmented Lagrangian Methods», §17.2 y teorema~17.3.

\bibitem{peng}
X.~B. Peng, M.~Andrychowicz, W.~Zaremba y P.~Abbeel.
\newblock Sim-to-real transfer of robotic control with dynamics randomization.
\newblock En \emph{IEEE International Conference on Robotics and Automation (ICRA)}, 2018.
\newblock \url{https://arxiv.org/abs/1710.06537}

\bibitem{sb3}
A.~Raffin, A.~Hill, A.~Gleave, A.~Kanervisto, M.~Ernestus y N.~Dormann.
\newblock Stable-Baselines3: Reliable reinforcement learning implementations.
\newblock \emph{Journal of Machine Learning Research}, 22(268):1--8, 2021.

\bibitem{rmd}
J.~B. Rawlings, D.~Q. Mayne y M.~M. Diehl.
\newblock \emph{Model Predictive Control: Theory, Computation, and Design}, 2.ª ed.
\newblock Nob Hill Publishing, 2017. Caps.~1 «Getting Started with Model Predictive Control», 2 «Model Predictive Control: Regulation» y~8 «Numerical Optimal Control». Versión electrónica libre.
\newblock \url{https://sites.engineering.ucsb.edu/~jbraw/mpc/}

\bibitem{aiact}
Parlamento Europeo y Consejo de la Unión Europea.
\newblock Reglamento (UE) 2024/1689, de 13 de junio de 2024, por el que se establecen normas armonizadas en materia de inteligencia artificial. Capítulo~III, sección~2 (artículos~8 a~15) y anexo~III.
\newblock \url{https://eur-lex.europa.eu/eli/reg/2024/1689/oj}

\bibitem{sculley}
D.~Sculley, G.~Holt, D.~Golovin, E.~Davydov, T.~Phillips, D.~Ebner, V.~Chaudhary, M.~Young, J.-F. Crespo y D.~Dennison.
\newblock Hidden technical debt in machine learning systems.
\newblock En \emph{Advances in Neural Information Processing Systems 28 (NeurIPS)}, 2503--2511, 2015.

\bibitem{seborg}
D.~E. Seborg, T.~F. Edgar, D.~A. Mellichamp y F.~J. Doyle~III.
\newblock \emph{Process Dynamics and Control}, 4.ª ed.
\newblock Wiley, 2016.

\bibitem{skogestad}
S.~Skogestad.
\newblock Simple analytic rules for model reduction and PID controller tuning.
\newblock \emph{Journal of Process Control}, 13:291--309, 2003.

\bibitem{ee363}
Stanford University.
\newblock EE363: Linear Dynamical Systems.
\newblock Primavera de 2026. Clases~16 «Linear quadratic regulator: Discrete-time finite horizon», 17 «LQR via Lagrange multipliers» y~18 «Infinite horizon linear quadratic regulator».
\newblock \url{https://ee363.stanford.edu/lectures.html}

\bibitem{nist}
K.~Stouffer, M.~Pease, C.~Tang, T.~Zimmerman, V.~Pillitteri, S.~Lightman, A.~Hahn, S.~Saravia, A.~Sherule y M.~Thompson.
\newblock \emph{Guide to Operational Technology (OT) Security}. NIST Special Publication 800-82, revisión~3.
\newblock National Institute of Standards and Technology, septiembre de 2023.
\newblock \url{https://csrc.nist.gov/pubs/sp/800/82/r3/final}

\bibitem{osqp}
B.~Stellato, G.~Banjac, P.~Goulart, A.~Bemporad y S.~Boyd.
\newblock OSQP: an operator splitting solver for quadratic programs.
\newblock \emph{Mathematical Programming Computation}, 12(4):637--672, 2020.

\bibitem{sb}
R.~S. Sutton y A.~G. Barto.
\newblock \emph{Reinforcement Learning: An Introduction}, 2.ª ed.
\newblock MIT Press, 2018. Cap.~3, ejercicios~3.15 y~3.16; §6.7 «Maximization Bias and Double Learning»; cap.~13 «Policy Gradient Methods».
\newblock \url{http://incompleteideas.net/book/the-book-2nd.html}

\bibitem{tobin}
J.~Tobin, R.~Fong, A.~Ray, J.~Schneider, W.~Zaremba y P.~Abbeel.
\newblock Domain randomization for transferring deep neural networks from simulation to the real world.
\newblock En \emph{IEEE/RSJ International Conference on Intelligent Robots and Systems (IROS)}, 2017.
\newblock \url{https://arxiv.org/abs/1703.06907}

\bibitem{gymnasium}
M.~Towers, A.~Kwiatkowski, J.~Terry, J.~U. Balis \emph{et al.}
\newblock Gymnasium: A standard interface for reinforcement learning environments.
\newblock Prepublicación, 2024.
\newblock \url{https://arxiv.org/abs/2407.17032}

\bibitem{acados}
R.~Verschueren, G.~Frison, D.~Kouzoupis, J.~Frey, N.~van Duijkeren, A.~Zanelli, B.~Novoselnik, T.~Albin, R.~Quirynen y M.~Diehl.
\newblock acados -- a modular open-source framework for fast embedded optimal control.
\newblock \emph{Mathematical Programming Computation}, 2021.
\newblock \url{https://docs.acados.org/}

\bibitem{wanghong}
Z.~Wang y T.~Hong.
\newblock Reinforcement learning for building controls: The opportunities and challenges.
\newblock \emph{Applied Energy}, 269:115036, 2020.

\bibitem{zn}
J.~G. Ziegler y N.~B. Nichols.
\newblock Optimum settings for automatic controllers.
\newblock \emph{Transactions of the ASME}, 64:759--768, 1942.

\end{thebibliography}


\end{document}
